
%%%%%%%%%%%%%%%%%%%%%%%%%%%%%%%% SCRIPT FOR E-RARA AND MDZ FILES     %%%%%%%%%%%%%%%%%%%%%%%%%%%%%%%%%%%%%%%%%%%%%%%%
%%%%%%%%%%%%%%%%%%%%%%%%%%% fini le 30.04.2025 par F. GOY            %%%%%%%%%%%%%%%%%%%%%%%%%%%%%%%%%%%%%%%%%%%%%%%%
% !TeX TS-program = lualatex
\documentclass{article}
\usepackage[T1]{fontenc}
\usepackage{microtype}
\usepackage[pdfusetitle,hidelinks]{hyperref}

\usepackage{polyglossia}
\setmainlanguage{english}
\setotherlanguages{latin,greek}
\usepackage[series={},nocritical,noend,noeledsec,nofamiliar,noledgroup]{reledmac}
\usepackage{reledpar}

\usepackage{fontspec}
\setmainfont{TeX Gyre Termes}

\usepackage{sectsty}
\usepackage{xcolor}

\usepackage{fancyhdr}
\pagestyle{fancy}
\fancyhf{}
\fancyhead[LE,RO]{\nouppercase{\leftmark}}  
\cfoot{\thepage}
\renewcommand{\headrulewidth}{0.4pt}

% Redefine \section to remove numbering
\usepackage{titlesec}
\titleformat{\section}[block]{\normalfont\scshape\color{gray}}{}{0pt}{} % no number in heading
\titleformat{\subsection}[hang]{\normalfont}{}{0pt}{} % also remove subsection number
\titleformat{\subsubsection}[hang]{\normalfont\footnotesize\color{black}}{}{0pt}{}

% Modify how section marks are stored to exclude numbers
\makeatletter
\renewcommand{\sectionmark}[1]{%
	\markboth{#1}{}} % Only store the section title, without number
\renewcommand{\subsectionmark}[1]{%
	\markright{#1}} % Only store the subsection title, without number
\renewcommand{\numberline}[1]{} % Hide the section number in TOC
\makeatother

\begin{document}

\date{}
        \title{Commentarii In Epistolas D. Pavli Ad Timotheum : [Bullinger, Heinrich], [1536]}
\maketitle
\tableofcontents
\clearpage
\begin{pages} 
\beginnumbering
        
\marginpar{[ p.92 ]}
\phantomsection
\addcontentsline{toc}{subsection}{\textit{ARGVMENTVM PRIMAE EPI- }}
\subsection*{\textit{ARGVMENTVM PRIMAE EPI- }}
\phantomsection
\addcontentsline{toc}{subsection}{\textit{stolae ad Timotheum. }}
\subsection*{\textit{stolae ad Timotheum. }}\pstart \huge\textbf{E}\normalsize X OMNIBVS epistolis Paulinis liquet Ti- motheum hunc absolutissimae cum pietatis tum e- ruditionis fuisse exemplar: unde iam colligi- mus epistolam hanc non tam in ipsius gratiam esse seriptam, qui uel non admonitus sapiebat et omnia faciebat recte, quam in omnium aliorum episcoporum utilitatem qui per successionem temporum fuerant ee clesiis praeficiendi. Proponit enim in hac absolutissi- inum exemplar boni pastoris et diligentis episcopi, orationis filum secundum materiarum praestantiam et episcopi officia atque attributa ducens.  \pend\pstart Principio enim cum potissima episcopi functio sit Docere, initio admonet utreiectis omnibus fabulis quae stionibusque curiosis et humanis traditionibus ea tra- dat quae sint charitatis et fidei uerae. In his enim duo- bus capitibus contineri doctrinam sanam.  \pend\pstart Deinde cum ad episcopi quoque curam pertineat praeesse coetui ecclesiastico, precibus sacris et eccle- siasticis ritibus, secundo loco disserit de eo quomodo deprecationes et interpellationes fieri debeant pro omnibus hominibus, item quid hic deceat uiros quid mulieres.  \pend\pstart Quia uero plurimum momenti habet exemplum  \pend
\vspace{0.5cm}\noindent
\vspace{0.2cm}\rule{1cm}{0.2pt}\\ 
\hspace{0.2cm}\textit{mg}
\footnotesize Occanio. 
\normalsize| 
\hspace{0.2cm}\textit{mg}
\footnotesize Scopus. 
\normalsize| 
\section*{AD TIM. }
\marginpar{[ p.93 ]}\pstart episcopi sanctum, et prophanum siue impurum offi- cit plurimum, ideo episcopum cum tota familia depin git qualis scilicet esse debeat et quid eum deceat qui domui dei praefectus est.  \pend\pstart Veram doctrinam plerumque sequitur falsa ut um- bra corpus. Nescit enim inimicus homo conquiescere, semperque malum semen agro dominico bono semine se minato inserere gliscit: specimen ergo huius exhibet apostolus, partim ut ab erroribus caueremus, partim ut tenacius adhaereremus doctrinae sanae. Atque huc ma xime uocat Timotheum, praescribens quomodo se ge- rat quilibet episcopus in rerum discrimine. Omnia ue£ ro maximo spiritu prosequitur.  \pend\pstart Haec autem de doctrina sana, de precibus sacris, de uita et moribus episcoporum, de doctrina falsa, et quod assidue laborandum sit fideli agricolae in uinea domini, potissima sunt huius epistolae capita, quae se- quuntur de uariis rebus praecipiunt. Praescribit enim quomodo se gerat erga seniores, anus, iuuenes, et pu ellas, item erga uiduas grandaeuas et adolescentiores id est aetatis suspectae mulierculas, rursus quomodo er ga episcopos siue presbyteros, quid praescribat seruis ac dominis. Ibidem et auaritiam accusat et ἀυτόρο nβαp laudat, docens quantis cum periculis coniun- ctum sit studium pecuniae, item quis sit usus opum, et quae sint opulentiorum officia.  \pend
\section*{AD TIM. CAP. I. }
\marginpar{[ p.94 ]}\pstart Vbique uero immiscet quod ipsis prae caeteris uoluit esse commendatissimum, ut cauerent sibi a quaestioni- bus et disputationibus curiosis spinosisque ́.  \pend\pstart Stilus congruit argumento. Institutiones enim fa- eili et dilucido stilo exarandae sunt. Et facilia et ex- posita in hac sunt omnia. Aiunt scriptam a Laodieea Carica. De qua re annotaui quaedam in 4, cap. ad Co- loss. Verum parum retulerit unde hanc scripserit siue- e Macedonia siue ex Asia. In hac certe non diu haesisse legitur a tumultu Ephesino. Protinus enim soluit ie Macedoniam inde transgressus in Graeciam, a qua ite rum in Macedoniam ueniens breuissimo temporis spa cio nauigauit in Syriam, ex hac uero captiuus Ro- mam ducitur.  \pend
\endnumbering\beginnumbering\section{CAP. I.}
\phantomsection
\addcontentsline{toc}{subsection}{\textit{\huge\textbf{P}\normalsize AVLVS apostolus le- su Christi, iuxta delegati- onem Dei seruatoris no- stri, et domini nostri lesu Christi, qui est spes nostra, Timotheo ger- mano filio in fide, Gratia, miseri- cordia, pax a Deo patre nostro et domino lesu Christo domino nostro. }}
\subsection*{\textit{\huge\textbf{P}\normalsize AVLVS apostolus le- su Christi, iuxta delegati- onem Dei seruatoris no- stri, et domini nostri lesu Christi, qui est spes nostra, Timotheo ger- mano filio in fide, Gratia, miseri- cordia, pax a Deo patre nostro et domino lesu Christo domino nostro. }}\pstart Inscriptione hac duo insinuat. Primum apostolum  \pend
\vspace{0.5cm}\noindent
\vspace{0.2cm}\rule{1cm}{0.2pt}\\ 
\hspace{0.2cm}\textit{mg}
\footnotesize Stilus. 
\normalsize| 
\section*{COMMENT. IN I. EPIST. }\pstart se dei esse non sua uoluntate, hoc est quod id muneris sibi ob compendium aut gloriam uindicarit, sed dei uocatione et ordinatione. Nam κατ ἐπιταγίμ idem pollet ac si quis dicat iuxta iniunctionem. ErrIrα TIouenim est iniungere, aliquid officii delegare et mandare quidpiam. Hoc ipso parat sibi authoritatem, ut quod paulo post propositurus est maius grauiusque ́ habeat pondus. Quasi diceret, Apostolus sum, proinde quae traditurus sum non meo sed illius spiritu pfectasunt qui me misit, ut iam horum meorum praeceptorum con temptus in illum redundet. Deinde summam totius e- uangelii perstringit. Praedicat enim Deum esse saluato rem beneficum et misericordem qui dederit nobis fi- Iium suum in argumentum concordiae quam inierit nobiscum et uitae sempiternae, quam nobis omnibus conferre uelit. Vnde ipse dominus Iesus spes et expe- ctatio salus et summum bonum est omnium fidelium. Maledicti autem sunt qui confidunt in hominem et diuos creaturasue alias spes suas uocant. Ierem.17. Ger- manum siue yvήσiop filium uocat uerum et ingenu um qui moribus et genio patrem omnino referebat. Nam ad Philippenses, Neminem, ait Paulus, habeo pa ri mecum animo praeditum qui germane res uestras cu raturus sit. Et ne quis carnalem filium esse Pauli Timo theum crederet, adiecit, In fide. Nam et de Corinthiis scribit apostolus, In Christo Iesu per euangelium ego  \pend
\section*{AD TIM. CAP. I. }
\marginpar{[ p.95 ]}\pstart uos genui. Et fides ac synceritas Pauli in Timotheo ut purissimo speculo relucebat. Theophylactus, Miseri- cordiam noue adiectam putat ex nimia amoris uehe- mentia. Domini uero nomen ideo geminauit, ut intelli geremus diuinitatem et potentiam Christi. Vnde Pa- raphrastes, Hic ut habet cum patre communia omnia ita non frustratur ope sua qui sese totos semel ipsius fi- dei concrediderunt tanquam serui fideles ab optimo simul et potentissimo hero modis omnibus pendentes.  \pend
\phantomsection
\addcontentsline{toc}{subsection}{\textit{DOCTRINA SANA TRA- denda ecclesiis. }}
\subsection*{\textit{DOCTRINA SANA TRA- denda ecclesiis. }}
\phantomsection
\addcontentsline{toc}{subsection}{\textit{Quemadmodum rogaui te ut remaneres Ephesi cum proficisce- rer in Macedoniam ita facito ut de- nuncies quibusdam ne diuersam sequantur doctrinam, nec atten- dant fabulis et genealogiis nun- quam finiendis, quae quaestiones praebent magis quam aedificatio- nem dei quae est per fidem. }}
\subsection*{\textit{Quemadmodum rogaui te ut remaneres Ephesi cum proficisce- rer in Macedoniam ita facito ut de- nuncies quibusdam ne diuersam sequantur doctrinam, nec atten- dant fabulis et genealogiis nun- quam finiendis, quae quaestiones praebent magis quam aedificatio- nem dei quae est per fidem. }}\pstart Prima et potissima doctoris ecclesiastici siue epi- scopi functio est docere. Neque quicquam ecclesia ha-  \pend
\vspace{0.5cm}\noindent
\vspace{0.2cm}\rule{1cm}{0.2pt}\\ 
\hspace{0.2cm}\textit{mg}
\footnotesize In quem fi- nem institu tum sit mini steriñ uerbi 
\normalsize| 
\section*{COMMENT. IN I. EPIST. }\pstart bet maius aut utilius. Et sathan huic imprimis insidia- tur ut adulteret uitiet subuertat. Prima ergo episcopi cura sit doctrinam sanam conseruare illibatam. Pau- lus certe de hac primo praecipit et breui expositione proponit in quem finem et usum Ephesi reliquerit Timotheum, ut scilicet doctrinam euangelicam asser- tione constanti tueretur aduersus omnium incursio- nes nebulas et subuersiones. Constat enim omnes uer- bi ministros in hoc esse a Domino in ecclesia institutos ut doctrinam tradant ueritatis, errores et superftitio nes destruant. Ita enim Paulus, In hoc reliqui te Ephesi ait, ut denuncies et cum authoritate ac grauitate tra das quibusdam curiosis et malae fidei hominibus μù ετεροδιιδεασκαλο͂ν, id est ne diuersa doceant, aut, ne diuersis utantur doctoribus. Est enim ambiguum uer- bum quod refertur ad utrumque . Est autem una uera et sancta ueritatis doctrina quae prodita est scripturis canonicis. Huius doctores et authores sunt prophe- tae et apostoli spiritusancto inspirati. Ab his non disce det pastor bonus uel latum unguem, nec patietur quen quam ecclesiae suae diuersos ab his doctores inferre cu- pientem. Superioris saeculi pastores hac in parte pec- carunt dum diuersos doctores intromittentes diuer- sam quoque ab apostolica doctrinam receperunt, qua re factum est ut tota ecclesia plusquam Iudaicis fabulis suerit oppleta. Qui enim ueritatis doctrinam deserunt  \pend
\section*{AD TIM. CAP. I. }
\marginpar{[ p.96 ]}\pstart aut non magnifaciunt uel hanc non unice obseruant incidunt in errores uarios. Inde sequitur in Paulo, Nec attendant fabulis et genealogiis nunquam finiendis. Fabulas Iudaicas non uocauit sacras literas aut scriptu ras propheticas sub figurarum inuolucro maxima con tinentes mysteria, sed constitutiunculas humanas de qui bus in euangelio legere est apud Matth.in15. et 23. cap.Item Marci 2.Huc pertinere crediderim quae ma- gno cum supercilio tradebant de iusti ficationibus car£ nis, de lotione hoc est uariis baptismatum generibus, de ciborum delectu, de circumcisione, de obseruatio- ne dierum, idque genus aliis nugis. Nam ad Titum, Re- darguito illos seueriter ait ut sani sint in fide, non at- tendentes ludaicis fabulis et praeceptis hominum a- uersantium ueritatem. Et perplexa erat series genea- logiae usque ab auis ac pauis tritauisque repetita. Glori- abantur Iudaei in carne, iactabant se semen esse Abrahae. Quilibet suos recensebat maiores, atque ita occupaban tur rebus nihili. Haec uero pro ratione illorum tempo rum praescripta sunt Timotheo, nos obseruare debe- mus similia. Similia apud nos sunt quae prodiderunt monachi et sacrifici de delectu uestium, ciborum, die rum, de lustratione, uotis, peregrinationibus, diuis, et id genus superstitionibus et fabulis innumeris. Haereti ci quoque , ut est apud Irenaeum et Tertullianum, non modo uerborum monstra commenti sunt sed nexus er  \pend
\section*{COMMENT. IN I. EPIST. }\pstart rorum et impietatis insolubiles nexuerunt. Scholastici quoque Theologi prodigiosam et fabulis fabulosiorem Theologiam ex Aristotelica philosopbia consuerunt. Ab his ineptiis quam alienissimos uult esse Christia- nos apostolus. Causa. Quaestiones praebent magis quam aedificationem dei, quae est per fidem. Duplex inquit incommodum oboritur ex illis. Prius inuoluuntur istis auditores dum quaestio ex quaestione ducitur. Posteri- us ad ueram pietatem nihil conferunt. Fide plantatur pietas, fide parantur corda, non disputationibus per- plexis quaestionibus et inexplicabilibus argumento- rum labirynthis. Aedificandi uerbo libenter utitur Paulus, de qua metaphora alibi disputatum est. Breui ter tradita est nobis uera pura et simplex a Deo per prophetas et apostolos doctrina, in hac perseuere- mus constantes et nulla parte declinantes ad huma- nas fabulas et doctores spinosos et quaestionibus in- flatos. Qui enim quaestionum et argumentorum no- dos alios super alios nectunt, quid aliud docent quam dubitare? Vera puritas et simplicitas fidem in Deum tradit quam purissime et simplicissime et ea quae ae- dificationis et charitatis sunt, hoc est quae hominem in pietate prouehunt ornant et perficiunt.  \pend
\phantomsection
\addcontentsline{toc}{subsection}{\textit{Porro finis praecepti est chari- tas ex puro corde et conscientia }}
\subsection*{\textit{Porro finis praecepti est chari- tas ex puro corde et conscientia }}
\section*{AD TIM. CAP. I. }
\marginpar{[ p.97 ]}
\phantomsection
\addcontentsline{toc}{subsection}{\textit{bona et fide non simulata, a quibus quod aberrarunt quidam deflexe- runt ad uaniloquium, uolentes es- se legis doctores non intelligentes quae loquuntur, neque de quibus as- seuerant. }}
\subsection*{\textit{bona et fide non simulata, a quibus quod aberrarunt quidam deflexe- runt ad uaniloquium, uolentes es- se legis doctores non intelligentes quae loquuntur, neque de quibus as- seuerant. }}\pstart Edisserit nunc quis sit scopus et quae summa pu- rae doctrinae, a qua prohibet nos deflectere, charitas ui delicet. Omnis inquit scriptura, et omnia dei instituta, omnes denique leges dei in hoc intendunt ut charitatem nobis commendent. Πlαραγγἐλεμ est praecipere, denuntiare interdicere et monere. Inde ταραγγελ Aiet, quo uerbousus est hic Paulus praeceptum signi- ficat aedictum denuntiationem et admonitionem. Id quod non tam legi Mosaicae quam uniuersae scripturae competit. Et Hebraei legem uocantes nn Thora quasi quis dicat directorium, frequentius non ipsum decalogum sed omnem doctrinam et institutionem intelligunt. Quod si quis contendat non omnem scri- pturam per praeceptum sed legem tantum esse intelli- gendam, iam hic erit uerborum Pauli sensus, Qui fabu- las Iudaicas docent et in umbris suam exercent pale- stram subinde quidem legem Mosis iactitant adducunt adhibentque demonstrationibus suis, caeterum nondum  \pend
\vspace{0.5cm}\noindent
\vspace{0.2cm}\rule{1cm}{0.2pt}\\ 
\hspace{0.2cm}\textit{mg}
\footnotesize Scopus le- gis et toti- us scriptu- rae est chas ritas. 
\normalsize| 
\section*{COMMENT. IN I. EPIST. }\pstart didicerunt in quem finem lex a deo nobis per Mosen data sit: non equidem ut rebus ludicris daremus ope- ram, ut externis iustificationibus purgaremus nosmete ipsos, ut traditiunculis pharisaicis seruiremus deo, ut in genealogiis enarrandis et huiusmodi fabulis aliis occuparemur, sedut Deum et proximum nostrum uere et ex animo diligeremus. Hic autem non legis modo Mosaicae sed omnium prophetarum quoque sco- pus est. Ideo in euangelio secundum Marcum legimus interroganti legisperito quodnam esset praecipuum et primarium omnium praeceptorum? respondisse, Primum omnium praeceptum est, Audi Israel domi- nus Deus noster dominus unus est. Et diliges dominum deum tuum ex toto corde tuo et ex tota anima tua et ex tota mente tua et ex totis uiribus tuis. Hoc est primum praeceptum. Et secundum simile hoc est, Di- liges proximum tuum ut teipsum. Maius his aliud prae ceptum non est. Additur apud Matthaeum, In his duo- bus mandatis uniuersa lex et prophetae pendent. Se- gtur in Marco, Et dixit legisperitus, Bene praeceptor in ueritate dixisti, quod unus sit Deus et non fit alius praeter illum quodque diligere eum ex toto corde et ex tota intelligentia et ex tota anima et ex totis ui- ribus ac diligere proximum ut seipsum plus sit quam uniuersa holocaustomata et uictimae. Hinc ergo stul- tissimos et scripturarum prorsus ignaros esse appa-  \pend
\vspace{0.5cm}\noindent
\vspace{0.2cm}\rule{1cm}{0.2pt}\\ 
\hspace{0.2cm}\textit{mg}
\footnotesize Mar. 12. 
\normalsize| 
\hspace{0.2cm}\textit{mg}
\footnotesize Matth. 22 
\normalsize| 
\hspace{0.2cm}\textit{mg}
\footnotesize 1. Reg. 15 
\normalsize| 
\section*{AD TIM. CAP. I. }
\marginpar{[ p.98 ]}\pstart ret doctores humanarum traditionum ceremoniarumque ma gistros quorum unicum studium estut miseram plebecu- lam onerent et implicitam ritibus superstitiosis in ex terna exerceant disciplina. Atqui unum erat necessa rium et sine istis per se absolutum, quo neglecto frustra in sumitur omnis et labor et sumptus, Charitas uidelicet eaque gemina. Charitas ergo gemina a nobis exigitur, charitatem geminam praecipit ubique dominus, charitatis denique opa reqret a nobis iudicaturus uiuos et mortu, os. Quin ergo desertis stultulis traditionibus ad chari tatem conuertimur, deumque ipsum in Christo diligimus, et proximum nostrum non modo non laedimus sed omni mo do beneficentia iuuamus. Ad haec nitendum omnibus. Haec inculcanda sunt omnibus et perpetuo. In his e- nim id est in fide et charitate consistit uera pietas. Fi dem autem oportet esseueram et charitatem synce- ram. Hinc in euangelio dominus non contentus dixis- se, Diliges dominum deum tuum et proximum tuum, insuper modum et rationem praescribens addit, Exto to corde tuo et c. Ad eundem modum et Paulus in praesentiarum non solum dixit, Finis et scopus prae- cepti est charitas, sed subiungit uerae charitatis condi tiones attributa et ingenium. Ex corde puro, id est in tegro et syncero adeoque toto. Tantum enim ualet haec loquutio quantum illa in euangelio, Diliges dominum ex toto corde et ex totis uiribus. Opponit enim scri-  \pend
\vspace{0.5cm}\noindent
\vspace{0.2cm}\rule{1cm}{0.2pt}\\ 
\hspace{0.2cm}\textit{mg}
\footnotesize Isaiae.28. 
\normalsize| 
\hspace{0.2cm}\textit{mg}
\footnotesize Lucae, u1. 
\normalsize| 
\hspace{0.2cm}\textit{mg}
\footnotesize Matth. 25. 
\normalsize| 
\hspace{0.2cm}\textit{mg}
\footnotesize Charitatls ingeniura, 
\normalsize| 
\section*{COMMENT. IN I. EPIST. }\pstart ptura cor purum siue mundum duplicato cordi aul lacero aut imperfecto. Ita enim saepe in sacris literis le gimus quosdam toto et perfecto corde ambulasse cum domino, hoc est uere et recte et synceriter uitam su am ad uoluntatem dei instituisse, alios autem corde du plici et imperfecto. 4.Regum 20. orat Ezechias ad dominum, Obsecro domine memento quaeso quomodo ambulauerim coram te in ueritate et in corde perfe- cto et quod placitum est coram te fecerim. Ita 4.Re- gum 10.cap. Iehuh ad Ionadab ait, Nunquid est cor tuum rectum cum corde meo sicut cor meum cum cor de tuo? Et apud Germanos usurpamus ista uocabula Reyn vno luter hoc modo, ut cum dicimus, Le ist reyn golo ( Luter bumberey (Æuterlich ourch Gott shoc est solius dei nomine, Nil aliud quam aurum et nequitia est. Ita charitas ex puro corde profici- scens ea est quae uni deo ex integro adhaeret et uni- cum complectitur ardentissime, nec quicquam huic in orbe antehabet. Id quod tot uerbis propter illos incul co qui nihil non tribuunt humanis uiribus, denique pu- ritatem ab homine nescio quam absolutam requirunt, non expendentes interim idiomata sacrae linguae et quod Solomon dixit Prouer.2o. Quis potest dicere mundum est cor meum, purus sum a peccato? Et Iob 25.Quid te eleuat cor tuum? Quid est homo ut immacu latus sit et ut iustus appareat natus de muliere? Ecce  \pend
\section*{AD TIM. CAP. I. }
\marginpar{[ p.99 ]}\pstart inter sanctos eius nemo immutabilis, et coeli non sunt mundi in conspectu eius: quanto magis abominabilis et inutilis homo, qui bibit quasi aquam iniquitatem? Sequitur autem in Paulo quod clarius exponit purita- tem istam, Conscientia bona. Quasi dicat, Ex puro cor- de, id est integra erit charitas si animus amantis bene sibi sit conscius. Id autem non fit meriti et operum nostrorum beneficio sed fidei. Qui enim, ait dominus, uenit ad me non esuriet et qui credit in me non sitiet unquam. Et Paulus ad Romanos, Nulla igitur nunc est condemnatio iis qui insiti sunt Christo Iesu. Insiti au- tem sunt Christo Iesu qui credunt uere. Hinc in Paulo protinus subditur uα) πιςεωσ ἀνυποκριτου et fi- de hypocrisi carente hoc est uera. Absoluit ergo cha- ritatem fides, quae purificat corda et tranquillat con scientias omnemque pellit simulationem. Proinde uera charitas est quae ex fide enata est, cuius anima et uis fides est, quae solam Deum arctissime complectitur, in situta eius magnifacit, et suaue putat quicquid dei no mine uel fugere uel cum labore persequi iubetur. Huc si reuoces uniuersam scripturam canonicam, ui- debis quadrare omnia, nec aliud tot libris tradi a sa- cris scriptoribus. His iam fundamentis iactis iterum inuehitur in falsos doctores, docens quanto cum peri- culo coniunctum sit si neglecto scopo Charitatis, in quam perpetuo fidem ueram includimus, quaestiones  \pend
\vspace{0.5cm}\noindent
\vspace{0.2cm}\rule{1cm}{0.2pt}\\ 
\hspace{0.2cm}\textit{mg}
\footnotesize In Mataeo- logos. 
\normalsize| 
\section*{COMMENT. IN I. EPIST. }\pstart curiosae, humanae constitutiones, et aliae his similes nu- gae tractentur. A quibus inquit quod aberrarunt qui- dam. Id est, ab hoc scopo quoniam aberrarunt quidam (allusit autem ad imperitos sagittandi) ideo deuolu ti sunt in ματαιολογ ίαμ id est in uaniloquium, id est dissertationem quae nihil differt ab anicularum nu- gis, quae tametsi satis longae sint inanes tamen et elum bes sine fructu cadunt, adeo ut tandem ne ipsae quidem quid dixerint norint. Vnde in Paulo sequitur, Volen- tes esse legis doctores non intelligentes quae loquun- tur aut de quibus asseuerant. Idem usuuenit hodie no- stris quoque Theologiae doctoribus. Nam licet strenue et subtiliter totos dies in scholis suis disputent de in- stantibus, de formalitatibus, et relationibus, de reali tatibus et friuolis argutiis suis, nullus ad auditores re dit fructus, nec ipsi sese interim intellexerunt. Disputa runt et apud haereticos quidam plus quam credi pos- sit acute, sed nullo cum fructu. Inuenias item hodie quosdam et inter eos qui uideri uolunt euangelicissimi, qui prolixe et prorsus seraphice pneumaticeque com- mentantur de deo, de religione, de regeneratione, de fide iustificante, de sacramentis et sanctimonia uitae, uerum ea ratione ut auditores intellexerint quidem de magnis rebus esse disputatum sed non probe perce perint quales illae aut quis eius disputationis usus. Nec mirum cum non raro ipsi scriptores uel oratores ut  \pend
\section*{AD TIM. CAP. I. }
\marginpar{[ p.100 ]}\pstart dicant et non quid dicant cogitent. Qui uero sapi- unt omnia ad charitatis regulam collimabunt. Non est quod in sacris eruditionem ostentes, gloriam pro- priam quaeras, nouarum rerum et quas nemo ante te uiderit praediceris author, scripturarum interpres o- culatissimus, legisperitusque prudentissimus. Populum Christo adducere debes. Hic uero non accedet nisi a- met Christum. Hunc non amabit nisi cognorit. Pinges ergo hunc suis coloribus ut toto in ipsum pectore ex- ardeat, hunc unum optet, huic soli adhaereat.  \pend
\phantomsection
\addcontentsline{toc}{subsection}{\textit{Scimus autem quod bona sit lex, si quis ea legitime utatur, sciens il- lud quod iusto lex non sit posita, sed iniustis et inobsequentibus, im- piis et peccatoribus, irreuerenti- bus et prophanis, patricidis et ma- tricidis, homicidis, scortatoribus. masculorum concubitoribus, pla- giariis, mendacibus, periuris, et si- quid aliud est quod sanae doctri- nae aduersetur secundum euange- lium gloriaebeati dei quod concre- }}
\subsection*{\textit{Scimus autem quod bona sit lex, si quis ea legitime utatur, sciens il- lud quod iusto lex non sit posita, sed iniustis et inobsequentibus, im- piis et peccatoribus, irreuerenti- bus et prophanis, patricidis et ma- tricidis, homicidis, scortatoribus. masculorum concubitoribus, pla- giariis, mendacibus, periuris, et si- quid aliud est quod sanae doctri- nae aduersetur secundum euange- lium gloriaebeati dei quod concre- }}
\section*{COMMENT. IN I. EPIST. }
\phantomsection
\addcontentsline{toc}{subsection}{\textit{ditum est mihi. }}
\subsection*{\textit{ditum est mihi. }}\pstart Occupatio est, meminerat enim legis, sugillauerat mataeologos et uanos legis doctores: in loco ergo sub iungit, Nequaquam uero haec dicimus quod Mosaicam legem damnemus. Bona est lex. Nam uoluntas dei ma la esse non potest. At lex dei uoluntas est. Proinde Pau lus legis abusum damnauit non ipsam legem dei bo- nam. Abusus legis quis sit nemo nescire potest ex supe rioribus. Cum enim nostro uitio Iudaeorum more bene dicta legis corrumpimus aut in diuersam sententiam trahimus, et aduersus legislatoris sententiam inter- pretamur, lege abutimur. Porro usus legis uerus est si sciamus iusto legem non esse positam sed iniustis, quae sententia propter breuitatem obscura est. Haec autem est ratio eius, Lex non in hoc a deo lata est ut iustificet, ergo ex lege non est iustitia. Caeterum data est ut nos deduceret ad Christum, Galat.3. Iam qui ad Christum uenerunt liberantur a peccatis, mundantur sanguine Christi, et iustificantur pura dei gratia. Qui uero iustificati sunt a lege liberi sunt, ipsum enim legis propositum et scopum Christum tenent. Deinde libe rati sunt a legis maledictione et onere per Christum cuius spiritu imbuti et regenerati impulsu charitatis ultro plus praestant quam ipsa lex praescribat. Amant igitur hilegem quam bonam imo optimam esse insu-  \pend
\section*{AD TIM. CAP. I. }
\marginpar{[ p.101 ]}\pstart is ipsorum cordibus senserunt. Atque hic usus legis ue- rus est. Sic inquam lex bona est, sic iusto non est posi- ta, ita lex nemini grauis est, sed suauis potius, quouis lapide pretioso gratior, melle et fauis dulcior. De qua re copiosius multis in locis disputauimus, ista ue- ro ad praesentia satis esse putaui. Caeterum qui sub hac Christi et euangelica gratia non sunt, adhuc sub lege sunt. Nihil ergo ad hos pertinent sermones de liberta te Christiana et abrogatione legis, iusto tantum nulla lex posita est. Commodum uero hoc in loco ubi de so na doctrina agit disputat de genuino legis usu. Aduer sarii enim leuitatis uanitatisque insimulabant praedica- tionem euangelicam, quasi hac fenestra aperiretur om nibus criminibus rumperenturque sancta legis uincula. Monet ergo quatenus et quibus lex abrogata sit. Ne- minem scilicet illa eximi nisi qui sua sponte quod lex praescribit facit. Facere autem ea qui iuxta euangeli- cam doctrinam Deum et proximum amet ex animo, id quod doceat euangelium, proinde lege teneri adhuc quotquot charitatis sensu et igne carent, imo ipsum sanctum dei euangelium illa ipsa damnare scelera quae lege sint olim damnata. Ita enim concludit hanc peri- odum, Et si quid aliud est quod sanae doctrinae aduer- setur secundum euangelium gloriae beati Dei quod con- creditum est mihi. Explicat autem quae sit doctrina sa na, nimirum quae est secundum euangelium gloriosum  \pend
\vspace{0.5cm}\noindent
\vspace{0.2cm}\rule{1cm}{0.2pt}\\ 
\hspace{0.2cm}\textit{mg}
\footnotesize Buangelii praedicatio non aperit fenestram sceleribus. 
\normalsize| 
\section*{COMMENT. IN I. EPIST. }\pstart beati dei, hoc est ipsum euangelium gloriosum. Glori osum appellat illud ideo quod gloriam dei et lucem puram praedicet, ut est 2. Corinth.3.4. Adiicit praete- rea Deo epitheton Beatus ut significet dei perfectio- nem qui sibi satis sit et nostris minime indigeat, quin potius salutem et foelicitatem suam nobiscum commis nicare cupiat. Particula, Quod concreditum est mihi appensa est ut ex ipsa subtexantur consequentia. Ca- talogus uero peccatorum qui a Paulo contextus hic legitur, in aliis locis prope expositus est, perstringe- mus tamen singula. Aυόμεοισ iniustis uertit interpres, aduerbum significat ex legibus, exponiturque ; haec uox per consequentem, ἀνυποτάκτοισ id est inobsequem tibus, ad uerbum inordinatis. Intelligit enim pertinacis malitiae nebulones qui nulla lege in ordinem cogi pos- sunt. ἀσεβέσι impiis. Hi sunt anumines qui nullum dei cultum nullam religionem curant sed affectus du- ces sequuntur. αμαρτωλοῖσ peccatoribus. Intelle- xit autem sub his elatos, auaros, inuidos, luxuriosos, et si quid aliud est huius ordinis. ανοσ ιoισ irreuerenti bus hoc est nephandis et atrocibus. Βεperήλοισ pro- phanis siue impuris atque pollutis. Graeci enim etέθη- Az, uocabant sacra ad quae et impuris liber erat ac- cessus Patricidis atque Matricidis capitalis quoque poe- na constituta est Leuit.2o. ανδ ροφονοισ homicidis, ut referas tam ad mares quam foeminas. Ad uerbum enim  \pend
\section*{AD TIM. CAP. I. }
\marginpar{[ p.102 ]}\pstart Jonat uiricidis. Sed et Hebraei hanc uoculam Vir fre- quentissime usurpant pro quouis homine. róevoie scortatoribus. De quibus alians saepeut et de et̓ρσενο- κoίταισ masculorum concubitoribus. ανδραποδι sαῖσ plagiariis siue furibus mancipiariis. Vt enim iu reconsultis sacrilegi et abigei dicuntur, qui res uel sacras uel publicas suffurantur, qui aliena pecora et iumenta abigunt et abducunt, sic plagiarii dicuntur qui aliena mancipia furto subducunt. Hoc enim furti genus plagium dicitur iureconsultis. Late patet hoc crimen, de quo legantur ueterum leges. Huc pertinet mendatium. Sequitur ergo Lεύς αις mendacibus, sub quibus omnes compraehendimus artes hypocriticas et impostorias. Periuri sunt qui fidem datam frangunt, sacramentum datum polluunt, denique et quouis mo- do fallunt. Colligimus autem ex hisce Pauli uerbis o- mnibus ad sanam episcopi doctrinam non id tantum pertinere ut omnes ad charitatem inuitet, libertatem et redemptionem Christi praedicet, sed ut libere quo- que constanter et cum authoritate accuset scelera. Nam ubi haec regnant ibi nullum locum habet uel charitas uel meritum Christi uel libertas Christiana, In scelera igitur primam impressionem faciat episco- pus fidelis qui charitatem, meritum Christi et liber- tatem recte fidelibus tradere cupit.  \pend
\section*{COMMENT. IN I. EPIST. }
\phantomsection
\addcontentsline{toc}{subsection}{\textit{Et gratiam habeo, qui me po- tentem reddidit Christo Iesu do- mino nostro, quia fidelem me iu- dicauit ponendo in ministerium qui prius eram blasphemus et per- sequutor et uiolentus, sed et mise- ricordiam adeptus sum, quod ig- norans fecerim per incredulitatem. Exuberauit autem supra modum gratia domini nostri cum fide et dilectione, quae est per Christum lesum. }}
\subsection*{\textit{Et gratiam habeo, qui me po- tentem reddidit Christo Iesu do- mino nostro, quia fidelem me iu- dicauit ponendo in ministerium qui prius eram blasphemus et per- sequutor et uiolentus, sed et mise- ricordiam adeptus sum, quod ig- norans fecerim per incredulitatem. Exuberauit autem supra modum gratia domini nostri cum fide et dilectione, quae est per Christum lesum. }}\pstart Dixerat sibi conereditum esse euangelium, porro ne uideretur hoc iactasse gloriosius, Deo cui solo om- nis honor debetur omnem gloriam transscribit. Ei in- quam gratias ago, inquit, qui imbecillo et huic mune ri longe impari uires suffecit, quibus gloriam Christi domini assererem, idem iudicauit me fore fidelem per suam gratiam, ideo in ministerium uerbi asciuit, ut mea opera in uerbi dispensatione uteretur. Ego enim nihil tale merueram benefitiis meis, qui fueram blasphemus etc. Iiaque multis nunc agit de salutifica dei gratia neq praeter occasionem. Perpetuo enim institutum et se-  \pend
\vspace{0.5cm}\noindent
\vspace{0.2cm}\rule{1cm}{0.2pt}\\ 
\hspace{0.2cm}\textit{mg}
\footnotesize De dei gr. tia. 
\normalsize| 
\section*{AD TIM. CAP. I. }
\marginpar{[ p.103 ]}\pstart ries disputationis Paulinae reuocanda est ob oculos, nempe quod Timotheum et in hoc omnem episco- pum fidelem adhortatur ut reiectis fabulis sanam do- ctrinam ecclesiis tradant, deinde quod certis quibus- dam notis praecipua sanae doctrinae capita adumbrat, in quibus primas tenet charitas, interim et deusu le- gis nonnihil disputauit, et docuit euangelio quoque ea proscribi et damnari peccata atque scelera quae dan- nentur per legem quoque . Sequitur nunc breuis disser- tatio de Gratia dei. Nam haec non postremas partes obtinet inter praecipua sanae doctrinae capita. Hic uero cum in docendo et demonstrando exemplis nihil sit euidentius, ipse non alieno sed suo exemplo proposito docet quanta sit propensa illa dei in nos uoluntas hoc est gratia et beneficentia. Vt uero omnia essent eui- dentissima proprium suum peccatum amplificat ut uis et opulens gratia ex effectu maximo quam propri- issime cognosceretur. Vocauit inquit me dominus ad ministerium uerbi, qui eram blasphemus, persequutor ηα) ὐ Βρις ήσ id est iniurius oppressor uiolentus et confidenter improbus petulans et contumeliosus. His ce ergo gradibus exaggerat peccatum suum. Primus gradus est conuitiari suggillare uerbum domini a- deoque maledicere et blasphemare. Secundus a male- dicis uerbis descendere ad malefacta, hoc est ad con- silia praua ad insidias et persequutionem, Tertius  \pend
\section*{COMMENT. IN I. EPIST. }\pstart gradus est opprimere palam et manus temerarias con fidenter atrociterque iniicere in bene meritos. His omni bus studuisse Paulum et sese blasphemiis, persequuti- one et innocenti sanguine polluisse testis est liber qui scriptus est de actibus Apostolorum. Ego uero, inquit, tantus conuitiator, sanctorum hostis et parricida ue- niam merui et misericordiam consequutus sum, neque  uulgariter. Non enim noxam mihi modo condonauit dominus aut in album amicorum recepit tantum, sed et summam in functionem et omnium honestissimam as- ciuit, ut scilicet nomen eius portarem coram regibus principibus omnibusque terrae populis. Addit, Quod ig norans fecerim per incredulitatem: qua sententia pia culum suum neque extenuat neque innocentiae suae suam uocationem tribuit, sed aequitatem et misericordiam dei praedicat, quae miserta incredulitatis et inscitiae peccantem sed errore potius quam maleuolentia res uocauit. Vnde protinus gratiae omnia illa significan- tius ascribens et hanc praedicans, Exuberauit autem supra modum inquit gratia domini nostri, condonans mihi illa peccata multo grauissima, exuberauit simul et fides id est fidelitas ueritasque eius et dilectio siue studium et amor eius in nos. Quae Iesu Christo domi- no nostro in carnem misso omnium apertissime oculis subiecit et quam propriissime exhibuit. In hoc enim est omnis plenitudo gratiae fidei et ueritatis amoris  \pend
\section*{AD TIM. CAP. I. }
\marginpar{[ p.104 ]}\pstart quoque in nos maximi. Misericordiam et gratiam pol- licitus est patribus nostris, idem in Christo fidelissime omnia praestitit. Atque huc pertinet quod sequitur.  \pend
\phantomsection
\addcontentsline{toc}{subsection}{\textit{Certus sermo et dignus quem modis omnibus amplectamur, quod  Christus lesus uenit in mundum, ut peccatores saluos faceret, quo- rum primus sum ego. Verum ideo misericordiam sum adeptus ut in me primo ostenderet lesus Chri- stus omnem clementiam, ad expri- mendum exemplar iis qui creditu- ri essent in ipso in uitam aeternam. }}
\subsection*{\textit{Certus sermo et dignus quem modis omnibus amplectamur, quod  Christus lesus uenit in mundum, ut peccatores saluos faceret, quo- rum primus sum ego. Verum ideo misericordiam sum adeptus ut in me primo ostenderet lesus Chri- stus omnem clementiam, ad expri- mendum exemplar iis qui creditu- ri essent in ipso in uitam aeternam. }}\pstart Planioribus et copiosoribus uerbis praedicat gra- tiam euangelicam suique exempli fructum ostendit, his ce omnibus fidem in Christum docens. Praemittit au- tem exordiolum quo paret auditorum animos ad res maximas. πιςόσ, ait, ολόγοσ, μαηι πάσησ ἀποδlo- χηῖς ἄξιοσ. Id est, hoc uerbum siue haec sententia quam propositurus sum fidelis uera et certissima est, di- gnaque multis nominibus quam summo studio et ob- uiis ulnis quod dici solet excipiamus amplectamur pe- ctorique nostro astringamus arctissime. Iam subiungit  \pend
\vspace{0.5cm}\noindent
\vspace{0.2cm}\rule{1cm}{0.2pt}\\ 
\hspace{0.2cm}\textit{mg}
\footnotesize Iuidens e- uangelii praedicatio 
\normalsize| 
\section*{COMMENT. IN I. EPIST. }\pstart certam illam et dignissimam de salute nostra senten- tiam, QVOD CHRISTVS IESVS VENIT IN MVN DVM VT PECCATORES SALVOS FACERET. Qua sententia nihil potest dici breuius sed eadem ni- hil maius nobilius salubrius. Singula uerba singula con- tinent mysteria. Omnes enim peccatores sumus, nul- lus suis potuit iustificari operibus. Eramus enim om- nes filii irae, et proinde aeternae damnationis manci- pia: unde et uates Isaias, Omnes inquit ueluti oues er rauimus, unusquisque in uiam suam declinauit, sed po- suit dominus in eum omnium nostrum iniquitates. Ve niens ergo in mundum carnis et sanguinis nostrae particeps factus obtulit se deo in expiationem omnium peccatorum nostrorum. Hic enim est agnus dei qui tol lit peccatum mundi:ut iam solus Christus omnium pec catorum spes salus redemptio et uita sit. Peccatorum inquam, non iustorum. Non enim egent medici opera sani, sed infirmi. Ista uero copiosissime expositum re- peries Matth. cap.1. et Lucae 2. item Matth, 9.11. 15. Lucae quoque 15. praecipue uero Ioannis tertio capite. Rom.3. Ephe.1.et 2.etc. Diuus Ambrosius sancte et erudite haec Pauli uerba explicans, Quid tam gratum inquit tamque iucundum quam peccato- rum indulgentiam praedicari? Quis enim sine peccato est ut hoc munus non cum omni gratiarum actione su scipiat, deum authorem huius praeferens et collaudans  \pend
\section*{AD TIM. CAP. I. }
\marginpar{[ p.105 ]}\pstart ln Christo? qui ut hominem peccatis ablueret de coelen- stibus ad terrena descendens carnem peccati accepit, terrenis se admiscuit ut eum coelestem efficeret. Mori se passus est, ut illum morte erutum paradyso redde- ret immortalem. Quis haec benefitia non in infinitum extollat quae deus praestitit homini, dominus seruo, con ditor creaturae? Haec Ambrosius. Peccant igitur in gratiam euangelicam qui ideo auertuntur a Christo ad diuos quod sint peccatores. Christus enim qui pec catoribus natus est dicit, Venite ad me omnes qui la- boratis et onerati estis et ego reficiam uos. Nec est quod quisquam obtendat scelerum suorum magnitu- dinem et numerum: expendat potius Pauli uerba con- solationis plenissima. Vbi enim pronuntiasset Chri- stum uenisse in mundum ut peccatores saluos faceret, protinus se exempli gratia adiecit dicens, Quorum primus ego sum. Caeterum uocula Primus in praesenti non est ordinis sed excellentiae perinde ualens ac prae cipuus et maximus. Fuerat enim blasphemus perse- quutor uiolentus parricida. Quid istis poterat exco- gitari atrotius? At haec omnia condonauit illi Christus dominus. Nemo ergo in sceleribus suis desperet. Imo ne quis obiiceret, Tametsi Paulo tanta obtigit gratia, eadem tamen non obtinget omnibus, praeocupat illud et ait, Ideo autem misericordiam adeptus sumut in me primo ostenderet Iesus Christus omnem clementiam  \pend
\section*{COMMENT. IN I. EPIST. }\pstart Ed exprimendum exemplar iis qui credituri essent in ipso in uitam aeternam. Quasi dicat, Hoc insigni ex- emplo ad similem spem ueniae prouocare uoluit om- nes, quamlibet prioris uitae maculis inquinatos, modo impia uita relicta et diffisi suis uiribus omnem fidu- tiam constituant in bonitate Christi, quae nunquam de stituit nos usque ad uitam aeternam. In Paulo igitur oro τύπωσις id est expressa imago et exemplum uiuum omnium hominum oculis exposita est quid sibi debe- ant sperare quicunque relictis moribus pristinis credi- derunt in Christum, gratuitam uidelicet peccatorum remissionem et spem indubitatam uitae sempiternae. Est et in ea uoce quam interpres noster clementiam uertit, Graece uero μακροθυμία dicitur, mysterium. Significat enim ad uerbum longanimitatem. Namscri ptura alians testatur deum ad uindictam lento proce- dere gradu, non quod peccata illi placeant, sed quod uolit mortem peccatoris, quin potius conuersionem e- ius uitamque ́. Ad exemplum Pauli pertinent exempla similia ut sunt Dauidis, Adae, Matthaei, Petri, Zachei, Mulieris peccatricis, Latronis, et si quae alia sunt hu ius generis.  \pend
\phantomsection
\addcontentsline{toc}{subsection}{\textit{Regi autem saeculorum immor- tali, inuisibili, soli sapienti Deo, ho- nor, gloria in saecula saeculorum, Amen. }}
\subsection*{\textit{Regi autem saeculorum immor- tali, inuisibili, soli sapienti Deo, ho- nor, gloria in saecula saeculorum, Amen. }}
\section*{AD TIM. CAP. I. }
\marginpar{[ p.106 ]}\pstart Praedicationem gratiae euangelicae orando cum lau de et gratiarumactione prosequitur, interim poten- tiam aeterni dei praedicans ne quis diffideret pollicitis tantis, quin cogitaret potius qualis sit qui promisit et totius euangelicae doctrinae negotiique salutis author est:innuit praeterea soli deo deberi omnem honorem et gloriam, ex hac munificentia nihil sibi homines uindicare posse. Rex saeculorum dicitur Deus, id est rex aeternus omnium aetatum et creaturarum condi- tor praeses rector et dominus. Idem dicitur immorta- lis siueαφθαρτόσ, id est incorruptibilis, qui ut corrum pi non potest et sempiternus est, ita nullis affectioni- bus et infirmitatibus obnoxius. Inuisibilis appellatur, non quod non sit, sed quod immensus sit, et maior quam qui sensibus compraehendi possit. Hic quidem nobis ob imbeeillitatem carnis nostrae inuisibilis est, ibi uero uidebimus eum sicuti est. Solus sapiens est quemadmodum solus Deus bonus est. Omne enim bo- num a solo deo est, et quicquid boni in creaturis re- peritur, id omne a Deo profectum est. Solomon et Da niel mortalium sapientissimi omnem sapientiam uni Deo tribuunt. Notus est enim locus Danielis 2. cap. Theophylactus hunc ipsum Pauli locum enarrans e- gregie scribit in haec uerba, Soli autem sapienti dixit, non quod pater a filio et spiritu sancto separatus sit, sed cum constet et angelos sapientes dici et homi-  \pend
\vspace{0.5cm}\noindent
\vspace{0.2cm}\rule{1cm}{0.2pt}\\ 
\hspace{0.2cm}\textit{mg}
\footnotesize Solus deus sapiens quo nodos 
\normalsize| 
\section*{COMMENT. IN I. EPIST. }\pstart nes, ostendit hunc solum praecipue sapere ut qui fons sit et causa sapientiae: reliqua uero quae sapientiam uendicant condita sunt et creata et participatione quadam hanc assequuntur.  \pend
\phantomsection
\addcontentsline{toc}{subsection}{\textit{Hoc praeceptum comendo tibi fili Timothee secundum praeceden- tes usque ad te prophetias, ut mili- tes in eis bonam militiam, habens fidem et bonam conscientiam. }}
\subsection*{\textit{Hoc praeceptum comendo tibi fili Timothee secundum praeceden- tes usque ad te prophetias, ut mili- tes in eis bonam militiam, habens fidem et bonam conscientiam. }}\pstart Epilogus est. Redit enim ad institutum et summam rei repetendo infert. Hoc igitur mandatum, ait, ô fi- li Timothee, quod non est nouum aut nuper domi nostrae natum, sed secundum prophetias id est scriptu ras propheticas quas nosti ante te fuisse et tibi ceu per manus traditas, docentes geminam dei et proximi charitatem, iusto legem non esse positam sed iniustis, gratia autem domini nostri Iesu Christi iustificari pec catores, hoc inquam praeceptum uetus diuinum et pro pheticum commendo tibi ut custodias illibatum et ee clesiis tradas purum et incontaminatum. Vt milites in eis bonam militiam. Allegorice id pronuntiauit, monens episcopum uigiliae et diligentiae. Praeesse e- nim ecclesiae et docere populum est peculiare quod- dam militiae genus, in quo sathan, scelera, falsi fra-  \pend
\vspace{0.5cm}\noindent
\vspace{0.2cm}\rule{1cm}{0.2pt}\\ 
\hspace{0.2cm}\textit{mg}
\footnotesize Colligitur sum̃a doctri nae sanae. 
\normalsize| 
\section*{AD TIM. CAP. I. }
\marginpar{[ p.107 ]}\pstart tres, haeretici hostili exercitu in populum dei mouent. Aduersus horum artes insidias et incursiones uigila- re et populum dei ducere debet episcopus armatus tota illa dei panoplia quam descripsit apostolus ad E- phe. 6.Et Paulus in praesentiarum adiicit in eandem sententiam, Habens fidem et bonam conscientiam. Prius enim monuit finem praecepti esse charitatem ex conscientia bona et fide non simulata. Hanc qui de- bet tradere aliis eius non sit expers oportet. Haec ex- pendant mihi episcopi nostrorum temporum qui ni- hil aliud quam prophanae et maledictae a Deo mili- tiae cruentissimi milites et faces sunt. Discant se sub coelestis et ecclesiasticae militiae signa uocatos duces uerbi et propugnatores populi dei aduersus satha- nam scelera et haereses agere debere. Desinant ergo facere quod faciunt, et faciant quod dominus iubet, aut horribile propediem expectent iudicium. Non igno ro autem haec Pauli uerba, Commendo tibi hoc praece ptum κατὰ τάσ τροαγούσας ἐπ ί σε προφuτεί- αe, id est iuxta ad te praecedentes prophetias, aliter exponi ab aliis, nimirum in hanc sententiam, Sicuti spiritu prophetiae admonitus feci te episcopum, ita diem Christi.  \pend
\section*{COMMENT. IN I. EPIST. }
\phantomsection
\addcontentsline{toc}{subsection}{\textit{Qua repulsa nonnulli circa fi- dem naufragium fecerunt, quorum de numero est Hymenaeus et Ale- xander, quos tradidi sathanae, ut discant non maledicere. }}
\subsection*{\textit{Qua repulsa nonnulli circa fi- dem naufragium fecerunt, quorum de numero est Hymenaeus et Ale- xander, quos tradidi sathanae, ut discant non maledicere. }}\pstart Contrario exemplo ad finem huius causae appen- so ad fidem et diligentiam denique constantiam pasto ralem animat. Congruit autem finis huius capitis cum principio. In principio dixit finem praecepti esse cha- ritatem ex conscientia bona et fide non ficta, a qui- bus quod aberrarint quidam deflexerint ad uanilo- quium adeo turpe ut non amplius intelligant seipsos, neque quid asserant norint. In fine autem dicit Hymenae um et Alexandrum cum nonnullis aliis repulisse fi- dem et charitatem et ideo circa fidem id est uera re ligionis capita fecisse naufragium impegisseque in sco pulos perfidiae. Verisimile autem est Alexandrum hoc ipsum quod Hymenaeum sensisse, de cuius sententia in posteriori epistola scribens, Hymenaeus inquit et Phi- letus circa ueritatem aberrarunt dicentes resurrectio nem iam esse factam, et subuertunt quorundam fi- dem. De qua re dicemus suo loco. Proprie autem alle goriam duxit ab iis qui faciunt naufragium. Tutissi- ma enim nauis in uasto hoc mundi errorum et scele-  \pend
\vspace{0.5cm}\noindent
\vspace{0.2cm}\rule{1cm}{0.2pt}\\ 
\hspace{0.2cm}\textit{mg}
\footnotesize Tuta nauis ueritas. 
\normalsize| 
\section*{AD TIM. CAP. II. }
\marginpar{[ p.108 ]}\pstart rum pelago est ueritas canonica fides pura et chari- tas syncera. In his qui perstant et pura simplicitate ad portum tendunt salutis, salui portum attingunt foe licitatis aeeternae. Porro qui flatibus arrogantiae et con tentionis uela dant abreptique malarum cupiditatum fluctibus a uia recta recedunt, et fidei lucidam cyno- suram negligunt, impingunt in perfidiae scopulos atque  intereunt. Quid uero sit sathanae tradere expositum est 1. Cor.5. cap. Finis huius traditionis est ut pudore probroque correcti, discant ab impiis parumque Christi- anis dogmatis temperare. Hactenus tractauimus nego tium de sana doctrina tradenda et uitandis fabulis, di uersis impiisque dogmatis atque opinionibus.  \pend
\endnumbering\beginnumbering\section{CAP.II.}
\phantomsection
\addcontentsline{toc}{subsection}{\textit{INSTRVITVR COETVS EC- clesiasticus. Adhortor igitur ut ante omnia fiant deprecationes, obsecratio- nes, interpellationes, gratiarum acti- ones pro omnibus hominibus, pro regibus, et omnibus in eminentia constitutis, ut placidam ac quietam uitam degamus cum omni pietate et honestate. }}
\subsection*{\textit{INSTRVITVR COETVS EC- clesiasticus. Adhortor igitur ut ante omnia fiant deprecationes, obsecratio- nes, interpellationes, gratiarum acti- ones pro omnibus hominibus, pro regibus, et omnibus in eminentia constitutis, ut placidam ac quietam uitam degamus cum omni pietate et honestate. }}
\section*{COMMENT. IN I. EPIST. }\pstart Proximum sanae doctrinae quod ecclesia Christi ha bet preces sunt sacrae, sub quas refero Coenam gratu- latoriam sacrificium scilicet laudis et gratiarum acti onis. Et in Actis cap.2.legimus de primitiua ecclesia, Erant autem perseuerantes in doctrina apostolorum et communicatione et fractione panis et precatio- nibus. Baptismi autem pauló ante haec uerba memine- rat Lucas. Proinde satis fuerit ecclesiae sanctae si habe- at doctrinam sanam id est euangelicam, praedicantem in nomine Christi poenitentiam et peccatorum remis sionem, item prophetiam id est lectionem ac expositi onem scripturarum piam, preces denique et coenam mysticam, quibus addimus eleemosynam atque Baptis- mum. De quibus diximus et in 2. et 20.cap. Act. a- posto. et in 1. ad Corin.11.cap. Certe apud uetustissi- mas ecclesias in coetibus sacris plura istis non habuere Christiani ad iustam institutionem ordinationemque ec- clesiasticam necessaria. Temporum successione et pa- storum negligentia factum est ut istorum uestigia uix appareant, de doctrina prophetia precibus sacris et coena mystica loquor. Nam haec corrupit auaritia su- perstitio et innouandi studium. Inde enim irrepsere ab usus missarum, sacrificia pro mortuis, murmur et cantus templorum. Olim in coetibus sacris paucissimae conspiciebantur Ceremontae, et quae erant uel ad do- ctrinam uel ad preces pertinebant. Omnia uero uul-  \pend
\vspace{0.5cm}\noindent
\vspace{0.2cm}\rule{1cm}{0.2pt}\\ 
\hspace{0.2cm}\textit{mg}
\footnotesize Coetus ec- clesiastici ueterum qua les fuerint. 
\normalsize| 
\section*{AD TIM. CAP. II. }
\marginpar{[ p.109 ]}\pstart gari lingua fiebant et recitabantur ad utilitatem om nium. Nihil hic erat uenale. Nullus uel emptor uel uenditor negotiabatur in aede sacra. Nullum erat di- scrimen inter plebem et sacerdotem, nisi quod hic pro ministerio suo et inflitutione sacris praeerat. Et ne quid hac in re dissimulem coetus ecclesiaflici ante annos mil- le fuerunt huiusmodi. Confluebat in aedem sacram cul- tus diuini gratia populus. Dum autem ingrederetur aedem, ab iis qui iam conuenerant Psalmi aliquot in a- liis quidem ecclesiis canebantur in aliis recitabantur tantum, donec uniuersus coetus plene coisset. Hoc ue- ro sacrorum initium ab introeundo appellarunt in- troitum. Iam ubi conuenisset ecclesia acclamabatur ab omnibus unanimi concentu κύριε ἐλέησον kyrie elei- son, Domine miserere. Addebatur a quibusdam Hym- nus qui dicitur angelicus, cuius initium est Gloria in excelsis Deo. Hic est gratulatorius et deprecatorius. Hoc finito collectam recitabat minister aliquis ecclesiae. Ea erat orationis species qua totius ecclesiae uota et necessitates collectae Deo recitabantur. Deinde prae- legebatur a Diaconis eruditioribus locus aliquis siue ex libris propheticis siue literis apostolicis pro tempo ris loci et populi ratione delectus. His absolutis in ae- ditiorem locum per gradus ascendebat episcopus euan- gelium Christi praedicaturus. Plebs interim carmine concepto sancti spiritus gratiam implorabat, quod a  \pend
\section*{COMMENT. IN I. EPIST. }\pstart gradibus Graduale appellari coepit. Hic uero magna cum authoritate praelegebat euangelium episcopus, deinde uero interpretabatur. Ad finem contionis sa- crae recitabat symbolum quod apostolorum uocant uel Nicenum uel Secundo conditum. Praeterea totam ecclesiam inuitabat ad misericordiam, ut quilibet pro suis facultatibus aliquid in aerarium conferret paupe- rum: eam sacri portionem nuncupabant Offertorium. Et Pontius Paulinus docet mensam in ecclesia poni solitam pro pauperibus cibo reficiendis quam et do- mini mensam uocat, et a domino positam. Hac inquam ratione hisce ceremoniis et ritibus exercebant pre- ces et doctrinam sanam ueteres. Attamen hic ritus non per omnia omnibus erat communis. Quidam e- nim sacrum coetum non a Psalmis auspicabantur, sed ab acclamatione κύριε ελέησον. Idem non habebant inusu uel Hymnum angelicum, uel collectas, uel car- men graduale. Huic rursus alii succinebant Hebraeo- rum Alleluiah ut alibi indicat Hieronymus. Apud quosdam episcopus ipse sine istis omnibus sacra con- tione et adiectis publicis precibus conuentum et in cipiebat et dimittebat. Nulla uero ecclesia cogeba- tur in alterius ritus et inftituta iurare, dummodo pre ces et sacrae contiones saluae essent, sancteque exerce- rentur. Et optimae perpetuo quam paucissimis conten tae potissimas impenderunt doctrinae et precibus. Por  \pend
\section*{AD TIM. CAP. II. }
\marginpar{[ p.110 ]}\pstart ro in Coena mystica hic ritus apud omnes penee eccle- sias seruabatur. Prodibat in coetum sacerdos. Instru- cta pane et uino stabat in prospectu populi mensa my- stica, huic iste assistens populo benedicebat dicens, Do- minus uobiscum. Is uero respondebat, Et cum spiritu tuo. Deinde hunc ad res maximas excitans mentesque ́ singulorum praeparans clamabat, Sursum corda. Re- spondebat plebs, Habemus ad dominum. Cyprianus e- nim in sermone de oratione Dominica, Ideo et sacer- dos inquit ante Orationem praefatione praemissa pa- rat fratrum mentes dicendo, Sursum corda, dum respon det plebs, Habemus ad dominum, admoneatur nihil a- liud se quam dominum cogitare debere. Haec Cypria- nus. Hinc uero ad gratiarum actionem inuitans sacer- dos ait, Gratias agamus domino deo nostro. Responde bat plebs, Dignum et iustum est. Hic subiiciebat sacer- dos, Vere dignum et iustum est aequum et salutare, nos tibi semper et ubique gratias agere domine sancte pater omnipotens aeterne Deus per Christum domi- num nostrum. Horum enim pene omnium meminit et Augustinus lib. de Bono perseuerantiae cap.13, dicens, Quod ergo in sacramentis fidelium dicitur, ut sursum corda habeamus ad dominum, munus est domini, de quo munere ipsi domino deo nostro gratias agere a sacerdote post hanc uocem quibus hoc dicitur, admo- nentur, et dignum et iustum esse respondent etc  \pend
\vspace{0.5cm}\noindent
\vspace{0.2cm}\rule{1cm}{0.2pt}\\ 
\hspace{0.2cm}\textit{mg}
\footnotesize Coena my- stica uete- rum. 
\normalsize| 
\section*{COMMENT. IN I. EPIST. }\pstart Post ista autem uerba subiiciebat sacerdos, Qui pridie quam pateretur accepit panem, gratias egit, fregit, de ditque discipulis suis, dicens, Accipite edite, Hoc est cor- pus meum quod pro uobis datur. Et reliqua quae legun- tur in euangelio. His magna cum religione peractis, dicebatur oratio Dominica. Id quod Hieronymus quo que testatur lib. aduersus Pelagianos3. Sed et uulgo receptum est Apostolos tantum ad orationis domini- cae preces consecrasse, hoc est celebrasse coenam my- sticam. Post orationem dominicam sacra mysteria per cipiebat populus, communioneque sacramentorum cor poris et sanguinis dominici in unum corpus coalesce- bat ecclesiasticum cuius caput Christus est. Caeterum rite ad hunc modum peractis omnibus dimittebatur coetus. Ex his autem omnibus facile colligit Lector non omnino stupidus unde abusus Missae papisticae fluxe- rit, et quam nihil hodie antiquae pietatis in hac nisi profunda obsoletaque prope nominum uestigia cernere liceat. Canonis cuius initium est, Te igitur clementissi me pater, hactenus nullam feci mentionem eo quod cer- tum sit hunc sero tandem receptum esse a quibusdam ecclesiis. Et proculdubio melius hodie haberet coetus ecclesiasticus si nunquam fuisset receptus. Verum non est institutum nostrum de his in praesentiarum agere prolixius, cum quaedam et in episto. ad Hebr. annota- rim. Haec in hunc finem hunc in locum annotare libuit  \pend
\section*{AD TIM. CAP. II. }
\marginpar{[ p.111 ]}\pstart ut intelligeret Lector quales fuerint olim coetus eccle- siastici, quodque piae antiquitatis pro se nihil habeant, qui hodie abusum Missarum omnibus ecclesiis ut ritum apostolicum obtrudere uolunt, adeoque et contendunt haereticos esse quotquot Missas suas uulgo receptas non adorant. Imo adiicio his quod illis nouum fore nil ambigo, Gregorium uidelicet eius nominis primum antistitem Romanum annis ab hinc ni fallor noningen- tis, cui et Missas maxime debemus, nullam tamen ec- clesiam ad ritum coëgisse uel suum uel praescriptum alias. Id quod inde liquet quod Augustino episcopo Cantuariensi scribenti percontantique ́, Cum una sit si des cur sunt ecclesiarum diuersae consuetudines et aliter consuetudo Missarum in sancta Romana eccle- sia atque aliter in Galliarum ecclesiis tenetur? Gre- gorius in haec uerba respondet, Nouit fraternitas tua Romanae ecclesiae consuetudinem in qua se meminit nutritum quam ualde amabilem te habeat. Sed mihi placet ut siue in Romana, siue in Galliarum, siue in qualibet ecclesia aliquid inuenisti quod plus omnipo- tenti Deo possit placere sollicite eligas in Anglorum ecclesia quae adhuc ad fidem noua est institutione prae cipua quae de multis ecclesiis colligere potuisti infun- das. Non enim pro locis res sed pro bonis rebus loca amanda sunt. Ex singulis ergo quibusque ecclesiis quae pia quae religiosa quae recta sunt elige et hac qusi  \pend
\vspace{0.5cm}\noindent
\vspace{0.2cm}\rule{1cm}{0.2pt}\\ 
\hspace{0.2cm}\textit{mg}
\footnotesize Olim nemo ad praescri- ptum Missa- rum ritum seruandum coactus. 
\normalsize| 
\section*{COMMENT. IN I. EPIST. }\pstart in fasciculo collecta apud Anglorum mentes in con- suetudinem depone. Haec Gregorius. Ea non profero quod praecipue iis innitar sed ut collatione istorum et nostrorum temporum euincam quam iniquis conditio nibus hodie nobiscum agatur, et quod contra uetus exemplum ritus nobis obtruditur qui non modo nihil utilitatis et pietatis habet sed impietatis et supersti- tionis plurimum. Neglecta igitur improbitate et con uitiis aduersariorum teneamus obsecro ea in coetu san- ctorum sine quibus sanctimonia non perficitur, nen- pe doctrinam sanam, et preces puras assiduasque ́, de doctrina sana disputatum hactenus, de precibus quoq sacris cum alias saepe dictum sit praesentia Pauli uer- ba perstringemus modo. Primum et ante omnia iubet fieri preces, hoc ipso ostendens precibus sacris nil per inde necessarium et utile esse. Praeterea tempus desi- gnat commodum, ut scilicet ipso statim diluculo con- ueniant et communes ad Deum profundant preces. Plinius certe in epistola quadam ad Traianum Impe- ratorem quae extat lib. episto. Plinii 10. meminit moris fuisse Christianis ut stato die ante lucem conuenirent carmenque Christo deo dicerent secum inuicem, seque  sa- cramento aliquo non in scelus aliquod obstringerent sed ne furta latrocinia et adulteria committerent. Et Tertullianus conuentus Christianorum non semel ante lucanos coetus appellat. Quanquam hic nullam statua-  \pend
\vspace{0.5cm}\noindent
\vspace{0.2cm}\rule{1cm}{0.2pt}\\ 
\hspace{0.2cm}\textit{mg}
\footnotesize Anteluca- ni coetus. 
\normalsize| 
\section*{AD TIM. CAP. II. }
\marginpar{[ p.112 ]}\pstart mus necessitatem neque preces certo tempori affigamus. Liberum enim est Christianis et locum et tempus de- signare precibus maxime idoneum. Certe matutinum tempus nemo improbare potest, ut scilicet ieiuni ac so brii simul cum luce Deo nostras preces offeramus et uerbo eius eruditi protinus ad operas et negotiatio- nes nostras digrediamur. Deinde enumerat Paulus cer- tas orationis species, sub his nimirum omnes intelligens. Erasmus ex Graecanicis scholiis tria priora sic discri- minat ut Paulus δέησιν uocarit qua rogamus ut libe- remur a malis et ideo Ambrosium eleganter uertisse deprecationem, προσευχήν autem esse qua preca- mur ut nobis contingant bona, ἔντευξιν uero cum querimur de his qui nos laedunt. Et multo uenustius haec omnia Paraphrasi explicans, Primum deprecen- tur, inquit, Deum ut auertat omnia quae turbant et inquietant statum religionis ac reip. Deinde postulent ab eo quae faciunt ad negotium pietatis et tranquilli tatem reipublicae. Sub haec aduersus eos qui persequun- tur gregem Christi, nihil aliud quam praesidium eius implorent. Postremo pro his quae Dei beneficio conti gerunt agantur gratiae. Habes in his quatuor specie- bus orationis opinor reliquas omnes. Habes item quid et pro quibus orandum sit, quanquam ea ipsa consequen- tibus copiosius exprimantur, nempe pro omnibus ho- minibus. Nam Christianorum est imitari exemplum  \pend
\vspace{0.5cm}\noindent
\vspace{0.2cm}\rule{1cm}{0.2pt}\\ 
\hspace{0.2cm}\textit{mg}
\footnotesize Orationie species. 
\normalsize| 
\hspace{0.2cm}\textit{mg}
\footnotesize Pro qui- bus et quid oran- dum sit, 
\normalsize| 
\section*{COMMENT. IN I. EPIST. }\pstart Dei patris et Iesu Christi filii sui. Ille autem sinit so- lem suum oriri super bonos et malos. Hic uero etiam pro persequutoribus orauit. Gentiles in sacris suis de uouent Christo addictos. At Christiani pro uniuersis hominibus etiam inimicis et peccatoribus orare de- bent. Variae item sunt necessitates hominum, sunt tem- tationes multae, infirmitates quoque cum animi tum cor- poris crebrae quibus affliguntur tum boni cum mali. Pro his omnibus orandus est dominus. Addit, Pro re- gibus et omnibus ἐμ ὑπεροχῆ id est in sublimitate et dignitate constitutis: qui uidelicet in republica pro- ceres sunt dignitateque et functione aliqua ciuili emi- nent. Loquitur autem et de regibus ethihcis et de magistratu a religione Christiana alieno. Ita enim et Ieremias ad eos qui apud idololatras in Babylone te nebantur captiui scribens, unde haec Pauli desum- pta creduntur, praecipit, Et quaerite pacem ciuitatis in quam migrare feci uos illuc et orate pro ea ad domi- num quia in pace eius erit pax uestra. Et Tertullianus ad Scapulam, Sacrificamus inquit pro salute Impera- toris, sed Deo nostro et ipsius, sed quomodo praecepit Deus pura prece. Non enim eget Deus conditor uni- uersitatis odoris aut sanguinis alicuius. Haec enim dae- moniorum pabula sunt. Et in Apologetico cap.39. Coimus, inquit, in coetum et aggregationem ut ad De um quasi manu facta precationibus ambiamus oran-  \pend
\section*{AD TIM. CAP. II. }
\marginpar{[ p.113 ]}\pstart tes. Haec uis deo grata est. Oramus etiam pro impera- toribus pro ministris eorum ac potestatibus, pro sta- tu saeculi, pro rerum quiete, pro mora finis. Item3o.ca pite eiusdem libri, Oramus semper pro omnibus impe- ratoribus, uitam illis prolixam, imperium securum, do- mum tutam, exercitus fortes, senatum fidelem, popu- lum probum, orbem quietum, et quaecunque hominis et Caesaris uota sunt. Sed et ipse Paulus quid oremus significautissime adiiciens ait, Vt placidam siue tran- quillam ac quietam uitam degamus: non equidem in uoluptatibus et impuritate sed in omni pietate siue ue- ro dei cultu genuinaque religione in iustitia καί σεμνό- τητι et honestate. Nam σεμνότοῦς honestatem signi- ficat, uerecundiam, sanctimoniam, et grauitatem, item morum praecipue seueritatem in adolescentia probe instituta. Opponuntur haec seditionibus turbis atque bel- lis quibus labefactantur omnia honesta studia, corrum- puntur mores, perit honestas, soluitur fides et pudi- citia. Potissimum itaque pro tranquillitate publica ora- bit ecclesia, aut si omnino bella pro defensione legum patriae et iustitiae gerenda suscipiuntur a magistrati- bus oret ecclesia dominum ut uictoria suis concessa iu- stam pacem oppressa omni iniustitia restituat. Harum rerum exempla innumera inuenias in Psalmis quibus omnium orantium uota plenissime comprehendun- tur. Habes pro quibus et quid omnino orandum sit.  \pend
\section*{COMMENT. IN I. EPIST. }
\phantomsection
\addcontentsline{toc}{subsection}{\textit{Nam hoc bonum est et acce- ptum coram seruatore nostro Deo, qui cunctos homines uult saluos fieri et ad agnitionem ueritatis ue nire. }}
\subsection*{\textit{Nam hoc bonum est et acce- ptum coram seruatore nostro Deo, qui cunctos homines uult saluos fieri et ad agnitionem ueritatis ue nire. }}\pstart Quia uero preces pro incredulis et ethnicis fusae uideri poterant inutiles contraque decorum et uolun- tatem dei fieri, ideo iam superiorum ueluti causam red- dens addit, Nam hoc est bonum et acceptum coram seruatore nostro Deo, qui cunctos homines uult sal- uos fieri et ad agnitionem ueritatis uenire. Hinc enim sequitur quod nostrum sit pro omnibus orare, nemini ueritatis doctrinam denegare. Qui enim sumus ut re- sistamus uoluntati dei? Vult ille omnes homines saluos fieri, neminem excludit a salute. Vult omnes ad agni- tionem ueritatis uenire. Inde scriptum est in prophe- tis, In omnem terram exiuit sonus eorum et in fines or bis terrae uerba eorum. Cui respondet illud domini praeceptum, Ite in orbem uniuersum et praedicate e- uangelium omni creaturae. Diuersa sunt cultoribus dei designata tempora. Non enim uniuersi prima sta- tim hora ueniunt, quidam sero tandem uineam domini ingrediuntur et eundem cum primis denarium reci- piunt. Ne ergo desperemus aut iudicemus de infideli- bus temere. Solet hoc loco necti quaestio spinosior  \pend
\vspace{0.5cm}\noindent
\vspace{0.2cm}\rule{1cm}{0.2pt}\\ 
\hspace{0.2cm}\textit{mg}
\footnotesize Quaestio praedestina tionis et li- beri arbi- trii. 
\normalsize| 
\section*{AD TIM. CAP. II. }
\marginpar{[ p.114 ]}\pstart quidem quam utilior, Cum deus uelit omnes homines saluos fieri, qui fiat ut non omnes saluentur? Nec igno- ro quid hic a plerisque responsum et quam odiose con- certatum sit maximo etiam cum dispendio et scanda- lo piarum mentium. Alii enim gratiam illam dei qua sola beamur prope ad blasphemiamusque defenderunt, reiicientes in Deum omnium peccatorum culpam. A- lii uero istud offendiculum uitare uolentes huc prola- psi sunt ut et ipsi non citra blasphemiam dei tribue- rint homini quae dei sunt. Vt enim illi nihil relinquunt homini ne peccata quidem: ita hi nihil non tribuunt uiribus hominis. Vtrinque peccatum est et utrisque u- su uenit quod prouerbio dicitur, Incidit in Scyllam cu piens uitare Charybdim. Multo prudentius et religi osius uidentur mihi hanc quaestionem expediisse uete- res, qui de ipsa aduersum Pelagianos disputarunt co- piosissime, atque ita ut quaecunque hominis sunt relique rint homini, quae uero dei sunt Deo. D. Ambrosius si- ue alius quispiam uir sanctus et doctus de hac re aedi- dit libros duos inscriptos De uocatione omnium Gen- tium, in iis tandem negotium totum in hoc contraxit compendium, remotis abdicatisque omnibus concerta- tionibus quas intemperantium disputationum gignit animositas tria esse perspicuum est quibus in hac quae- stione debeat inhaereri. Vnum quod profitendum est, Deum uelle omnes homines saluos fieri et in agnitio  \pend
\section*{COMMENT. IN I. EPIST. }\pstart nem ueritatis uenire. Alterum quod dubitandum non est ad ipsam cognitionem ueritatis et perceptionem salutis non suis quenquam meritis sed ope atque opera diuinae gratiae peruenire. Tertium quod confitendum est Altitudinem iudiciorum dei humanae intelligentiae penetrabilem esse non posse, et cur non omnes saluet qui omnes uult saluos fieri non oportere disquiri. Quo niam si quod cognosci non potest non quaeratur inter primam et secundam definitionem non remanebit cau sa certaminis, sed secura ac tranquilla fide utrumque prae dicabitur utrumque credetur. Deus quippe apud quem non est iniquitas et cuius uniuersae uiae misericordia et ueritas omnium hominum bonus conditor iustus est ordinator, neminem indebite damnans neminem in debite liberans, nostra plectens cum punit noxios, sua tribuens cum facit iustos ut obstruatur os loquentium iniqua et iustificetur Deus in sermonibus suis et uin- cat cum iudicatur. Nec damnati iusta quaere monia nec iustificati uerax est arrogantia, si uel ille dicat non me- ruisse se poenam uel iste asserat meruisse se gratiam. Congruunt his Ambrosianis Augustiniana quae le- guntur lib. de Bono perseuerantiae cap.19. Hic tamen non negarim rationem commodiorem inueniri posse qua hic nodus dissoluatur. Neque hoc dissimulat Aurel. August. cuius uerba quod mire huc faciant ex 53. Tra cta. in Ioannem nunc ascribo. Si quis istam quaesionem  \pend
\section*{AD TIM. CAP. II. }
\marginpar{[ p.115 ]}\pstart inquit liquidius et melius nouit se posse et confidit exponere, absit ut non sim paratior discere quam do- cere, tantum ne audeat quisquam liberum arbitrium sic defendere ut nobis orationem qua dicimus, Ne nos inferas in tentationem, conetur auferre: rursus ne quis neget uoluntatis arbitrium et audeat excusare pecca- tum. Sed audiamus dominum et praecipientem et o- pitulantem et iubentem quid facere debeamus, et adiuuantem ut implere possimus. Nam et quosdam ni- mia uoluntatis suae fidutia extulit in superbiam, et quos dam nimia uoluntatis suae diffidentia deiecit in negli- gentiam. Illi dicunt, Vt quid rogamus Deum ne uinca mur tentatione quod in nostra est potestate? isti di- cunt, Vt quid conamur bene uiuere quod in dei est po testate? O domine ô pater qui es in coelis ne nos inferas in quamlibet istarum tentationum, sed libera nos a malo. Audiamus dominum dicentem, Rogaui prote Petre ne deficiat fides tua. Nec sic existimemus fidem nostram esse in libero arbitrio ut diuino non egeat ad iutorio. Audiamus et euangelistam dicentem, Dedit eis potestatem filios dei fieri ne omnino existimemus in nostra potestate non esse quod credimus uerum in utroque illius beneficia cognoscamus. Nam et agendae sunt gratiae quia data est potestas et orandum ne suc- cumbat infirmitas. Ipsa est fides quae per dilectionem operatur sicut eius mensuram dominus cuique partitus  \pend
\section*{COMMENT. IN I. EPIST. }\pstart est, ut qui gloriatur non in seipso sed in domino glorie- tur. Hactenus Augustini uerba recensui. Ex quibus o- mnibus liquet quod bonam illam dei uoluntatem atti- net ipsum quidem uelle cunctos homines saluos fieri et ad agnitionem ueritatis uenire, et proinde quod alii ueniunt non esse humani meriti sed diuinae uolun- tatis et gratiae, quod uero alii non ueniunt humanae malitie et corruptionis. Salus igitur et iustificatio non erit neque uolentis neque currentis sed miserentis dei, credentque quotquot ordinati sunt ad uitam aeternam. Sunt autem ad aeternamuitam ordinati quicunque in Christum credunt. Deus enim omnia agit lege certa et ordine iusto. Elegit nos antequam iacerentur funda- menta mundi, sed in Christo et ad sanctimoniam. ltaque  qui Christo fidunt et student sanctimoniae electi et ad uitam ordinati sunt. Electio enim dei a uoluntate non est separanda neque sola praeter ordinem et con- tra media est iactanda. Petrus enim electis scribit dei secundum praefinitionem dei patris per sanctificatio- nem spiritus in obedientiam et aspersionem sanguinis Iesu Christi. Et Paulus, Quos inquit praesciuit Deus eos et praefiniuit ut conformes fiant imagini filii eius, etc. lam cum constet non omnes conformes fieri ima- gini filii eius, colligunt alii non omnes ad uitam esse praescitos, proinde tropo hunc Pauli locum exponunt Deum uelle omnes credentes saluos facere. Quod et  \pend
\section*{AD TIM. CAP. II. }
\marginpar{[ p.116 ]}\pstart apud Augustinum contra Pelagianos inuenies, et pla- ne uerum est. Impii enim non saluantur sed pii tantum. Caeterum illi non habent quod Deo qui pro bonitate sua uoluit eos ad agnitionem ueritatis uenire, suae dan- nationis et impietatis culpam impingant. Non enim sequitur praescientiam dei nos ad res impias affige- reut culpam peccati in Deum reiicere possimus. Re- cte enim D. Augustinus Tracta. in Ioan.53. Non pro- pterea quenquam deus ad peccandum cogit quia fu- tura hominum peccata iam nouit. Ipsorum enim prae- sciuit peceata, non sua non cuiusquam alterius sed ipso- rum. Fecerunt ergo peccatum Iudaei, quod eos non compulit facere, cui peccatum non placet, sed facturos esse praedixit quem nihil latet. Et ludas non ideo pro- didit dominum quod ita de eo scriptum esset, sed ideo de eo scriptum erat quod deus praesciret auaritia et desperatione animum ipsius fore corrumpendum. Ob- iicis, Atqui ratione pari consequetur ut quia uiderit deus fidelium bona opera praedestinarit eos ad uitam aeternam, unde consequatur electionem dei pendere a nostris meritis. Nequaquam. Nam mala quae ex corru- ptione et labe primaria uere nostra sunt puniuntur in nobis iustitia dei: sed gratia dei electi sumus ad uitam uon propter opera nostra, sed propter sanguinem Christi. 2. Timoth. 1. Haec in praesenti ex uetustis scri- ptoribus praecipuisque religionis nostrae antistitibus pro  \pend
\section*{COMMENT. IN I. EPIST. }\pstart ferre libuit. Quod si alii melius quiddam reuelatum fu- erit certe cum Augustino discere quam docere malo, dummodo Deus non blasphemetur et in illum omnis cul- pa malitiae et scelerum nostrorum non reiiciatur, aut nostris uiribus tribuatur quod reuera illius gratiae est. Dominus enim apud Oseam planis uerbis pronuuti- at, Perditio tua Israel, sed in me est auxilium tuum. Et in euangelio dicit Christus, Hierusalem Hierusalem, quoties uolui te congregare sicut gallina congregat pullos suos sub alas, et tu noluisti. Item Daniel, Tibi domine inquit iustitia, nobis autem confusio facierum. Et Dauid, Ne intres in iudicium cum seruo tuo quia non iustificabitur in conspectu tuo omnis uiuens. Non nobis domine non nobis, sed nomini tuo da gloriam. Iu stus est dominus in omnibus uiis suis, et sanctus in o- mnibus operibus suis. Et Paulus, Si adhuc uelatum est euangelium nostrum in his qui pereunt uelatum est in quibus Deus huius saeculi excaecauit sensus incredu lorum ne illucesceret illis lumen euangelii. Vt ergo u- no uerbo dicam omnia, Lux et salus a sola dei gratia est, Excaecatio et damnatio a corruptione nostra. Quod uero deus interim excaecare dicitur iustitiae e- ius est, secundum quam merito punit inobedientiam et perfidiam nostram induratione desperatione et caecitate. Ipsi soli gloria in saecula, Amen.  \pend
\section*{AD TIM. CAP. II. }
\marginpar{[ p.117 ]}
\phantomsection
\addcontentsline{toc}{subsection}{\textit{Vnus enim Deus, unus etiam mediator dei et hominum homo Christus lesus, ꝗ dedit semetipsum precium redemptionis pro omni bus, ut esset testimonium tempo- ribus suis in quod positus sum e- go praeco et apostolus? ueritatem dico in Christo non mentior do- ctor gentium cum fide et ueritate. }}
\subsection*{\textit{Vnus enim Deus, unus etiam mediator dei et hominum homo Christus lesus, ꝗ dedit semetipsum precium redemptionis pro omni bus, ut esset testimonium tempo- ribus suis in quod positus sum e- go praeco et apostolus? ueritatem dico in Christo non mentior do- ctor gentium cum fide et ueritate. }}\pstart Ad confirmationem superiorum etiam haec perti- uent. Nam si unus Deus est qui neque Iudaeorum Deus tantum est, neque huius uel illius gentis peculiaris, sed omnium ex aequo communis, si unus totius carnis me- diator est inter Deum et hominem, siis semetipsum precium redemptionis impendit pro omnibus, profe- cto orandum nobis erit pro omnibus. Quando quidem uero insignis inter discipulos contrauersia fuerat de eo num et gentibus natus esset Christus, et an Deus Gentium quoque Deus esset, ideo uocationem et mini sterium suum hic praedicat apostolus, ne quis testimo- nio eius, quo testabatur Christum omnibus redimen- dis uenisse et Deum non Iudaeorum tantum sed et Gentium esse Deum, diffideret. Haec in summa tractat  \pend
\section*{COMMENT. IN I. EPIST. }\pstart praesens Pauli locus, iam ut singula repetamus et ex- plicemus diligentius, Vnus inquit est Deus, quo non excludit uel spiritum sanctum uel filium qui non sepa- ratur a substautia patris, substantia autem unus est deus, persona trinus, unitas ergo illa nontam ad sub- stantiam dei refertur quam ad homines, qui unum id est communem habent Deum. Ita etiam unus id est con- munis omnium mediator est Christus Iesus, ut hunc ne- mo sibi peculiariter uendicare possit. Non enim pro una aliqua gente, sed pro omnibus sese impendit. Ve- rum uidentur haẽc quod ipsauis euangelii in ipsis late- at altius esse repetenda. Mediator non est inter consen- tientes sed dissidentes. Qui enim concordant nullo o- pus habent conciliatore. Et mediator et conciliator equipollent, ubi itaque mediator est ibi dissidium atque  partes sunt. Et in praesentiarum partes sunt Deus et Homo. Hae inter se dissident. Causa uero dissidii est pec- catum. At hoc non expiatur citra mortem et sangui- nem. Oportet itaque huius dissidii mediatorem cum ho- minibus participare carne et sanguine adeoque et ue- rum hominem esse qui mori et sanguinem effundere possit. Iam si nihil aliud quam homo sit et sanguis e- ius nihil differat a reliquo corruptorum hominum san- guine, non expiabit peccatum dissidii causam, imon ni- si Deus quoque sit hic mediator ad deum nequaquam ac- cedere poterit. Oportet enim mediatorem utrique par-  \pend
\vspace{0.5cm}\noindent
\vspace{0.2cm}\rule{1cm}{0.2pt}\\ 
\hspace{0.2cm}\textit{mg}
\footnotesize Vnus deus 
\normalsize| 
\hspace{0.2cm}\textit{mg}
\footnotesize Vnus me- diator. 
\normalsize| 
\section*{AD TIM. CAP. II. }
\marginpar{[ p.118 ]}\pstart ti communem esse. Vere enim Chrysostomus, Homo inquit purus mediator nunquam penitus fieret. Opor- tebat enim huiusmodi mediatorem cum Deo colloqui. Deus item solus mediator esse non posset. Neque enim suscepissent eum hi quorum mediator accederet. Pro inde diuino consilio uerbum caro factum est. Et Pau- lus in praesenti non citra emphasim dixit, Homo Chri- stus Iesus. Quasi dicat, In hoc hominem induit Chri- stus Iesus ut mediatorem agere posset. Verus ergo De us et uerus homo est Christus. Hic dedit semetipsum praecium redemptionis pro omnibus. Semetipsum in- quam. Non enim per alium offerri potuit. sed per spi- ritum aeternum seipsum obtulit immaculatum deo. He- brae.7.9. Et propheta, Oblatus est inquit quoniam ipse uoluit. Denique ipse dominus, Panis quem ego dabo ait caro mea est quam ego impendam pro mundi uita. Et uerbum quo hic utitur apostolus ἀντίλυτρον id si- gnificat praecium quo redimuntur captiui ab hostibus, eamque commutationem qua capite caput uita redimi- tur uita. Ita ergo impendit se pro nobis dominus ut sua morte nos uiuificaret, suo sanguine expiaret, et ex maledictis faceret benedictos dei filios. Neminem ex- clusit hic nisi qui sese excludit sua ingratitudine et perfidia. Hinc Paulus Impendit se ait pro omnibus. Quod autem alii hinc colligunt omnes etiam perfidis- fimos fore saluandos, petulanter blasphemant. Qui e-  \pend
\section*{COMMENT. IN I. EPIST. }\pstart nim eredit habet uitam aeternam, qui non credit dan- nabitur. Hoc quidem ab aeterno praeuidit Deus, uo- luit autem praedicari suis temporibus. De quibus et prophetae, et Petrus 2. Pet. 2.Et Paulum delegit ut omni- bus praedicaret salutem. Id uero iuramenti quadam spe- cie confirmans addit, Veritatem dico per Christum, non mentior quod Deus uidelicet me delegit ut sim Do- ctor omnium gentium, ut omnibus annuntiem gratiam euangelicam, idque  fideliter et uere. Haec uero eximia Pauli laus est.  \pend
\phantomsection
\addcontentsline{toc}{subsection}{\textit{Volo igitur orare uiros in o- mni loco, sustollentes puras ma- nus absque iracundia et disceptatione. }}
\subsection*{\textit{Volo igitur orare uiros in o- mni loco, sustollentes puras ma- nus absque iracundia et disceptatione. }}\pstart Redit ad preces unde paulo ante fuerat digressus, tollitque loci superstitionem et ostendit qualis debeat esse oratio. Quasi dicat, Quod iubeo conuenire in coe- tum atque in hoc publicas profundere preces non est quod inde quis colligat preces nullibi quam in conuen- tu efficaces esse, imo uolo orare uiros in omni loco u- bicunque res postulat. Omnis enim locus purus est si a- nimus non sit impurus. Et oratio ex mente non ex loco estimatur. Legitur Christus saepius in precibus perno- ctasse sane non in templo, frequentius orauit in mon- te. Tollitur ergo hisce Pauli uerbis loci superstitio, ut  \pend
\vspace{0.5cm}\noindent
\vspace{0.2cm}\rule{1cm}{0.2pt}\\ 
\hspace{0.2cm}\textit{mg}
\footnotesize De oratio- nis loco. 
\normalsize| 
\section*{AD TIM. CAP. II. }
\marginpar{[ p.119 ]}\pstart et apud Matthaeum in 6. et Ioan.in 4. Porro inde non sequitur, quod tamen uaesani quidam colligunt, in templo non esse orandum. Non enim temere pro- nuntiauit dominus, Domus patris mei domus orationis est. Et Paulus in templo non impie legitur in Actis o- rasse. Quod si in templo et aede sacra hoc est doctri- nae precibus et sacramentis dicata orare non licet, cer- te oratio erit loco affixa. Atqui libera est, licet itaque et in aede sacra et in quouis loco ubicunque res et deco- rum ipsaque necessitas postulat orare. Audiamus nunc qualis esse debeat oratio. Volo inquit uiros orare su- stollentes puras manus absque ira et disceptatione. Cae- terum exteriori orantis habitu adumbrat qualis esse debeat pie precantis animus. Eleuantur manus ad coe- lum, eleuetur et mens. Obliuiscamur rerum terrena rum sola coelestia cogitemus. Proferimus et excutimus manus quoties palam testari uolumus nos a furto et rapina esse quam alienissimos, orantes igitur excu- tiamus ex pectore omnem iniustitiam crudelitatem et impuritatem, oremus autem mente pura et syncera. Nam apud Isaiam in cap.2.clamat ad Israelem domi- nus, Cum extenderitis manus uestras abscondam ocu- los meos a uobis et cum multiplicaueritis orationem non exaudiam, eo quod manus uestrae plenae sint san- guinibus. Lauamini igitur et mundi estote et tollite malitiam conatuum uestrorum e conspectu meo et  \pend
\vspace{0.5cm}\noindent
\vspace{0.2cm}\rule{1cm}{0.2pt}\\ 
\hspace{0.2cm}\textit{mg}
\footnotesize Qualis esse debeat ora- tio. 
\normalsize| 
\section*{COMMENT. IN I. EPIST. }\pstart cessate malefacere. Apud Ieremiam in cap.n.dicit do- minus, in singulis ciuitatibus et plateis posuerunt a- ras ad offerendum Baal. Tu igitur noli orare pro po- pulo hoc, quoniam ego non exaudiam. Huc pertinet quod Paulus adiungit χωρις ὀργῆς καί διαλογισ- μoῦ sine ira et haesitatione siue disceptatione. Postre- mum illud referunt quidam ad fidutiam inconstantio rem et uacillantem, ut hic locus congruat cum illo Ia- cobi, Postulet cum fidutia nihil haesitans, etc. Alii re- tulerunt ad animum implicatiorem odio inuidiaque ui- tiatum. Certe congruunt huc quae dominus apud Mat- thaeum hisce uerbis docuit. Si obtuleris munus tuum ad aram, et ibi recordatus fueris quod tibi cum fratre male conuenit, relinque munus, abi et reconciliare fratri tuo. Munus nostrum oratio est. Pura enim prece uult sibi litari dominus. Ara autem nostra cor nostrum est. Si igitur constituamus apud nos precari, tollamus prius ex animo quicquid inuidiae et malitiae contra fratrem concepimus, ut ex pura bonaque ; conscientia dicere queamus, Dimitte nobis peccata uostra sicut et nos remittimus debitoribus nostris. Tertullianus in 3o. cap. Apolog. describens exteriori quoque orantium ha- bitu mentis iustam compositionem. Ad coelum inquit sufpicientes Christiani manibus expansis, quia inno- cuis, capite nudo, quia non erubescimus, deniq sine monitore, quia de pectore oramus precantes sumus  \pend
\section*{AD TIM. CAP. II. }
\marginpar{[ p.120 ]}\pstart omnes semper pro omnibus. Et paulo post, Offerimus deo opimam hostiam orationem de carne pudica, de a- nima innocenti, de spiritu sancto profectam. Haec imi- temur obsecro qui hodie Christiani uideri uolumus et inuicem batalogia lacessimus. Non enim in uerbo- rum copia non in numero, sed in pietate et fide in assi duitate et charitate consistit efficatia et uirtus o- rationis.  \pend
\phantomsection
\addcontentsline{toc}{subsection}{\textit{Consimiliter et mulieres in a- mictu modesto cum uerecundia et castitate ornantes semetipsas, non tortis crinibus, aut auro, aut mar- garitis aut uestitu sumptuoso, sed quod decet mulieres profitentes pietatem per opera bona. }}
\subsection*{\textit{Consimiliter et mulieres in a- mictu modesto cum uerecundia et castitate ornantes semetipsas, non tortis crinibus, aut auro, aut mar- garitis aut uestitu sumptuoso, sed quod decet mulieres profitentes pietatem per opera bona. }}\pstart Nihil in hunc locum adfertur a receptis interpre- tibus quod non in breuem summam contraxerit Pa- raphrastes, nec quicquam dici potest necessarium quod idem non sit hisce uerbis complexus, quae hic malui a- scribere quam post planam illam expositionem pro- lixius cornicari. Ad exemplum uirorum inquit pre- centur et mulieres. Si quid est in animo muliebrium affectuum id prius eiiciant, pro ludaicis purgationi- bus morum innocentiam adferentes, ad hanc uicti-  \pend
\vspace{0.5cm}\noindent
\vspace{0.2cm}\rule{1cm}{0.2pt}\\ 
\hspace{0.2cm}\textit{mg}
\footnotesize Mulierum ornatus et habitus in coetu sacro. 
\normalsize| 
\section*{COMMENT. IN I. EPIST. }\pstart mam animi non corporis cultum adhibeant, non nudi- tate corporis uirorum oculos ad libidinem sollicitan- tes, sed amictu tectae et eo amictu qui modestiam qui pudorem qui pudicitiam prae se ferat. Absit autem ut Christianae mulieres eo cultu prodeant in coetum sa- crum quo uulgus prophanarum mulierum solet ad nuptias aut theatrum exire, quae se multo studio pri- us ornant ad speculum arte contortis crinibus, aut auro intertexto, aut pendulis ab auribus collonue mar- garitis, aut alioqui holoserica purpureaue ueste quo simul et formam spectatoribus lenocinio commendent et suas ostentantes opes tenuioribus suam inopiam exprobrent. Quin is sit potius Christianarum habitus qui uitae respondeat qui dignus uideri possit iis mulie- ribus quae ueram pietatem ac dei cultum profiteantur, non ostentatione diuitiarum sed benefactis quibus opi- bus unice delectatur Deus, cuius oculis foedum est quod mundo uidetur splendidum et magnificum. Haec paraphrastes. Paria propemodum reperies in 2.ad Co- rinth.12.et 1.Pet.3. cap. Tertullianus scripsit libel- lum de Cultu foeminarum eruditissimum et utilissi- mum, quem hisce uerbis mire uidelicet ad praesentia facientibus concludit, Prodite uos iam medicamentis et ornamentis extructae apostolorum, sumentes de simplicitate candorem, de pudicitia ruborem, depictae oculos uerecundia et spiritus taciturnitate, inseren-  \pend
\section*{AD TIM. CAP. II. 120 }\pstart tes in aurem sermonem dei, annectentes ceruicibus iu- gum Iesu Christi. Caput maritis subiicite et satis orna- tae eritis. Manus lanis occupate, pedes domi figite et plus quam in auro placebunt. Vestite uos serico pro- bitatis, byssino sanctitatis, purpura pudicitiae. Taliter pigmentatae Deum habebitis amatorem.  \pend
\phantomsection
\addcontentsline{toc}{subsection}{\textit{Mulier in silentio discat cum o- mni subiectione. Caeterum mulie- ri docere non permitto, neque autho- ritatem usurpare in uiros sed esse in silentio. Adam enim prior for- matus est, deinde Eua, et Adam non fuit deceptus, sed mulier sedu- cta fuit per praeuaricationem: sal- ua tamen fiet per generationem li- berorum, si manserit in fide ac di- lectione et sanctificatione cum ca- stitate. }}
\subsection*{\textit{Mulier in silentio discat cum o- mni subiectione. Caeterum mulie- ri docere non permitto, neque autho- ritatem usurpare in uiros sed esse in silentio. Adam enim prior for- matus est, deinde Eua, et Adam non fuit deceptus, sed mulier sedu- cta fuit per praeuaricationem: sal- ua tamen fiet per generationem li- berorum, si manserit in fide ac di- lectione et sanctificatione cum ca- stitate. }}\pstart Quo facilius a mulierculis quod praescripserat im- petraret Legem dei profert admonens illas sui officii et humilitatis. Lege autem dei praecipitur mulieribus obedientia silentium et humilitas. Est autem hoc ge-  \pend
\vspace{0.5cm}\noindent
\vspace{0.2cm}\rule{1cm}{0.2pt}\\ 
\hspace{0.2cm}\textit{mg}
\footnotesize Lex de mu- lierum di- sciplina. 
\normalsize| 
\section*{COMMENT. IN I. EPIST. }\pstart nus hominum natura garrulum fastuosum et recalcitrans. Hinc Paulus, Mulier inquit in silentio discat cum omni subiectione. Id est mulieres silere et per omnia subes- se et a uirorum authoritate pendere debent. Sequi- tur enim huius legis expositio, Mulieri docere non per mitto, neque ; ἄυθεντεῖν ἀνδρός id est authoritate uti in uiros et illos ad suum cogere affectum legesue prae- scribere authenticas, sed omnino esse in silentio. Hanc legem repetiit et in 1. ad Corinth.14. cap. et trans- sumpsit ex Gene. 3. cap. ubi in haec uerba pronuntiat dominus, Multiplicando multiplicabo dolorem tuum et conceptum tuum. In dolore paries filios, et ad ui- rum tuum erit desyderium tuum, et ipse dominabi- tur tibi. Id est nihil constitues priuato affectu sed ad nu- tum uiri omnia ages. Proinde cum in rebus priuatis re- uerenter debeant auscultare mulieres, multo minus publicis functionibus praefici debent. Inuertunt enim naturae et dei ordinem qui in iudiciis exercendis inque ́ dicendis sententiis, in ferendis item legibus, denique  et in ministeriis et functionibus ecclesiasticis ullum locum permittunt mulieribus. Vnde consequens est et humanarum et diuinarum rerum ex aequo fu- isse ignaros qui Abbatissas (sic enim uocarunt princi- pes inter uirgines Vestales, in quarum numero ta- men non deerant scorta Corinthiaca) praefecerunt rei pub. legibus adde et sacris functionibus. Causas di-  \pend
\section*{AD TIM. }
\marginpar{[ p.121 ]}\pstart uinae legis subiicit apostolus ne quis muliercularum causam agens quereretur istam ex audacia et libidi- ne uirorum esse promulgatam. Prior causa est, Adam prior formatus est deinde Eua, adeoque ex uiro formata est foemina, adde quod in adiutorium uiri a Deo con- dita, ut portio uerius uiri sit quam caput:non e capite sed latere uiri sumpta. Posterior causa est, Adam non fuit deceptus sed mulier seducta fuit. Vir enim neque  serpentis promissionibus, neque pomi illecebra decipi potuisset: sola charitas uxoris pertraxit ad pernitio- sum obsequium. Denique Adam homo a quo nihil hu- mani alienum non saxum rigidum erat, ut nihil dicant qui obiiciunt uirum debuisse constantiorem fuisse quam quem illecebrae fregissent uxoris. Fons ergo praeuaricationis a foemina existit. Hinc non temere ad mulierem Tertullianus de habitu muliebri, Tu inquit es diaboli ianua, tu es arboris illius resignatrix, tu es diuinae legis prima desertrix, tu es quae eum suasisti quem diabolus aggredi non ualuit, tu imaginem dei ho- minem tam facile elisisti, propter tuum meritum id est mortem etiam filius dei mori habuit etc. Quando er- go magisterium Euae non uni Adae sed toti mortalium generi semel cessit pessime, subesse posthac et iure ad uirum respicere iubetur mulier quaelibet. Agnoscat imbecillitatem et lapsum ueterem, ac satis habeat quod quae olim fuit dux ad impietatem et ruinam  \pend
\section*{COMMENT. IN I. EPIST. }\pstart nunc erroris memor sequatur uiri pietatem et resur- gat in pia obedientia. Non enim haec dicuntur in con- tumeliam omnium mulierum, neque ideo excluduntur a consortio sanctorum, nec quisquam putet sibi ideo in istas licere quidlibet. Imo cum muliercularum ge- nus alias sit timidulum, ne in desperationem prorsus praeceps rueret, addit nunc consolationem, dicens, Sal- ua tamen fiet per generationem liberorum etc. Atqui certum est salutem non deberi ullis meritis nostris, tantum abest ut debeatur procreationi liberorum. So- la enim dei gratia per Christum saluamur. Abutitur ergo saluandi uerbo pro reparandi, quemadmodum et Germanis mos est usurpare uerbum heylen für wi- derbzingen vnd ersetzen. Quasi dicat, Audistis quantam muliebris sexus sibi inusserit maculam, sed ea eluetur si non desit suo officio. Officium autem mu- lierum est liberos gignere, labores et dolores quos secum trahit partus et educatio liberorum pacienter tolerare, liberis omnem operam et diligentiam im- pendere, fidem item cum erga deum tum maritum ser- uare, utrumque diligere ex animo. Sic autem fiet ut excus sa omni superstitione cui tamen hoc genus maxime ob noxium est uere et indiuulse uera fide Deo et Christo adhaereant, fidem quoque coniugalem non uiolent, deum item ante omnia, mox maritos diligant quemadmo- dum ecclesia solet diligere Christum sponsum suum. Huc  \pend
\vspace{0.5cm}\noindent
\vspace{0.2cm}\rule{1cm}{0.2pt}\\ 
\hspace{0.2cm}\textit{mg}
\footnotesize Gene.3. 
\normalsize| 
\section*{AD TIM. CAP. II. }
\marginpar{[ p.122 ]}\pstart pertinet sactimonia et castitas siue pudicitia et so- brietas. Debent enim mulieres animo et corpore inte- grae et sanctae esse, deinde castae et sobriae in dictis fa- ctis moribus denique toto corporis gestu et habitu, Ti- tum 2. Mulier enim impudica monstrum uerius hominis quam homo est. Hisce autem caeterisque uirtutibus sarcient inquit Paulus, quod in primo illo lapsu diminutum est, quod scilicet hominis partes et officia attinet, alias ne- mo negare potest omne reparationis precium et meritum passioni dominicae deberi. Doceant ergo episcopi san- cti quod non puduit sanctissimum Christi apostolum Pau- lum docere. Quidam non sustinentes doctrinam sanam ad fa- bulas conuertuntur et nugantur de uita monastica ac uir- ginitate quam interim sanctissimam deoque gratissimam esse credimus, si modo eiusmodi sit qualem describit a- postolus 1. Cor. 7. Caeterum coactam illam et inuoluntari- am adeoque et immundam non damnare non possumus, quod certo sciamus deo probari placereque sanctimoni- am cum animi tum corporis. Vnde et apostolus praecepit corpora nostra possideamus in sanctimonia et hono- re et non in concupiscentia, alibi quoque commendans connubium sanctum et immunditiam damnans. Honora bile ait est comnubium apud omnes et cubile impollutum, scortatores autem et adulteros iudicabit deus.  \pend
\phantomsection
\addcontentsline{toc}{subsection}{\textit{DESCRIPTIO EPISCOPI ET familiae eius. }}
\subsection*{\textit{DESCRIPTIO EPISCOPI ET familiae eius. }}
\vspace{0.5cm}\noindent
\vspace{0.2cm}\rule{1cm}{0.2pt}\\ 
\hspace{0.2cm}\textit{mg}
\footnotesize Coelibatus monastico- rum. 
\normalsize| 
\section*{COMMENT. IN I. EPIST. }\pstart In ecclesia dei hoc est sanctorum coetu ueluti anima est minister uerbi, utpote quem Paulus alibi appellat co- oparium dei, et per cuius ministerium deus populum suum gu- bernat. Est igitur hic a deo constitutus ut ab ore eius cum- cti expectent et audiant uerbum dei, idem a domino ueluti lu- cerna quae piam toti domui est expositus ad quem intenden- tes qui uersantur in domo quid agant norint. Nam in euangelio dominus ad uerbi ministros, Vos, inquit, e- stis lux mundi. Sic luceat lux uestra coram hominibus ut uideant uestra bona opera glorificentque patrem qui in coelis est. Proinde Paulus uitam et mores Episcopi informat ne ex his ullae in ecclesiam offundantur tenebrae.  \pend
\endnumbering\beginnumbering\section{CAP. III.}
\phantomsection
\addcontentsline{toc}{subsection}{\textit{Indubitatus sermo, Si quis epi- scopi munus appetit bonum opus desiderat. Oportet igitur episco- pum irreprehensibilem esse, unius uxoris maritum, uigilantem, sobri um, modestum, hospitalem, aptum ad docendum, non uinolentum, non percussorem, non turpiter lu- cri cupidum, sed aequum, alienum a pugnis, alienum ab auaritia, }}
\subsection*{\textit{Indubitatus sermo, Si quis epi- scopi munus appetit bonum opus desiderat. Oportet igitur episco- pum irreprehensibilem esse, unius uxoris maritum, uigilantem, sobri um, modestum, hospitalem, aptum ad docendum, non uinolentum, non percussorem, non turpiter lu- cri cupidum, sed aequum, alienum a pugnis, alienum ab auaritia, }}\pstart Chrysostomus et huius abbreuiator Theophyla- ctus hanc particulam πιςτὸς ὁ λόγος retulerunt ad  \pend
\section*{AD TIM. CAP. III. }
\marginpar{[ p.123 ]}\pstart superiora quae Paulus uoluerit indubitata esse, quod uidelicet dixit mulierem seruari gignendis et recte in stituendis liberis. Videtur autem referenda esse ad con sequentem Episcopi descriptionem ut sit occupatio quaedam. Ambiebant enim quidam hanc functionem ob dignitatem ac quaestum et non uocati se ingere- bant, clamantes honestum esse quod appeterent. Hic uero Apostolus Episcopi munus sane (inquit) nemo ne scit honestum esse munus, et qui hoc concupiscit ut prosit, honesto quidem illum desiderio flagrare. Hic uero addendum quod inabsolutam sententiam perfi- ciat, Verum ad hanc functionem non quilibet sed se- lectissimus quisque appositus est:neque omnes illam appe- tunt ut prosint. Quidam enim quaestum, honores, et tyrannidem spectant. At Episcopus non tam dignita- tis quam officii ac ministerii beneficiique nomen est. Quae si uera sunt ut falsa esse nequeunt quae spiritussan ctus prodidit, concidit hic omnis episcoporum Ponti ficiorum authoritas, damnati sunt praeterea ambitus quotquot ex doctis hodie quoque non uocati ob hono- res et opes assequendas huic sancto se dei ministerio ingerunt. Nec est quod illi nobis quoque functionis pre cium et honestatem obiiciant, curent illi potius ut iu sta uia in hoc ministerium uocentur et ingrediantur, ut interim sancti sint, ut uni Deo et uerae pietati con secrati, inque lege domini exercitati uocationem iustam  \pend
\vspace{0.5cm}\noindent
\vspace{0.2cm}\rule{1cm}{0.2pt}\\ 
\hspace{0.2cm}\textit{mg}
\footnotesize Contra am bitum. 
\normalsize| 
\section*{COMMENT. IN I. EPIST. }\pstart praestolentur, utque uocati iam sanctos et Deo dignos agere possint episcopos. Haec qui faciunt studiosi bo- num deoque gratissimum opus faciunt. Atque hunc sco- pum praefixum sibi habebunt quotquot hodie ex ec- clesiae aerario ad ministeria ecclesiastica aluntur. De nomenclatura episcopi annotauimus quaedam in 2o. caput Act. aposto. et in alia nonnulla Pauli loca. Se- quitur nunc simulachrum ueri episcopi, quod aposto- lus suis attributis hoc est coloribus egregie pinxit. Im primis iubet episcopum esse ἀνεπίλντομ irreprehen sibilem, hoc est qui tam sit uitae inculpatae ut nullo iu- re reprehendi possit. Excipimus enim hic iniurias et calumnias, quibus impii uel optime meritos de pietate et bonis omnibus nunquam non grauarunt. Apostoli Pauli studium reprehendebatur ab impiis omnibus, inte rim tamen ipse irreprehensibilis erat apud omnes san ctos, qui rectius iudicabant de uirtutibus et uitiis. Ex cipimus praeterea uitia atque peccata. Nam nemo alias uiuit qui non aberret et peccet. Christum solum ne- mo peccati arguere poterat. De sceleribus et impu- dentiore conuersatione intelligo hunc Pauli locum, ut uidelicet in episcopo exigat summum innocentiae studium, ne quod recte et sancte docet, praua et im- pura uita destruat. Quis enim patientibus auribus au diat impurum aliquem de sanctimonia interna et ex terna sancte licet disputantem? Proinde apostolus im  \pend
\vspace{0.5cm}\noindent
\vspace{0.2cm}\rule{1cm}{0.2pt}\\ 
\hspace{0.2cm}\textit{mg}
\footnotesize Irreprehe sibilis. 
\normalsize| 
\section*{AD TIM. CAP. III. }
\marginpar{[ p.124 ]}\pstart primis et in genere flagitat ut qui caeteris destinati sunt censores morum et magistri innocentiae, accusato res scelerum et omnis iniquitatis extirpatores, caueant sibi quam diligentissime ne in ipsis haereat tenatius quod aliis uolunt ademptum. Iamuero per partes et ueluti species subiectas exponit quomodo fiant irre- prehensibiles episcopi, si uidelicet sint unius uxoris mariti. Quo membro commendatur pudicitia et ca- stitas coniugalis, damnatur impudicitia, scortatio, adul terium et impuritas carnis. Quibus si episcopi pollu ti ac contaminati sint, sordebit plebi uniuersa ipsorum institutio, aperieturque tum fornicationi tum adulterio porta patentissima. Nam qua fronte qua authoritate quo denique affectu scortationem damnabit aut adulte ros accusabit qui omnium maximus et scortator et adulter est? Caeterum si quis pudice uiuit coelibemque ́ uitam agit non cogitur ulla dei lege connubio se irre- tire. Bonum enim est mulierem non tangere praecipue ob incommoditates uarias quae uerbi ministrum perse quutionis procella ingruente quatiunt atque obturbant. Verum si ea gratia homini data non estut sine uxo- re commode possit uiuere, melius est contrahere ma- trimonium quam uri. Neque uero sacrosanctum ministeri um polluit coniugium sanctum. Nam apostolus, Honorabile est, inquit, comnubium apud omnes et cubile impollutum. Son itaque sanctum domini institutum polluit sanctum dei  \pend
\vspace{0.5cm}\noindent
\vspace{0.2cm}\rule{1cm}{0.2pt}\\ 
\hspace{0.2cm}\textit{mg}
\footnotesize Vnius uxo ris maritus 
\normalsize| 
\section*{COMMENT. IN I. EPIST. }\pstart ministerium. Quo circa mirari non possum qua fidutia exposuerint quidam hunc Pauli locum non de coniugio epi- scoporum, sed de ecclesia potius et sacerdotio, quod ui- delicet Paulus nolit episcopo uni plures committi eccle sias. Quorsum enim (ut multa alia transeam) referent quod sequitur, Qui suae domui bene praesit, ꝗ liberos habeat subiectos et per omnia reuerentes? Si uero et haec exponant de membris ecclesiae, hoc est fidelibus, reuincentur per ea quae sequuntur protinus, Si quis propriae domui praeesse non nouit, quomodo ecclesi- am dei curabit? Hic enim nota manifestissima ostendit se loqui de familia et uxore non spirituali sed carna- li sanctaque ́. Iam quod obiiciunt apostolos quidem ha- buisse uxores, sed non circumduxisse secum in ministe rio, confutatur ex hoc loco, qui palam quid fieri de- beat praescribit, non quid gestum sit exponit. Episco- pus uerbi dei minister est. At hunc oportet unius uxo ris maritum esse, eo scilicet quia episcopus est, san- ctumque connubium aliis commendare et a fornicatio ne absterrere debet. Deinde clarissimis uerbis ipse a- postolus in epistola ad Corinth. 9.cap. An non habe- mus, ait, potestatem sororem mulierem circumducen di quemadmodum et caeteri apostoli et fratres domi ni et Cephas? Hunc in modum conuicti comminiscun tur aliud, quibus contendere quam ueritati obsequi iu cundius est, primum enim matrimonium admittunt,  \pend
\section*{AD TIM. CAP. III. }
\marginpar{[ p.125 ]}\pstart repetitum damnant uel improbant: hinc a Paulo di- serte Vnius uxoris maritum, unius inquam dictum putant. Quasi uero secundae nuptiae primis sint pollu tiores, et Paulus uerum non dixerit cum in Corinthiis inconiugatis et uiduis praescripsit bonum esse si sic manserint, caeterum si se non contineant contrahant matrimonium. Vnius ergo dixit ut excluderet uxorum gregem et concubinarum turbam, sanctum uero et iustum matrimonium stabiliret, quod non inter plures sed duos tantum confit. Id quod et dominus apud Mat thaeum docuit in cap.29. Verum plura de sacerdotum connubio annotauimus alibi. Sane planum est non mo do ex monumentis ueterum sed neotericorum quoquam omnes Galliarum et Germaniae episcopos sero tan- dem professionem coelibatus recepisse fecisseque ́. Ab an tiquo enim uxores habuisse longe certissimum est. Id quod et Lambertus Schaffnaburgensis historicus in rebus Germanicis luculenter exponit eo praecipue in loco ubi gesta praeclara anni 1o74. describit. Mira tum erat in episcopis eius temporis constantia. Confiden- ter enim et negligebant et damnabant sanctiones et decreta Romani praesulis uxoresque secundum uer- bum domini retinebant. vnφάλιoμ uigilantem retu- lere alii ad animum. Theophylactus enim, Sobrium ait esse oportet episcopum, hoc est perspicacem et meditabundum nullis huius saeculi curis uel moerore  \pend
\vspace{0.5cm}\noindent
\vspace{0.2cm}\rule{1cm}{0.2pt}\\ 
\hspace{0.2cm}\textit{mg}
\footnotesize Vigilantem 
\normalsize| 
\section*{COMMENT. IN I. EPIST. }\pstart offusum. Ego uero de tota pastorali cura expono. Spe culator enim et uigil est episcopus qui quoquo uer- sus oculos uigiles circumagendo hostes tum domesti- cos tum exteros obseruare cauereque debet ne irruptione facta periclitetur gregis dominicae salus et pietas. Huc pertinet Ezech.ca.3. et 33.item Acto.2ο. σώ͂φρονα, sobrium. Id uero ad animum refero episcopi, cuius pro prium est sapere et matura prudentia rebus consule re quibuslibet. Et dominus in euangelio prudentiam et fidelitatem requirit a dispensatore mysteriorum dei. Mens uero praua et prauis imbuta opinionibus recte putatur esse ebria, cui opponitur sobrietas. κὁσ μῖομ modestum siue ornatum hoc est compositum et bene moratum, id quod ad totum corporis habitum et cultum refertur. Plurimum enim momenti ista in uita hominis habent. Vestis igitur ne sit inculta, rursus ne sit sumptu osa laboriosaue. Nam ut hoc luxum habet, sic illud nescio quid hypocriticum resipit. Quidam existimant se pietatem Christianam non posse nisi militari ueste uel rufticana et prorsus sordida exprimere, horum studium unice huc spectat ut ueste maxime prophana utantur. Sunt alii qui in nouo quoque priuato et maxime hypocritico uesti- endi genere bonam pietatis partem existimant sitam, his sordent uulgares sed honestae uestes. At prudens episco pus honestissimos quosque in plebe uestium genere forma et usu imitatur, ita hanc rem temperans ut nihil in ipso  \pend
\vspace{0.5cm}\noindent
\vspace{0.2cm}\rule{1cm}{0.2pt}\\ 
\hspace{0.2cm}\textit{mg}
\footnotesize Sobrium. 
\normalsize| 
\hspace{0.2cm}\textit{mg}
\footnotesize Modestun 
\normalsize| 
\section*{AD TIM. CAP. III. }
\marginpar{[ p.126 ]}\pstart haereat incultum nihil hypocriticum, nihil item nouum prophanum aut luxuriosum. Hic enim eius quod ubique  locum habet meminisse oportet modum esse in re qualibet seruandum. Et episcopi praecipua functio est luxum et hy- pocrisim damnare omnem. Frugalitas enim et parsimo- nia uectigal omnium maximum est: sumptus uero immo- dicus uel uestibus uel uictui impensus ecclesiam et pol luit et malis artibus opplet. Curabit ergo episcopus cum doctrina tum exemplout totus populus honestis frugali bus atque decoris utatur uestibus. Mores episcopi sint candidi, graues, puri, alieni ut a leuitate uanitateque mum- dana, ita a stoicismo et tetrica pharisaeorum specie. Bre uiter in omni uitae consuetudine in quam sermones et colloqa maxime putamus pertinere, prae se feret epi- scopus grauitatem iustam et suauem, ut sit alienus ab omni leuitate morumque ineptia quae minuunt doctoris et authoritatem et reuerentiam. Proinde ut tetricos quosdam et misanthropos pastores probare non pos sumus, ita episcopos faciles omniumque horarum bel- los homunculos, leues atque nimium ciuiles facetosque ́ non damnare non possumus. Modestiam approbamus. σιλόξεvoμ hospitalem, et qui amet peregrinos, li- beralis et humanus sit. De hospitalitate disputauit Ci- cero quoque in Offi. In religione nostra inter primas nu meratur uirtutes. Habet et suam promissionem, habet exem pla uaria in Abraham praecipue, Loth, in Sunamitide,  \pend
\vspace{0.5cm}\noindent
\vspace{0.2cm}\rule{1cm}{0.2pt}\\ 
\hspace{0.2cm}\textit{mg}
\footnotesize Hospitalem 
\normalsize| 
\section*{COMMENT. IN I. EPIST. }\pstart in Martha, in Philippo. Acto.21.et Publio Acto.2s. Maiores nostri episcopis amplissimas opes in hoc donarunt ut hospitalitatem exercerent liberaliter, cae terum iniuria temporum et auaritia impiorum epi- scoporum factum est, ut non modo nulla amplius exercea tur ab illis hospitalitas, uerum etiam illis opibus nunc non ut episcopi et dispensatores praesideant sed ut principes imperitent iisque ad luxum et tyrannidem pro animi uoto in detrimentum ecclesiae maximum abutantur. διδ ακτικὸμ, id est docibilem, id quod a- ctiue intelligendum est, nempe qui docere possit. Haec prima et potissima est episcopi functio. Qualis uero esse debeat doctrina et alibi et in primo huius epi- stolae cap. ostendimus, hic de facultate et modo docen di loquitur. Inuenias qui non nesciant quid sit docen- dum, sed docere interim nesciant. Oportet ergo epi- scopum non tam doctum esse quam docibilem, hoc est appositum ad docendum, qui uidelicet sciat quid qui bus quatenus et quomodo docenda sint quaelibet. E- ant ergo nunc amusi bardi stupidique aphronii et dam nent dicendi artem sine qua nemo apposite docebit, ni si cui natura et deus ipse sarsit quod arte diminutum est. Contemptis ergo illiteratorum nugis et quere- lis qui consultum uolunt ecclesiis fideliter instituant iuuentutem ecclesiasticam in oratoriis praeceptis hoc est in recta dicendi arte. Agnoscimus enim et hanc  \pend
\vspace{0.5cm}\noindent
\vspace{0.2cm}\rule{1cm}{0.2pt}\\ 
\hspace{0.2cm}\textit{mg}
\footnotesize Aptum ad docendum. 
\normalsize| 
\section*{AD TIM. CAP. III. }
\marginpar{[ p.127 ]}\pstart non sine munere et afflatu dei ad nos esse delatam, ut prorsus nullum apud nos locum habeant clamores isto- rum de philosophia et gentilium literis contemnen- dis. Norunt sancti quatenus istis utantur in sacris. Aù rάροιvop, id quod recte uertit Ambrosius non obno xium uino, Non enim uini usum omnino damnauit, sed abusum ingurgitationem et uini studium assiduum, quo qui exercentur uino mancipati nihil aliud quam desipere possunt. Leuitis praecipiebat dominus absti nerent ab omni potu inebriante cum ingressuri essent in sanctum. Voluit autem illos esse sana mente et in- tegris sensibus. Vinum enim dementat et peruertit sensum adeo ut ne membra quidem corporis suum sa- ciant officium. Cauenda tamen imprimis spiritualis e- brietas exaecatio, inuidia, fastus, cupiditas nimia rerum pessimarum. Et Theophylactus, Non dicit hoc loco, inquit, qui mero madeat, sed contumeliosum et teme rarium. Ego uero Paulum omnino de temulentia et uini abusu loquutum existimo. μηὶ πλήκτημ non per cussorem, id quod non pertinet ad uiolentiam manuum, sed ad acerbitatem linguae, ne saeuus et improbus sit obiurgator. Est enim et in obiurgando modus, qui seruatur si obiurgatio sic sit temperata ut reus intelli- gat peccati sui magnitudinem, et quod paterna ob- iurgatio ad saluandum non conuiciandum tantum in stituta sit. Quidam plus nimio feroces licentius agunt  \pend
\vspace{0.5cm}\noindent
\vspace{0.2cm}\rule{1cm}{0.2pt}\\ 
\hspace{0.2cm}\textit{mg}
\footnotesize Non uino- lentum. 
\normalsize| 
\hspace{0.2cm}\textit{mg}
\footnotesize Non percus sorem, 
\normalsize| 
\section*{COMMENT. IN I. EPIST. }\pstart omnia, nimiumque frena conuiciatricis linguae laxan- tes reorum animos tam a se quam a uero prorsus alie nant et repellunt, idem illi alios quoque pusilanimita- tis atque perfidiae damnant qui non pari improbitate et seuitia utuntur in obiurgando. Caeterum Paulus iubet e- piscopum μηι πλήκτηpesse. Nam uult eum abhorrere a contentione et saeuitia ne uideatur odio non amore facere quod facit. Interim tamen ad scelera facineroso rum conniuere istum nolumus, ita nec crimina extenu are sed intrepide et sancta cum moderatione argue- re et obiurgare iubemus. μιι αιαρομερδιη non turpis lucri cupidum. Tametsi enim hospitalitas ab episcopo requiratur attamen huius praetextu nullum sectabitur lu crum inhonestum: quod faciunt qui cum pietatis et religionis detrimento munera accipiunt, deinde rebus turpibus in honestis illiberalibus et episcopo indecoris quaerunt ui ctum. Non damnatur omne lucrum sed indecorum modo et in iustum. Plura annotaui de hac re in s.ca.1. epi. Pet.et12 ἐπισκῆ sed aequum id est moderatum pbum humanum man suetum et frugi hominem. Hoc enim superioribus uitiis op£ ponit. αμα αχoν alienum a pugna id est absque iurgio et litigandi contendendiue ardore. Recte enim dictum est ueritatem nimia contentione amitti. ἀφι λάργυρομ abar genti cupiditate alienum, qui uidelicet studium omne suum intenderit corradendis pecuniis et congerendis the- sauris. Auaritia enim radix est omnium malorum. Nec  \pend
\vspace{0.5cm}\noindent
\vspace{0.2cm}\rule{1cm}{0.2pt}\\ 
\hspace{0.2cm}\textit{mg}
\footnotesize Non turpis lucri cupi- dum. 
\normalsize| 
\hspace{0.2cm}\textit{mg}
\footnotesize Sed aequum 
\normalsize| 
\hspace{0.2cm}\textit{mg}
\footnotesize Alienuma pugna. 
\normalsize| 
\hspace{0.2cm}\textit{mg}
\footnotesize Alienum ab auaritia. 
\normalsize| 
\section*{AD TIM. CAP. III. }
\marginpar{[ p.128 ]}\pstart quisquam Deo et mamonae inseruire potest. Debet igi tur episcopus ab auaritia, muneribus, a quaestu turpi, a lucro inhonesto, et cumulandi industria ac diligen- tia quam alienissimus esse. Hodie industrii et egregii creduntur episcopi qui auri ac argenti cumulos ma- gnaque pondera posteris ad bella contra Christum instau randa et alenda relinquunt. Tam probe episcopi et proceres ecclesiastici hodie ex sacris instituti suum fa ciunt officium scilicet.  \pend
\phantomsection
\addcontentsline{toc}{subsection}{\textit{qui suae domui bene praesit, qui liberos habeat in subiectione cum omni reuerentia. Quod si quis pro- priae domui praeesse non nouit, quomodo ecclesiam dei curabit? }}
\subsection*{\textit{qui suae domui bene praesit, qui liberos habeat in subiectione cum omni reuerentia. Quod si quis pro- priae domui praeesse non nouit, quomodo ecclesiam dei curabit? }}\pstart Quemadmodum uero auaritia uitio datur episco- pis ita et profusio et rei familiaris neglectus. Auaro non tam deest quod habet quam quod non habet. Pro digo uero non tam prodest quod habet quam quod non habet, utpote cui nihil satis est quique ; abliguriret omnes omnium facultates si quis modo illos suppe- ditaret. Vterque in extremis quod dici solet haeret, u- terque culpabilis est, qui uero medium tenet is iusta am bulat. Quocirca Paulus designat episcopum qui aua- rus non sit et tamen suae domui propriae (τοϋἴδlιoυ ὁικον καλως προίςάμενον) bene praesit, hoc est,  \pend
\vspace{0.5cm}\noindent
\vspace{0.2cm}\rule{1cm}{0.2pt}\\ 
\hspace{0.2cm}\textit{mg}
\footnotesize Frugi ho- mo sit epi- copus. 
\normalsize| 
\section*{COMMENT. IN I. EPIST. }\pstart qui opes et facultates omnes recte et cum iudicio di spenset. Nulla enim familia nulla domus sine opum et facultatum usu et dispensatione sustentari administrari ue potest. Id quod ppter eos moneo qbus subinde in o- re est uerbi ministros leberide nudiores esse debere, atque ideo hunc Pauli locum non de opum iusta et frugali di spensatione, sed de liberis exponunt, ut Paulus per ipsam domum intellexerit liberos, quod quidem ut non infrequens est ita nemo negare potest ad liberorum educationem alimoniam et curam facultates quoque necessarias esse. Pertinent igitur ad rem domesticam curandam pecu- nia, commutatio, locatio, contractus, emptio, uenditio, et si quae alia sunt rei familiaris. In his requirit apo- stolus hominem frugi et sanctum, non prodigum et qui rerum suarum nullum habeat pensum. Et expe- rientia magistra didicimus huiusmodi decoctores mi- nisterio sancto non semel maculas ignominiae altissimas inussisse. Idcirco Paulus frugi hominem designat. Vnde consequitur quod minime anxia diffidens et peruersa, sed syncera sancta iustaque rei familiaris cura episcopum ornat, neglectus uel auaritia defoedat. Deinde requiri tur ut episcopus liberos habeat obedientes et reueren tes, qui pudore ingenuo moribusque compositis pietate singulari educationem patris sanctam comprobent a- liis, iisque exemplum uitae praeferant melioris. Et liberi quidem parentum uiua sunt simulachra, in quibus ple-  \pend
\section*{AD TIM. CAP. III. }
\marginpar{[ p.129 ]}\pstart rumque parentum genius, disciplina, probitat et dili- gentia relucet. Haec uero tanti sunt apud apostolum momenti, ut ex his ceu prognosticis indiciisque aestimari ac censeri uelit episcopum, subiiciens comparatum quoddam, Quod si quis propriae domui praeesse non nouit, quomodo ecclesiam dei curabit? Opponit enim priuatam domum et dispensationem episcopi deligen di, ecclesiae et ministerio siue gubernationi publicae. Qui, inquit, in curanda re familiari et liberorum edu eatione egregiam nauarit operam, is specimen de se praebuit foreut publicam omnium curam gerere gna uiter possit. Caeterum qui iners est in administrations rerum exiguarum is nulla cum laude maiora obibit munia. Quia uero summa quaeque ministeria uanissi- mis inertibus et helluonibus committuntur hinc sini stre geruntur omnia.  \pend
\phantomsection
\addcontentsline{toc}{subsection}{\textit{Non nouitium, ne inflatus in condemnationem incidat calum- niatoris. Oportet autem illum et bonum habere testimonium ab ex- traneis ne in probrum incidat et laqueum calumniatoris. }}
\subsection*{\textit{Non nouitium, ne inflatus in condemnationem incidat calum- niatoris. Oportet autem illum et bonum habere testimonium ab ex- traneis ne in probrum incidat et laqueum calumniatoris. }}\pstart veóφurog iuxta etymon est nuper natus aut nu- per insitus, diciturque ́; non tam de iis qui nuper in fidei contuberium asciti Christi nomen profiteri coeperunt  \pend
\vspace{0.5cm}\noindent
\vspace{0.2cm}\rule{1cm}{0.2pt}\\ 
\hspace{0.2cm}\textit{mg}
\footnotesize Lucae.16. 
\normalsize| 
\hspace{0.2cm}\textit{mg}
\footnotesize Non nouï- tium. 
\normalsize| 
\section*{COMMENT. IN I. EPIST. }\pstart quam de iis qui docti quidem sunt et periti diuinae le gis, sed non solide. Solet autem eruditio immatura et inflare et confidentes facere. Addit igitur Paulus cau sam cur nolit nouitium, Ne inflatus et plus nimio sibi placens incidat in maledicorum condemnationem, hoc est, ne illi qui religionem nostram unice ut uelli- cent obseruant, iustam a nouitio conuitiandi ansam ac cipiant, meritoque et fastum et inconstantiam eius ac cusent. Nam solent literatores et neophyti procatius erumpere et parum expensa temere proponere, quae demum retractare et nonnunquam plane coguntur recantare, id quod authoritatem eorum imminuit haud exigue sed et ministerium ipsum nonnihil infame red dit, quocirca apostolus reiicit neophytos, et uel erudi tione iusta insignes, uel diutino usu exercitos ecclesiae iubet praeficere. Addit et aliud, Oportet autem illum et bonum habere testimonium ab extraneis. Per ex- traneos uero intellexit alienos a Christiana religione ethnicos uidelicet. Tanta inquam innocentia conspi- cuum uult esse episcopum ut ne gentes quidem habe- ant quod iure obiiciant flagitium. Et Aelius Lampridi us in uita Alexandri Seueri scribit Chriftianos si quan- do sacerdotes ordinare uoluissent prius ipsorum pro posuisse nomina palam, hortatosque populumut si quis quid haberet criminis probaret manifestis rebus. Qui mos ex hoc nimirum praecepto Pauli inoleuit et ad id  \pend
\section*{AD TIM. CAP. III. }
\marginpar{[ p.130 ]}\pstart utilis fuit quod ethnici de pietate Christianorum non male sentientes ab impietate sua facilius desciuerunt. Hodie uera religio ob improbitatem quorundam mi- nistrorum non bene audit. Et haec quidem de ministro rum uita et moribus dilucide et breuiter disputauit apostolus. Ad hanc tractationem pertinent quae se- quuntur.  \pend
\phantomsection
\addcontentsline{toc}{subsection}{\textit{Ministros itidem compositos, non bilingues, non multo uino de- ditos, non turpiter lucri auidos, te- nentes mysterium fidei cum pura conscientia. Atque hi probentur prius, deinde ministrent sic ut ne- mo possit illos criminari. }}
\subsection*{\textit{Ministros itidem compositos, non bilingues, non multo uino de- ditos, non turpiter lucri auidos, te- nentes mysterium fidei cum pura conscientia. Atque hi probentur prius, deinde ministrent sic ut ne- mo possit illos criminari. }}\pstart Diaconi, quod et Ambrosius sentit, sunt episcopo rum ministri et ueluti membra adeoque et seges epi- scoporum. Hi sane duplicis generis et functionis inue niuntur in sacris. Alii enim aerario praefecti ecclesiasti co pauperum curam gerunt, ut uidere est in Actis ca. σ. de quibus et hic locus per omnia exponi potest. Alii uero literis honestis disciplinis et scripturae sa- erae inuigilant adeoque in rebus exercentur sacris assi- due ut aliquando muniis praefecti ecclesiasticis eccle- siae dei quamplurimum prosint. Ex his enim deligun-  \pend
\vspace{0.5cm}\noindent
\vspace{0.2cm}\rule{1cm}{0.2pt}\\ 
\hspace{0.2cm}\textit{mg}
\footnotesize Diaconi li- teris manci pati iuue- nes. 
\normalsize| 
\section*{COMMENT. IN I. EPIST. }\pstart tur episcopi siue sacerdotes ecclesiarum ministri. Re- ceptum fuit et prophetarum temporibus adulescen- tes ad ministeria alere sacra. Nota est historia de Sa- muele. Traduntur item fuisse scholae cum in Iericho cui praeerat Heliseus, tum in Ramoth Galaad. Harum alumni et Nazarei et prophetarum filii appella- bantur. Occurrunt ubique in prophetis istarum rerum egregia specimina. Et sane religioni et pietatis stu- dio sine disciplina et scholis huiusmodi non bene con sulitur. Coeperunt ergo ab ipsis protinus incunabulis Christianae fidei, magistris et authoribus apostolis, ali ingenia egregia educarique homines de quibus spes e- rat fore ut aliquando doctrina et uita ecclesiam a ma ioribus plantatam iusta cura et agricolatione proue herent ad maturam frugem. Talibus ergo nunc prae- scribit leges, uidelicet ut imprimis sint compositi, quod ad morum candorem pertinet, ad pudorem liberalem et ingenuum, ad grauitatem et reuerentiam, ut scili cet toto uitae habitu disciplinam quandam prae se fe- rant Christianam. Deinde praecipit ne sint bilingues, hoc est inconstantes, uani, futiles, leues et mendaces, adulatores et palpones. Notati sunt hi et a ueteribus elegantissimis prouerbiis, cuius generis sunt, Aliud stans aliud sedens loquitur, Duos parietes de eadem dealbat fidelia, Ex eodem ore calidum et frigidum es flat. Foedum sane hoc uitium est cum in omnibus tum  \pend
\vspace{0.5cm}\noindent
\vspace{0.2cm}\rule{1cm}{0.2pt}\\ 
\hspace{0.2cm}\textit{mg}
\footnotesize Non bilin. gues. 
\normalsize| 
\section*{AD TIM. CAP. III. }
\marginpar{[ p.131 ]}\pstart in ministris uerbi praecipue, aut iis quoquam qui uerbi mi nisterio destinantur. Horum enim ora et linguae spi- ritussancti hoc est constantiae ueritatis iustitiae et pu- ritatis organa sunt. Quod uero paulo ante breuius et ideo obscurius posuit μηι τάροινoμ non uinolentum, id nunc pluribus et explicatius exponit. μηι ὁινφ τολλῶ τροσέχουτας non multo uino addictos. Non enim uini usu interdicit in uniuersum sed abusu et stu dio nimio cui raro abest ebrietas. At hanc celeberrimi e nostris authoribus uoluntariam esse insaniam dixere, quae si in plures paulo extendatur dies habitus is etiam manifestarius euadat furor. Celebre est et illud ignem gestare in sinu, qui assiduo utatur uiuo. Et Solon unus e septem morte principem esse mulctandum lege san- xit si deprehenderetur ebrius. Pittacus quoque ebrios si peccassent duplici dignos incommodo statuit. Sed et dominus ipse magna cum indignatione contra Israeli tas Amosz, cap. Suscitaui, ait, de filiis uestris prophe tas, et de iuuenibus uestris Nazareos. An non ita se habet res? Vos autem fecistis Nazareos bibere uinum. Accusat autem Israëlitarum luxum et quod iuuentu tem ad prophetiam delectam educatione molli cor- rumperent. Omnium itaque episcoporum erit diligen ter curare quo adolescentia Christiana et sacro mi- nisterio dicata sancte et in literis puris et in mori- bus integris insituatur synceris, ut tempestiue et ab  \pend
\vspace{0.5cm}\noindent
\vspace{0.2cm}\rule{1cm}{0.2pt}\\ 
\hspace{0.2cm}\textit{mg}
\footnotesize Non multo uino dedi- tos. 
\normalsize| 
\section*{COMMENT. IN I. EPIST. }\pstart ipsis incunabulis bona amare et mala fugere ac auer sari assuescat. Sequitur enim quod huc pertinere ui- detur, Tenentes mysterium fidei cum pura conscientia. Neque enim satis est si sint compositi, si in omni erudi- tione et doctrina praeclari, nisi et fidei mysterium teneant, hoc est pii sint et totam religionis rationem probe teneant, adeoque et ex animo credant alieni ab omni fuco et hypocrisi. Nam in hunc finem conscien tiam puram diserte expressit. Inuenies enim studiosos quidam literarum sed parum pios, Inuenias et pie in sacris eruditos qui suauiter de religione philosophen tur, sed non ex conscientia pura uerum ex hypocri- si. Addit igitur et aliud, Atque hi probentur prius, de- inde uero ministrent atque ita ut nemo possit illos cri- minari. Quasi dicat, Cum uero hypocrisis maximum quidem sed occultum malum, non raro corruptissimos ad ministerium subuehat sacrum, hinc nolim protinus committi sacrarum rerum administrationem uel exi- miae spei studiosis nisi multo iam tempore exploratis. Ita enim fiet ut nemo possit eximios et probatos cri- minari temere. Haec si nobis essent sacrosancta pro- culdubio melius haberent res ecclesiasticae et maior uberiorque esset prouentus pietatis.  \pend
\phantomsection
\addcontentsline{toc}{subsection}{\textit{Vxores similiter modestas, non ca- lumniosas, sobrias, fidas in omnibus  }}
\subsection*{\textit{Vxores similiter modestas, non ca- lumniosas, sobrias, fidas in omnibus  }}
\section*{AD TIM. CAP. III. }
\marginpar{[ p.132 ]}\pstart Multum refert quam quisque nactus sit uxorem. Nam in uniuersa uiri uita haec momentum non leue habet. Periclitatur enim ex illius improbitate ipsa uiri exi- stimatio et gloria, eadem augescit ex sludiis et mori bus eius sanctis. Vnde Sapiens Hebraeus, Mulier sedu- la, ait, corona est mariti sui, sed quasi putredo in ossi- bus eius est mulier turpis. Mulier sapiens aedificat do- mum, stulta uero subuertit eandem. Item, Mulier bona pars bona, oblectat uirum suum et annos uitae illius in pace implebit, Gratia super gratiam mulier sancta et pudorata. Nobilis in portis uir eius quando sede- rit cum senatoribus terrae. Disciplina illius datum dei est. Sicut sol oriens in mundo in altissimis dei, sic et mulieris bonae species in ornamentum domus eius. Proinde apostolus iubet et episcoporum et diacono rum uxores esse σεμνας compositas modestas, non le ues aut ineptas, sed matronali pudore graues. Deinde αηι διαθόλους non calumniosas, quod plerumque ad- haerescit mulierculis. Dum enim curiosae per alienas perreptant aedes, dum negotiosae et garrulae alienis se immiscent causis, nascitur obtrectandi susurrandi ca- lumniandique negotium. Ab iis uitiis abesse uult uxo- res ministrorum ueritatis concordiae et pacis. Et silen tium praecipuum est omnium mulierum ornamentum. vuφ et Xioue id est uigilantes, quod tam ad sobrietatem animi temperantiamque quam ad diligentiam et pru-  \pend
\vspace{0.5cm}\noindent
\vspace{0.2cm}\rule{1cm}{0.2pt}\\ 
\hspace{0.2cm}\textit{mg}
\footnotesize Vxores mi nistrorum uerbi qua- les esse de- beant. 
\normalsize| 
\section*{COMMENT. IN I. EPIST. }\pstart dentiam, hoc est solertem rei familiaris administra- tionem referri potest, quemadmodum et sequens por tio τιςας ἐμ πᾶσι fidas in omnibus pertinet, et ad si dem coniugalem et ad fidelem facultatum dispensa- tionem, uidelicet ne sint adulterae, impurae, futiles, fura ces et inconstantes. Quemadmodum enim episcopi uita uitaeque totius habitus aliorum institutio, sic mulie rum omnium exemplar et pudicae probaeque uitae ty- pus esse debet purior uxoris episcopi conuersatio.  \pend
\phantomsection
\addcontentsline{toc}{subsection}{\textit{Diaconi sint unius uxoris ma- riti, qui liberis recte praesint et pro- priis familiis. Nam qui benc mini- strauerint gradum sibi bonum ac- quirunt et multam libertatem in fide quae est in Christo lesu. }}
\subsection*{\textit{Diaconi sint unius uxoris ma- riti, qui liberis recte praesint et pro- priis familiis. Nam qui benc mini- strauerint gradum sibi bonum ac- quirunt et multam libertatem in fide quae est in Christo lesu. }}\pstart Ne diaconorum periclitaretur pudicitia et uel a- dulterio uel scortatione polluerentur atque ita maculam sibi opprobriumque asciscerent perpetuum, permittit ipsis connubium. Nam melius est matrimonium contrahe re quam uri. Et quod hac in re praescripsit episcopis ut unius uxoris mariti sint, ut liberis et familiae bene praesint, idem praecipit et diaconis. Addit hoc loco spem et egregiam pollicitationem, qua admoueret cal car currentibus, hoc est excitaret diligentiam ultro-  \pend
\section*{AD TIM. CAP. III. }
\marginpar{[ p.133 ]}\pstart neam. Sensus enim est, Tametsi functio diaconi paulo inferior sit functione episcopi, id tamen neminem red- dat ignauum. Spectanda sunt proposita praemia. Quis quis enim in hoc ordine se gesserit gnauiter et inte- gre, non modo gradum sibi sternit ad munina excel- lentiora, uerum etiam auctionem accipiet a deo spiri tusssancti, ωoλιίμ inquam παρρησίαμ multam li- bertatem in fide, id est fidei religionis sapientiaeque coe lestis apertam liberalemque cognitionem. Duo itaque sunt quae sirenue suum officium facientibus pollice- tur apostolus: prius prope terrestre est, fore scilicet ut excellentiori praeficiatur officio: posterius plane spiri tuale et coeleste est, fore ut fidei et cognitionis auge atur donum. Recte enimuereque Paraphrastes, ut in re publica prophana inquit uocatur ad praeturam qui in aedilitate strenue se gesserit, rursus a praetura ad tribu natum aut consulatum qui illic quoq uicerit seipsum, ita diaconi functio declarat qui sit sacerdotis aut apo stoli loco dignus. Et in euangelio disertis uerbis de dono rum auctione loquens dominus id proposuit parado- xon, Omni habenti dabitur et abundabit, qui uero non habet etiam id quod habet auferetur ab eo. Mat- thaei2s. Neque uero nouum est quod egregia inges nia spebus et pollicitationibus ducuntur incitantur atque capiuntur. Debetur sane uirtuti suum praemium et honor iustus. Proinde praeter ius et phas spe sus  \pend
\section*{COMMENT. IN I. EPIST. }\pstart frustrantur aliquoties uiri aut adolescentes industrii, qui cum singulari diligentia et eruditione egregiam spem de se praebuere amplioribus praemiis ad maiora erant prouehendi, nunc uero negliguntur aliquoties prouehunturque ad honores qui non tam indocti quam indigni sunt, quo demum fit ut et illi animum despon deant et asini, quod dici solet, portent mysteria.  \pend
\phantomsection
\addcontentsline{toc}{subsection}{\textit{Haec tibi scribo sperans fore ut ueniam ad te cito, quod si tardius uenero ut noris quomodo opor- teat in domo dei uersari quae est ec- clesia dei uiuentis columna et ba- sis ueritatis. }}
\subsection*{\textit{Haec tibi scribo sperans fore ut ueniam ad te cito, quod si tardius uenero ut noris quomodo opor- teat in domo dei uersari quae est ec- clesia dei uiuentis columna et ba- sis ueritatis. }}\pstart Institutum et usum praecedentium ostendit prae- sentibus, ex quibus facile intelligimus Paulum hisce capitibus tribus totius ecclesiasticae oeconomiae sum- mam adumbrasse:ut uani et ingrati sint spirituisan- cto quicunque formulam oeconomlae petunt aliam. Iam ut diligentiam et studium Timothei excitaret ma- gnifice splendideque de coetu sanctorum loquitur, quem- admodum et in Actis cap.2o.dicens, Attendite uobis et cuncto gregi in quo uos spiritussanctus posuit e- piscopos ad regendum ecclesiam dei quam acquisiuit sanguine suo. Ita enim et in praesentiarum coetum il-  \pend
\section*{AD TIM. CAP. III. }
\marginpar{[ p.134 ]}\pstart Tum fidelium appellitat domum dei uiuentis, alludens quidem ad Hierosolymitanum templum, sed ipsam in- terim ueritatem praedicans et typis praeferens. Inha- bitat enim deus uiuitque non in templis saxeis sed in cor dibus fidelium, in quibus et uires suas exerit et to- tus regnat in ipsis. Nam qui sese Christo per fidem con secrarunt, non iam ipsi uiuunt, sed uiuit in ipsis Chri stus. Huiusmodi uero ecclesia est ςύλος κα) εδραίοͅ αα, της ἀληθείας hoc est columna et firmamentum siue basis ueritatis. In coetu enim sanctorum praedica- tur illud magnum pietatis mysterium, huic item inniti tur ecclesia, in hoc seruatur, in eodem perstat et per- seuerat ueluti ahaeneus murus contra omnes diaboli mundique insultationes. Vnde in euangelio dixit domi nus Petro confitenti, Tu es Christus filius ille dei ui- uentis, Tu es Petrus et super hanc petram aedificabo ecclesiam meam, et portae inferorum non praeuale- bunt aduersus eam. Abusi sunt hisce locis pontificii. Cum enim mendacia sua impias item et tyrannicas le ges auaritiam quoque et quaestum suum scripturis defende re non potuerunt, confugerunt ad ecclesiae patrocini um, eolligentes ecclesiam esse fundamentum uerita- tis, iccirco et leges ab ipsis latas esse ueritati consen- taneas, cum ecclesiae sint praesides adeoque ipsa eccle- sia. Caeterum fallatia in hoc est quod uariant in nomi ne eeclesiae. Paulus eam ecclesiam uocat ueritatis co-  \pend
\vspace{0.5cm}\noindent
\vspace{0.2cm}\rule{1cm}{0.2pt}\\ 
\hspace{0.2cm}\textit{mg}
\footnotesize Ecclesia do mus dei est. 
\normalsize| 
\hspace{0.2cm}\textit{mg}
\footnotesize Heclesia est columna ue ritatis. 
\normalsize| 
\section*{COMMENT. IN I. EPIST. }\pstart lumnam quae Christo nititur, solius pastoris uocem au dit, nec ipsa sibi uiuit, sed Christo qui et regnat totus in ipsa. Isti uero sibi adhuc uiuentes et suam gloriam et utilitatem quaerentes pastoris uocem non modo non agnoscunt, sed et legibus latis testantur se Christi ad uersarios esse, interim audent sanctae praetextu eccle- siae suas traditiones ineptissimas uenditare pro sacro- sanctis. Huc scilicet impudentiae uenit hominum auda cia. Ecclesia quae Christo petrae nititur, Christi spiritus alitur, Christo item capite regitur, Christi uerbo duci- tur, quae in hoc uiuit et ipse in ea, ipsa sane columna et basis ueritatis est. Veritati enim Christo nititur uincitque in eauirtus et gratia Christi. Hic enim toti corpori et membris singulis salutem dat, quemadmo dum expositum est in 4.et 5.cap. episto. ad Ephe. Porro qui hac exciderunt gratia, qui non ducuntur uerbo Christi, aut qui non ad hunc quem diximus mo- dum reguntur spiritu et fide Christi, hoc est in qui- bus neque Christus uiuit neque ipsi in eo, non est quod in negotiis suis ecclesiae quaerant patrocinium.  \pend
\phantomsection
\addcontentsline{toc}{subsection}{\textit{Et citra contrauersiam magnum est pietatis mysterium, Deus mani- festatus est in carne, iustificatus est in spiritu, uisus est angelis, praedi- }}
\subsection*{\textit{Et citra contrauersiam magnum est pietatis mysterium, Deus mani- festatus est in carne, iustificatus est in spiritu, uisus est angelis, praedi- }}
\section*{AD TIM. }
\marginpar{[ p.135 ]}
\phantomsection
\addcontentsline{toc}{subsection}{\textit{catus est gentibus, fides illi habi- ta est in mundo, receptus est in gloria. }}
\subsection*{\textit{catus est gentibus, fides illi habi- ta est in mundo, receptus est in gloria. }}\pstart Dixerat ecclesiam esse basim ueritatis, id uero nunc explicans quae nam illa sit ueritas cui innitatur eccle- sia et quam certa quantaeue sit authoritas exponit, idque exordiolo praemisso commendat dicens, uαι ὀμο λογουμένως id est et confesse, et citra conntra- uersiam, siue ut Ambrosius uertit, Et quidem omnium confessione magnum est illud pietatis Christianae my- sterium, summarium uidelicet religionis nostrae. Hic addit quod sensum absoluat, Vt ecclesiam hoc tenentem mysterium non temere uocauerim columnam et funda mentum ueritatis. Iam sequitur breuis huius mysterii expositio, Deus manifestatus est in carne, quod quid a liud est quam quod loannes quoque scripsit, Et uerbum ca ro factum est? Pertinent huc quaecunque in sacris tradun tur de diuinitate Christi et incarnatione eius, deque  causis et efficatia eius. In hoc filius nostram induit car nem peccati, ut per peccatum damnaret peccatum cre- dentesque efficeret aeternae uitae haeredes, haec uidelicet ecelesiae sanctae basis est, et haec prima et certa ueri tas est quod deus in Christo manifestauit se homini- bus, hoc est per incarnationem euidenter apperuit to ti mundo quali sit in genus humanum praeditus ani-  \pend
\vspace{0.5cm}\noindent
\vspace{0.2cm}\rule{1cm}{0.2pt}\\ 
\hspace{0.2cm}\textit{mg}
\footnotesize Summa ue ritatis. pri- mae, 
\normalsize| 
\section*{COMMENT. IN I. EPIST. }\pstart mo, optimo clementissimo saluifico. Proinde relucel in Christo ipsissima dei imago, adeo ut patrem uiderit ipsum qui uidit filium, Ioan.14. Sequitur, Iusti ficatus in spiritu, quod loquutionis genus idioma resipit He- braicum. Hoc enim significat, Per spiritum declaratus est deus et is esse qui a patre nobis datus sit, iustitia, sanctificatio, redemptio et salus uera, iustificans omnes qui p eundem spiritum ipsi fidunt soli relictis omnibus spebus mediis et elementis uisibilibus externisue. Huc iam pertinent sequentia scripturae testimonia, Ro.1. Qui declaratus fuit filius dei per potentiam et spiritum san ctificationis. Item loan.1. Vidi spiritum, inquit Bapti- sta, descendentem specie columbae de coelo et mansit super eum, et ego non noueram eum, sed qui misit me ut baptizarem aqua is dixit mihi, Super quem uideris spiritum descendentem ac manentem super eum, hie est qui baptizat spiritu sancto. Et ego uidi et testatus sum hunc esse filium dei. Huc pertinet postrema pars 3. capitis ad Rom. Praeterea testatur spiritus omnium sanctorum Iesum esse Christum et filium dei, agnum inquam illum dei qui tollit peccatum mundi. Innuit item hic Pault locus uerum dei cultum ueram iustitiam non in rebus externis et carnalibus, sed in spiritu et fide consistere. In his ergo duobus membris quod deus manifestatus est in carne, quod iuftificatus in spiritu, contmnetur totum pietatis Christianae mysterium, in iisdem  \pend
\section*{AD TIM. CAP. III. }
\marginpar{[ p.136 ]}\pstart von minus atque in sequentibus explicatur quam ma- gnum hoc est celebre et augustum quam denique cer- tum et indubitatum sit istud mysterium, quod uide- licet ab ipso patre aeterno per filium unigenitum per agitur, quodque omni modo spiritussancti testimonio comprobatur. Quo illud quoque referendum est quod uisus est angelis dominus noster. Nam et ipsa Christi domini incarnatio, natiuitas, tentatio, passio, resurre- ctio et ascensio ad coelos non sine angelorum ministe rio, hoc est praesentia et testimonio peracta est. Vbiq enim creatori suo inseruiunt spiritus coelestes, ubique ap£ probant nobis mysteria sacra, iustificationem et sa- lutem per Christum partam. Huc pertinet apostolorum testimonium et quod in hoc totus mundus conspira- uit. Hoc enim Paulus significauit dicens, Praedicatus est gentibus, fides illi habita est in mundo. Postremo uero receptus est in gloria. Deuictis enim triumpha- tisque daemone inferis et peccatis considet nunc ad de- xteram patris, regnatque in orbe toto. Quae quidem o- mnia, aeternum uidelicet consilium, potentia dei, testi monium spiritus, angelorum item et sanctorum omni um apostolorum praecipue praedicatio, denique et con sensus orbis publicus et uis Christi regnantis e coelo potentissima, declarant quantum sit pietatis mysteri- um, quam certum et indubitatum, quam denique co- lumnae firmae et fundamento ueritatis solido innita-  \pend
\vspace{0.5cm}\noindent
\vspace{0.2cm}\rule{1cm}{0.2pt}\\ 
\hspace{0.2cm}\textit{mg}
\footnotesize Celebritas et certitu- do euangeli- cae praedica tionis. 
\normalsize| 
\hspace{0.2cm}\textit{mg}
\footnotesize Lucae.1.2 Matt.4.26 Lucae.24. Acto.1. 
\normalsize| 
\section*{COMMENT. IN I. EPIST. }\pstart tur sancta dei ecclesia, ut uideas nunc ad expositionem superiorum haec esse adiecta. Cedant his quaecunque di cuntur ab hominibus ad authoritatem parandam nu- per inuentis religionibus, ecclesiastica pietas uera de mum et infallibilis pietas est. Huius generis et alia quae dam leguntur apud Augustinum contra Manicheum de utilitate credendi ad Honoratum cap.17. Haec uero a Paulo prolata sunt hoc consilio, ut Timothei industri am et excitaret et animum eius confirmaret. Et sane non potest animus non confirmari in fide cum audit tantum esse pietatis mysterium, nemo non cupit inse- ri ecclesiae quam audit columnam et basim ueritatis esse.  \pend
\phantomsection
\addcontentsline{toc}{subsection}{\textit{DE DOCTRINIS DAEMONIORVM. }}
\subsection*{\textit{DE DOCTRINIS DAEMONIORVM. }}\pstart Episcopus uerus non tantum quid uerum sit sed quid falsum quoque in re qualibet scire debet. Hinc ue- ro Paulus subiungit quaedam de doctrinis peruersis, hae ipsa tractatione ad uigiliam et diligentem industri- am excitans episcopos. Nam qui periculum sibi certo imminere norunt accuratius aduigilant.  \pend
\endnumbering\beginnumbering\section{CAP.IIII.}
\phantomsection
\addcontentsline{toc}{subsection}{\textit{Spiritus autem certo loquitur quod in posterioribus tempori- bus desciscent quidam a fide, atten- dentes spiritibus impostoribus ac doctrinis daemoniorum, per simu- }}
\subsection*{\textit{Spiritus autem certo loquitur quod in posterioribus tempori- bus desciscent quidam a fide, atten- dentes spiritibus impostoribus ac doctrinis daemoniorum, per simu- }}
\phantomsection
\addcontentsline{toc}{subsection}{\textit{}}
\subsection*{\textit{}}
\section*{AD TIM. CAP. IIII. }
\marginpar{[ p.137 ]}
\phantomsection
\addcontentsline{toc}{subsection}{\textit{lationem falsiloquorum cauterio notatam habentium conscientiam, }}
\subsection*{\textit{lationem falsiloquorum cauterio notatam habentium conscientiam, }}\pstart In genere dicit uenturos prophetas falsos. Quo plus autem fidei haberet prophetia, Non meum hic, inquit, somnium referam sed quod spiritus ille dei cu- ius oculis patent omnia pκτῶς diserte praescripteque ́ suggerit, nempe quod temporibus sequentibus aposto lorum uidelicet saeculum et ad iudicii usque diem pro- currentibus uenturi sint pseudoprophetae. Atque hos iam suis coloribus distinguit. Primo αποςήσουται, id est deficient a fide, a uera inquam religione et do ctrina apostolorum, non seruantes puritatem et sim- plicitatem pietatis uerae, sed attendentes spiritibus im postoribus, hoc est doctrinis daemoniorum. Monuimus enim in2. cap.2. ad Thess. epistolae omnino dignos esse ut spiritu seducantur diabolico qui spiritum Christi do ctorem audire noluerunt. Hi deinde malo spiritu im- buti doctores in hypocrisi loquentur docebuntque men dacia hoc est a scripturis aliena, quae uidelicet inspira uit impostor sathanas. Illi denique ipsi, quod tertio sub iungitur loco, apud simplices se uenditabunt persona sanctimoniae, cum apud deum habeant conscientiam insynceram et compunctam. Id quod metaphora cau terii a corpore ad animum deflexa adumbrauit. Cau- terium enim authore Coelio lec, antiq lib.5. cap.31.  \pend
\vspace{0.5cm}\noindent
\vspace{0.2cm}\rule{1cm}{0.2pt}\\ 
\hspace{0.2cm}\textit{mg}
\footnotesize Descriptio pseudopro- phetarum. 
\normalsize| 
\section*{COMMENT. IN I. EPIST. }\pstart instrumentum est quo putridae carnes aduruntur et ulcera. Inde (addit idem) cauteriata conscientia apo- stolo Paulo ueluti cauterio siue cautere ambusta dici- tur. In quam sententiam propemodum et S. Ambron- sius hunc locum exponens, Illorum, inquit, cauteria tam dicit conscientiam hoc est simulatione corruptam, quia sicut cauterium corium corrumpit et notam in- figit, ita et fallatia conscientiam quae dolo maleuolen tiae aliud scit et aliud profitetur, denotat ad perditio nem. Haec Ambro. Iam et malefici candenti ferro hoc est cauterio notantur apud nos stigmatis impressis cu- ti. Significauit ergo apostolus falsos illos doctores ani mum habituros sceleris conscientia et ardore ambustum, qui conquiescere nesciat et perpetuis sese excruciet tormentis, ut ut foris sancte et hylariter componant et faciem et totum corporis habitum.  \pend
\phantomsection
\addcontentsline{toc}{subsection}{\textit{prohibentium contrahere ma- trimonium, iubentium abstinere a cibis quos Deus creauit ad su- mendum cum gratiarumactione fi- delibus, et his qui cognouerunt ue- ritatem quod quicquid creauit De- us bonum sit et nihil reiiciendum si cum gratiarumactione sumatur. }}
\subsection*{\textit{prohibentium contrahere ma- trimonium, iubentium abstinere a cibis quos Deus creauit ad su- mendum cum gratiarumactione fi- delibus, et his qui cognouerunt ue- ritatem quod quicquid creauit De- us bonum sit et nihil reiiciendum si cum gratiarumactione sumatur. }}
\section*{AD TIM. CAP. IIII. }
\marginpar{[ p.138 ]}
\phantomsection
\addcontentsline{toc}{subsection}{\textit{Sanctificatur enim per sermonem Dei ac precationem. }}
\subsection*{\textit{Sanctificatur enim per sermonem Dei ac precationem. }}\pstart Iam speciatim subiungit quaedam dogmata ex qui bus quae iam proposuit manifesta fiant et reliquae ipso rum doctrinae aestimentur. Bina autem exempla pro- ponit, alterum quidem quod coniugio et alterum quod cibis interdicturi sint. De utroque quid sentiat ueritas nemo nescit. De coniugio pronunciat scriptura Hono rabile esse connubium apud omnes et cubile impollutum. De cibo eadem censet Puris omnia esse pura, neque id quod in os ingreditur polluere hominem, sed quod ex corde egreditur. In qua simplicitate ut uera fides per stat ita hypocrisis rimas quaerit per quas elabatur. Ac cedit hic spiritus impostor, et quoniam is bona specie imposuit primis parentibus, ideo nunc quoque suggerit egregium esse si quis tam a mulieribus quam a carni- bus et uino abstineat, inde enim enasci in mundo glo riam et laudem maximam. Hanc qui spectant a fide deficientes spiritibus impostoribus et daemoniorum doctrinis attendunt, labunturque in hypocrisim et mendacium. Tametsi enim abstineant tam a mulierum complexu quam a carnium esu, animus tamen insyn- cerus est. Nam multarum sibi conscii sordium consci- entiam habent inustione uitae impurioris ineluibili compunctissimam. Foris quidem non uescuntur carni  \pend
\vspace{0.5cm}\noindent
\vspace{0.2cm}\rule{1cm}{0.2pt}\\ 
\hspace{0.2cm}\textit{mg}
\footnotesize Dogmata pseudopro phetarum. 
\normalsize| 
\section*{COMMENT. IN I. EPIST. }\pstart bus, intus uero ardent ira, inuidia, aemulatione et contentione. Abstinent hi ab uxoribus, sed animus interim aestuat foedis cupiditatibus. Deus autem exter- na specie hoc est hypocrisi non placatur, animum poscit sibi purgatum uitiis offerri. Sathan ergo ut a ue ro abducat et erroribus implicet nouas huiusmodi do ctrinas inspirat. Exposuerunt hunc locum alii de Ma nicheis et Tatianis quos et Encratitas appellant. De his enim Augustinus ad Quoduultdeum de haeresi- bus scribit quod damnarint nuptias atque has omnino pares fecerint fornicationibus, quodque nullas ederint carnes imo usum eorum damnarint prorsus. Egoue- ro uideo Pauli uerba non de una aliqua gente aut fa- milia haeretica, non de una aetate, sed de omnibus aeta- tibus et hominibus a fide deficientibus, spiritibus im postoribus attendentibus prohibentibusque nuptias et cibos, esse exponenda. Intelligo ergo apostolum etiam haec de Apostolicis, quos alio nomine Apotactitas ap- pellauit uetustas, quasi quis dicat extraordinarios, lo- quutum. Hi enim Augustino authore in communio- nem suam non recipiebant utentes coniugibus et res proprias possidentes. Interpretor haec et de fide Ma- humetica et Turcis. Praecipue autem expono ea de Pontifitiis et nostri saeculi religiosis uirisque Essenis si ue Monasticis. Manichei enim et Encratitae una cum Apotactitis damnarunt nuptias et ciborum usum ma  \pend
\vspace{0.5cm}\noindent
\vspace{0.2cm}\rule{1cm}{0.2pt}\\ 
\hspace{0.2cm}\textit{mg}
\footnotesize Est in uoc hac Cata- chresis. 
\normalsize| 
\section*{AD TIM. CAP. IIII. }
\marginpar{[ p.139 ]}\pstart gis quam prohibuerunt. Pontifitii uero prohibuerunt non damnarunt. Caeterum fides et religio syncera quoa iam monuimus tradit liberum esse homini cuilibet du- cere uxorem aut nubere uiro, hoc uel illo cibo uesci. At ab ista simplicitate fidei defecerunt nostri specieque ́ delusi sanctimoniae, magnifica suasore sathana et iu- bente hypocrisi de coelibatu, de uotis, de uirginitate, de ieiuniis, de delectu et abstinentia ciborum tantum non omnem his salutem tribuentes disputarunt, imo et leges promulgarunt, quas uiolare aut refigere uel no- tare saltem licebat nemini. In ficta autem haec sancti- monia uiuentes ceu sancti ac diui quidam coelitus de- lapsi adorabantur a nobis, quod foris omnia ipsorum ad ostentationem essent instructissima et apparerent purissima. Domi interim quales fuerint ipsi melius no- runt, et nos non nescimus ipsorum conscientias fuisse adustas et multa immunditia pollutissimas, Nihil hic fingo, nihil ex suspitione effutio, ipsa experientia lo- quitur quid factum sit et quid etiam hodie fiat. Con- templare enim mihi Essenos nostros hoc est Cartusia- nos, contemplare omnes in uniuersum monasticos ui- ros et mulieres, adde his et sacrificos pontifitios, et cogeris fateri in tanta turba otiosorum hominum non inueniri multos uere continentes. Iam ut maxime inueniatur unus atque alter qui uel ob dei donum rarum, uel ob uires effoetas aut ualetudine fractus uiuat conti-  \pend
\vspace{0.5cm}\noindent
\vspace{0.2cm}\rule{1cm}{0.2pt}\\ 
\hspace{0.2cm}\textit{mg}
\footnotesize In Mona- chos et uo ta monasti- ca. 
\normalsize| 
\section*{COMMENT. IN I. EPIST. }\pstart nentius, uidemus interim superesse multas hypocrita- rum myriades qui nihil minus sint quam quod audi- unt coelibes, monachi, nonnae, spirituales, monastici, puri, sancti. De optimis hic loquor, nihil de impudenti bus dico nebulonibus qui palam uel scortantur, perpo tant, alea ludunt, belligerantur, adulteriis polluunt omnia, quorum maior propemodum numerus, uel foe diora et abominabiliora designant, indigni enim sunt quorum fiat aliqua mentio honesta apud uiros bonos: de mundioribus hic magisque reuerentibus opti misque dissero quod ne horum quidem conscientiae syn cerae sint. Nam purum esse nequit quod sine fide fit. At hic Defecerunt ait Paulus a fide. Et cum haec mani festiora sint quam quae negari possint, inuenias tamen quos non pudeat damnata illa praecepta de ciborum delectu deque coelibatu immundo, deque uotis monasticis angelica nominare praecepta:rectius uero Paulus ap- pellauit sathanica. Scio quam insuauem hic multorum auribus cantilenam canam, sed adulari non didici, ue- ritatem dissimulare non potui, neque alia hic quam Pau lina uerba usurpaui. Proculdubio resipiscent quibus cordi est uera pietas agentque monitori gratias. Intelli- gent enim me hac in disputatione non damnare uel uir ginitatem et castimoniam sanctam uel ieiunia mode rata, castigationem carnis iustam et abstinentiam syn ceram, sed hypocrisim magis et impuram illam tem-  \pend
\section*{AD TIM. CAP. IIII. }
\marginpar{[ p.140 ]}\pstart perantiam. Sed redeo nunc ad Paulum, qui copiosius sed in transcursu demonstrat liberum esse fidelibus quouis cibo quolibet uesci tempore. Deus enim, ait, cibos creauit (arguens a prima dei inftitutione et uo- luntate) in usus hominum ut hos scilicet fideles su- mant cum gratiarumactione. Ad cuius interpretatio- nem pertinet quod mox sequitur, His inquam qui co- gnouerunt ueritatem. Hi enim fideles sunt. At ueritas ipsa quae a fidelibus agnoscitur hoc in negotio est, quod quicquid creauit Deus bonum est nec quicquam in cibis abominandum et reiiciendum. Quaecunque  enim condidit deus bona sunt, Gen.1. Abusus cibos im puros siue malos reddit. Proinde Paulus de sancto eo- rum usu scribensiubet hos sumere cum gratiarumacti one, et sententiam elegantem subiungens dicit, San- ctificatur enim per sermonem dei ac precationem. Vbi non est quod confingas tibi benedictionem quan- dam anilem conceptis ex uerbo dei sententiis pronun ciatam super cibum, sed Sanctificari purificari est et ab usu separari prophano. Deinde Verbum dei institu tio dei est illaque sententia qua pronunciat Puris omnia esse pura, omnia item quae ab ipso condita sunt esse bona inque usum instituta hominum. Per precationem autem intelligit gratiarumactionem et usum fidelem. Ita sententia erit, ipsum enim domini uerbum, ipsa in quam dei concessio et sententia qua pronunciat om-  \pend
\vspace{0.5cm}\noindent
\vspace{0.2cm}\rule{1cm}{0.2pt}\\ 
\hspace{0.2cm}\textit{mg}
\footnotesize Libertas et delectus ci borum. 
\normalsize| 
\section*{COMMENT. IN I. EPIST. }\pstart nia esse puris pura, fides item recipientis et gratia- rumactio fidelis facit ut nihil impuritatis habeat cibus Christianorum. Atque haec expositio uidetur omnium esse simplicissima maximeque genuina. Paraphrastes uero ex interpretibus uetustis hanc Pauli sententiam eadem propemodum ratione in haec uerba explana- uit, Quod etiam si quid impuritatis haberet cibus, ta- men hymnis quibus laudatur diuina largitas sacrisque ́ uerbis ac precationibus, sanctum purumque redde- tur quod prius erat immundum. Concidunt ergo hie semel quaecunque mercatores Romani praeceperunt de ciborum delectu, de temporum et dierum discrimine. Christiano enim licet quouis tempore quibuslibet ue- sci cibis. Quatenus uero scandali ratio habenda sit o- stensum est in cap.14. ad Roma. Discimus praeterea quanta cum reuerentia ad mensam accedere et abea dem debeamus recedere. Conuiuae domini sumus, do- minus pascit nos, in huius oculis uersamur, illius mu- nificentia alimur. Precandum ergo et gratias agen- das esse ante et post cibum nemo est qui non uideat, sed et ita in mensaut in oculis dei uersandum. Aberit itaque omnis amaritudo, contumelia, calumnia, rixa, obtrectatio, denique luxus et ebrietas obscoenitas et intemperantia omnis, nil foedum dictu aut auditu tur- pe in mensa proponatur sed quod honestum et iu- cundum sit. Reuerentia hofpitis qui praesens est et  \pend
\section*{AD TIM. }
\marginpar{[ p.141 ]}\pstart nos alit perpetuo retineat in officio et disciplina san cta. Haec autem de pseudoepiscopis et peruersa imo daemonica doctrina paucula quidem in loco obiterque ́ inseruit quando de officiis boni pastoris disputauit, re dit nunc ad institutum suum, et pergit iusta episcopi munia informare, praecipue praescribens quanta uigi- lia et assiduitate episcopus quisque doctrinam inculcet sanam, daemoniacam uero caueat atque impugnet.  \pend
\phantomsection
\addcontentsline{toc}{subsection}{\textit{OFFICIA ET INSTITVTIO Episcopi. }}
\subsection*{\textit{OFFICIA ET INSTITVTIO Episcopi. }}
\phantomsection
\addcontentsline{toc}{subsection}{\textit{De his si commonefeceris fra- tres bonus eris minister lesu Chri- sti, enutritus in sermonibus fidei bonaeque doctrinae, quam usque se- quutus es. Caeterum prophanas et aniles fabulas reiice, quin potius exerce temetipsum ad pietatem. }}
\subsection*{\textit{De his si commonefeceris fra- tres bonus eris minister lesu Chri- sti, enutritus in sermonibus fidei bonaeque doctrinae, quam usque se- quutus es. Caeterum prophanas et aniles fabulas reiice, quin potius exerce temetipsum ad pietatem. }}\pstart ταῦτς, inquit ὐποτιθέαενθ id est, Haec sugge rens atque proponens fratribus. Haec inquam quae de ue ra doctrina, de ecclesiastico coetu sancto, de moribus episcopi et studiis, de incarnato uerbo et iuftificato in spiritu, de pietatis precio et certitudine, deque ca- uenda daemoniorum doctrina, de sancto item coniu-  \pend
\section*{COMMENT. IN I. EPIST. }\pstart gio et sancto ciborum usu, haec inquam si tradideris bona fide fidelibus bonus eris minister Christi. Hone- sto enim et bono excitat ad uigiliam et fidelem do- cendi diligentiam. Aequum est, inquit, ut fideliter do mino qui te ad hoc ministerium uocauit, seruias. Id fe- ceris autem si ista quae praescripsi proposueris fideli- bus. Hodie uero haereseos insimulamur quoties ista ita ut Paulus praecepit tradimus ecclesiis: unde coniicere licet quales habeamus aduersarios, quales item doctri nae nostrae praeceptores. Deinde accedit aliud quoquam quod et in officio retineat Timotheum et excitet ad maiora. Ab ipsa adolescentia inquit in fide euange lica bonaque doctrina es ueluti enutritus, sanam quoque  doctrinam et fidem ueram es hactenus constanter se quutus, ideoq perges porro summa cum laude quod cepisti absoluere, tui per omnia similis esse, imo temet ipsum in studio isto pulcherrimo superare plurimum. His iam addit praeceptum breuissimum in quo comple ctitur omnia, Ad pietatem temetipsum exerce, pro- phanas et inanes fabulas reiice. De pietate autem quid sit ne inter Gentes quidem qui non intelligerent defuerunt. M. enim Cicero in oratione pro domo sua ad Pontifices non impie de pietate sentiens, Equidem, inquit, sic accepi Pontifices in religionibus suscipien- dis caput esse interpretari quae uoluntas deorum im- mortalium esse uideatur. Nec est ulla erga deos pietas  \pend
\vspace{0.5cm}\noindent
\vspace{0.2cm}\rule{1cm}{0.2pt}\\ 
\hspace{0.2cm}\textit{mg}
\footnotesize 2. Timo.3 
\normalsize| 
\hspace{0.2cm}\textit{mg}
\footnotesize Pietas quid 
\normalsize| 
\section*{AD TIM. CAP. IIII. }
\marginpar{[ p.142 ]}\pstart nisi honesta de numine eorum ac mente opinio, quum expeti nihil ab iis quod sit iniustum et inhonestum arbitrere. Haec Cicero. Et sane uox ἐυσέθεια quae La tinis Pietas dicitur ex etymo sonat bonus siue uerus cultus, religionem ueram uertit alibi D. Hierony. Et Aurel. Augustinus epist. ad Macedonium 52.Pietas est inquit, uerax ueri dei cultus, unde omnia recte uiuen di duci oportet officia. Et ad Hierony. epist.29. Pietas ait, quam Graeci uel ἐυσέθααpuel expressius et ple nius θροσέθεια p uocant, aliud non est quam iustus dei cultus. Si ergo pietas uoluntatis diuinae iusta est cogni tio, si uerus dei cultus est, Voluntas autem dei est ut credamus in eum quem misit ille, et ut mutuo diliga mus inuicem, cultus idem dei consistat in fide et cha- ritate, necesse estut apostolus in fide et charitate uo- luerit nosmetipsos exerceri cum in pietate exerceri praecepit. Ita consequetur etiam prophanas et aniles fabulas esse quaecunque a fide et charitate aliena sunt. Prophanum enim est quicquid a sacrosancta fide ali enum est. Et apud Graecos quoque usurpatum est pro- uerbium γραῶμ ὐθλος, quod Latinis sonat Anicu- larum deliramenta. Id quod dici solebat de nugis ina- nibus cuiusmodi effutiunt uetulae, cum iam sexus uitio uitium aetatis accedens nugacitatis morbum conduplicat. Est apud nos quoque frequenti in usu ut de rebus nihili et uanissimis dicamus, Is sye ollter myber tanot.  \pend
\vspace{0.5cm}\noindent
\vspace{0.2cm}\rule{1cm}{0.2pt}\\ 
\hspace{0.2cm}\textit{mg}
\footnotesize Ioan.6. 1.loan.3e 
\normalsize| 
\section*{COMMENT. IN I. EPIST. }
\phantomsection
\addcontentsline{toc}{subsection}{\textit{Nam corporalis exercitatio pau- lulum habet utilitatis, at pietas ad omnia utilis est, ut quae promissio- nes habeat uitae praesentis et fu- turae. }}
\subsection*{\textit{Nam corporalis exercitatio pau- lulum habet utilitatis, at pietas ad omnia utilis est, ut quae promissio- nes habeat uitae praesentis et fu- turae. }}\pstart Causam siue rationem adiungit praeceptioni. Ideo inquit iubeo imprimis et assidue pietatem inculcare, quod haec sit ad omnia utilis, exercitatio uero corpo- ralis paululum habeat utilitatis. Pietas autem quid sit iamiam exposui. Caeterum per Exercitationem corpo ralem Ambro. intellexit ieiunia et omnia carnis fre- na, disciplinas uidelicet externas quibus ferocientem edomamus carnem. Et sane uoluit Paulus semel omnem illum exteriorem cultum et disciplinam notare per corporalem exercitationem ut hic non tantum ieiu- nia, uigilias, delectum ciborum et uestium, corporis castigationem, sed omne intelligat quod externa ob- seruatione in corpore fit. Haec quidem si moderate et in fide fiant non prorsus inutilia sunt. Praeparant e- nim corpora nostra ad officia pietatis, uerum in eis non consistit solida pietas, imo nisi prius adsit pietas ista proderunt ad nihilum. Sed et omnes dei in scri- pturis promissiones fidelibus et charitatem per ope- ra pietatis exprimeutibus pollicentur tam praesentis  \pend
\section*{AD TIM. CAP. IIII. }
\marginpar{[ p.143 ]}\pstart bona quam futurae et aeternae uitae gaudia. Nullibi le- gas quemquam ob corporalem exercitationem apud deum esse iustificatum. Fide placuerunt deo quotquot a mundi exordio placuerunt. Et charitatis opera non corporis exercitia exiget dominus cum uenerit iudi- care uiuos et mortuos. Denique et sancti patres omnes approbarunt se deo fide innocentia et charitate et non ceremoniis et disciplina externa. Id quod Ste- phanus quoque demonstrat Acto. cap.7. Vbi ergo nunc sunt praecepta Monachorum de cultu dei et uitae san ctimonia? Confixa hic ruunt omnia, neque enim aliud sunt quam prophanae et aniles fabulae. Paraphrastes ea quae de pmissionibus scripsit Apostolus paulo diuersi us exposuit. Cuius uerba si quis requirat haec sunt, Cae terum euangelica pietas quae nititur syncera fide ueraque ́ charitate nullo non tempore in nullauitae parte non estu- tilis quae compendio profitetur quicquid uel in prae- senti tempore sit expetendum uel sperandum in futu- ra, adeo ut non sit opus ullum aliunde praesidium ue- nari nobis. Quae expositio ut superiori non repugnat ita nihil habet impietatis. Deligat ergo lector pius quod sequatur.  \pend
\phantomsection
\addcontentsline{toc}{subsection}{\textit{Indubitatus sermo dignusque ́ qui omnibus modis approbetur. Nam in hoc et laboramus et pro- }}
\subsection*{\textit{Indubitatus sermo dignusque ́ qui omnibus modis approbetur. Nam in hoc et laboramus et pro- }}
\vspace{0.5cm}\noindent
\vspace{0.2cm}\rule{1cm}{0.2pt}\\ 
\hspace{0.2cm}\textit{mg}
\footnotesize Hebtae. ni. 
\normalsize| 
\hspace{0.2cm}\textit{mg}
\footnotesize Matth.15. 
\normalsize| 
\section*{COMMENT. IN I. EPIST. }
\phantomsection
\addcontentsline{toc}{subsection}{\textit{bris afficimur quod spem fixam ha- bemus in Deo uiuente qui est ser- uator omnium hominum maxi- me fidelium. }}
\subsection*{\textit{bris afficimur quod spem fixam ha- bemus in Deo uiuente qui est ser- uator omnium hominum maxi- me fidelium. }}\pstart Confirmat quod dixit pietatem esse utilem ad o- mnia ut quae habeat promissiones praesentis ac futurae uitae, Quod itaque proposui inquit de pietate uerum est et indubitatum planeque dignum quod recipiatur ab omnibus. Quod autem uerbis tantopere inculcamus et pro indubitata ueritate uenditamus apud omnes hoc ipsum re ipsa corpore et afflictionibus nostris in dubitatum esse comprobamus. Nisi enim persuasissi- mum mihi esset in solo deo uiuente salutem esse quam nobis Christi contulerit meritum, certe non tam infra cto animo huius uitae aerumnas, carceres, impiorum plagas et uarias mortes perpeterer. Nunc autem post habita exercitatione corporali, neglectis ceremoniis et umbris Iudaicis imo reiectis semel omnibus quae mundus miratur externis solam euangelii pietatem praedico docens Christum dominum uenisse non uni a- licui genti sed seruandis omnibus qui modo perfidia seipsos non fecerint indignos hac salute. Sola enim fi- de haec percipitur, charitate exprimitur. Atque huius quidem doctrinae gratia multa perpetior, sed uolens et lubens, ita simul testificans uerum esse quod doceo.  \pend
\section*{AD TIM. CAP. IIII. }
\marginpar{[ p.144 ]}\pstart Tam uidelicet certa est pietatis doctrina, ut huic meri to cedant quaecunque de exercitatione corporali prodi ta sunt a superstitiosis hominibus. Haec enim certa imo indubitata sunt, idcirco digne iam subiungit.  \pend
\phantomsection
\addcontentsline{toc}{subsection}{\textit{Praecipe haec et doce. Nemo tu- am iuuentutem despiciat, sed esto forma fidelium in sermone, in con- uersatione, in dilectione, in spiritu, in fide, in puritate. Donec uenero attende lectioni, exhortationi do- ctrinae. }}
\subsection*{\textit{Praecipe haec et doce. Nemo tu- am iuuentutem despiciat, sed esto forma fidelium in sermone, in con- uersatione, in dilectione, in spiritu, in fide, in puritate. Donec uenero attende lectioni, exhortationi do- ctrinae. }}\pstart Infert nunc quod potissimum hac intendit disputa- tione. Haec, inquit, tam uidelicet certa et indubitata adeoque ; ab ipso profecta Christo, fideliter et constan- ter praecipe, hoc est multa et matura cum authorita te trade et sic trade ut totus in hoc sis quo omnes in- telligant solam pietatem esse ad salutem consequen- dam utilem. Ista uero omnibus quoque nobis dicta puta- bimus quicunque diuina ordinatione praesumus eccle- siis Christi. Caeterum cum communis quaedam sit et praesumpta opinio contemptibilem esse adolescentiam et Timotheus admodum tenere esset adolescentiae, per occupationem subiungit, Nemo tuam iuuentutem despiciat. Quasi dicat, Neque est quod iuuenis licet in  \pend
\vspace{0.5cm}\noindent
\vspace{0.2cm}\rule{1cm}{0.2pt}\\ 
\hspace{0.2cm}\textit{mg}
\footnotesize Doctrinz episcopi sa na. 
\normalsize| 
\hspace{0.2cm}\textit{mg}
\footnotesize Adolescen tia cana non contemnen da. 
\normalsize| 
\section*{COMMENT. IN I. EPIST. }\pstart re tanta uerearis aut cedas improbitati diuersa docent tium. Alibi sit humanitati locus, hic ubi salutis pericu lum uertitur proferenda est authoritas. Neque spectan dum est quot annos uixeris sed quod munus geras. Hic senex est quisquis uitae integritatem quisquis morum praestat grauitatem: ita ergo uersandum tibi atque ui- uendum est ut nemo iuuentutem tuam possit despice. re. Id autem qua ratione fiat consequentibus exponit. Si uidelicet ita se componat ut fideles in ipso habeant exemplar ἐμ λόγῳ in sermone uel doctrina. Hoc est si nihil aliud doceat quam pietatem ueram. Deinde si conuersatio eius sancta sit, pudica, sobria, et irrepre hensibilis, quae doctrinae respondeat per omnia. Hine si dilectione praeditus sit haud uulgari sed eximia ex qua faciat omnia officia. Item si non desit spiritus pru- dentiae fortitudinis et constantiae. Quibus si accedat puritas morum et animi et maturitas nihil habens commixtum iuuenilium ineptiarumm cupiditatumue, efficient quo minus aetas alians suspecta ueniat in con- temptum. Ex quibus colligimus immerito quosdam ho die conqueri iacere authoritatem ministorum uerbi ut quos rudis plebecula despiciat nimium, Cur enim non obsequuntur Pauli consilio et uitam doctrinamque ́ suam ita instituunt ut merito non possint contemni a quoquam. Propria sane ipsorum culpa et inertia fit quod contemnuntur cum uon modo prophani sine  \pend
\section*{AD TIM. CAP. IIII. }
\marginpar{[ p.145 ]}\pstart sed indocti quoque et ab omni germana pietate alieni. Discimus praeterea hominum ingenia non a uetustate ut uina esse aestimanda, sed ab ipsa probitate et erudi tione. Iuuenes fuerunt Daniel, Dauid, Samuel, Ieremi- as, et alii multi in utroque testamento egregii uiri. Im probantur autem in sacris literis non semel senes impru dentes et impudici. Interim uero non negamus uitae usu multaque rerum experientia parari prudentiam, unde etiam agnoscimus sanctam senectutem sancte esse colendam. Iam et studia in quibus se exerceant episcopi egregie praescribit et informat. Nullam ue- ro hic agriculturae mentionem facit, neque episcopos manuariis sellulariisque artibus aut opifitiis affigit, quo magis miror quae illos hodie agat daementia qui ma- lunt commendari se a stiua, malleo, dolabra aut lima, quam ab assiduo sacrarum et honestarum literarum studio. Obiicient fortassis Paulum quoque exercuisse opifitium manuarium. Sed cur non expendunt quod reliqua Pauli dona nondum assequuti sunt. Paulus e- nim summam eruditionem summae coniunxerat elo- quentiae, et ut uas dei selectissimum erat ita nulla sa- crarum literarum nesciebat mysteria. Primitias enim spiritus una cum apostolis reliquis acceperat. Isti ue- ro egregii si Diis placet Pauli imitatores Musas ac li- teras sacras nondum a limine salutarunt, pariter et uulgari eruditione et eloquentia destituti sunt, quae  \pend
\vspace{0.5cm}\noindent
\vspace{0.2cm}\rule{1cm}{0.2pt}\\ 
\hspace{0.2cm}\textit{mg}
\footnotesize Studia et exercitia e- piscopi. 
\normalsize| 
\section*{COMMENT. IN I. EPIST. }\pstart interim si saperent assiduo assequi studio conaren- tur. Quod si quicquam apud eos ualet Paulina autho- ritas quare non audiunt praecipientem Attende lecti oni exhortationi doctrinae? Vere enim Theophyla- ctus, si, inquit, Timotheo iniungit ut sacris literis ua cet, quid nobis erit agendum? Constat enim Timothe- um a cunis in sacris uersatum literis iustam habuisse apud omnes absolutae eruditionis laudem. Et Graeca Latinis multo sunt euidentiora. πρόσεχε, inquit, ri ἀναγνωσα. Est autem ἀναγνωσις lectio quidem sed non quaelibet imo accuratissima et certissima quasique ; persuasio quaedam. Et πρoσέχοͅ adhaereo est, dedico, intendo, adeoque animum deuoueo. Praecepit ergo a- postolus ut se totum consecret sacris studiis, nec ullam rem aliam sibi perinde habeat commendatam atque le- ctionem et cognitionem rerum sacrarum. Proinde grauiter peccant qui succisiuis modo horis sacra per- tractant ministri, iustas horas et sese totos impendunt opifitiis. Imo arbitror apostolum Paulum succifiuas mo do horas elocasse artifitio, iustas sacrarum studio. E- quidem succisiuas aut secundarias horas impendisse uel opifitio uel exercitio alicui honesto laudamus pluri- mum. Caeterum praesentis instituti non est pluribus de lectione sacra disserere, cum quod alio uocet me insti- tutum meum, tum quod alii de hac re melius scripse runt. Noluit autem apostolus hoc episcopi studium e£  \pend
\section*{AD TIM. CAP. IIII. }
\marginpar{[ p.146 ]}\pstart priuatum, hinc adiecit exhortationem et doctrinam, quibus speciebus complexus est omne inftituendi ge- nus. Hunc ergo finem habeant studia nostra ut quam pluribus prosimus, utque pro rei negotio uel docere, uel hortari, uel accusare, uel consolari possimus. Haec au- tem sunt propria episcoporum officia.  \pend
\phantomsection
\addcontentsline{toc}{subsection}{\textit{Ne neglexeris quod in te est do- num quod datum est tibi per pro- phetiam cum impositione manu- um authoritate sacerdotii. }}
\subsection*{\textit{Ne neglexeris quod in te est do- num quod datum est tibi per pro- phetiam cum impositione manu- um authoritate sacerdotii. }}\pstart Admonet Timotheum officii sui currentique admo uet calcar. Multa sunt inquit quae admonent ut gnaui ter prouinciam quam nactus es administres. Habes exi mium et ad hoc negotium appositissimum a deo da- tum donum, quod proculdubio excitabis ni malis ac- cusari neglectus et inertiae. Accedit iusta ad episcopa tum uocatio. Neque enim temere ingessisti te, quod mul ti solent, sed prophetia et arcanus quidam afflatus id muneris tibi imposuit teque ad hoc sacrum ministerium designauit. Accedit praeterea et manuum seniorum impositio, quae ista comprobauit et publico muneri publice praefecit. Haec inquam omnia admonent ut su- scepto muneri respondeas per omnia. Colligimus au- tem nos ex hoc Pauli loco tria in cuiusuis episcopi or- dinatione necessaria esse. Primum ut habeat donum  \pend
\section*{COMMENT. IN I. EPIST. }\pstart dei, hoc est ut sit irreprehensibilis bonarum et lite- rarum et rerum non tam studiosus quam peritus. De inde ut iuste et sancte hoc est secundum uerbum dei eligatur ac uocetur. Postremo ut delecto et ecclesiae quod aiunt praesentato imponantur manus. Hinc Ti- gurina ecclesia posteaquam quaestuosae et inordinatae episcopi Pontifitii ualedixisset ordinationi, uiros e do ctis et uerbi ministris, e senatoribus item et diacosi- is, hoc est e plebe aliquot delegit qui ex optimis hoc est honestissimis atque doctissimis diaconis ( quoties ec- clesiae alicui sanctus praeficiendus est episcopus) cer- tos senatui et plebi statuant. Ex his uero unum ali- quem libera electione designant, quem protinus eccle siae cui praeficiendus est mittunt et commendant per legatum ac senatorem honestissimum. Atque hic quidem praestantissimus eius ditionis loci aut urbis antistes pro publica contione perorat, fiunt ab ecclesia ad Deum preces publicae, interque precandum et in conspectu totius populi imponit antistes ille sacrorum manus fu turo illi episcopo cui et ecclesiam commendat fideli- ter. Haec uero plurimum momenti habere et episco- po plurimum authoritatis conciliare nemo est qui ne- sciat. Certe cum ea fiunt ex institutione domini fiunt. Et apostolus hinc magna petiit argumenta quibus Ti- motheum sui admoneret officii. Eadem optimum quem- que non poterunt non commonere grauiter quoties  \pend
\vspace{0.5cm}\noindent
\vspace{0.2cm}\rule{1cm}{0.2pt}\\ 
\hspace{0.2cm}\textit{mg}
\footnotesize Mos Tigu rinae eccle- siae in ordi nandis epi scopis. 
\normalsize| 
\section*{AD TIM. CAP. IIII. }
\marginpar{[ p.147 ]}\pstart ste secum expendit qua quidque ratione quóue consie lio fiat. Iam et Catabaptistae ultro et sua sponte cur- rentes hac ordinatione proturbabuntur ualidissime, si ea modo sancte et iuste in ecclesiis obseruetur. Illud quoque priusquam hinc digrediar monebo, hunc Pauli locum sic quoque reddi posse, Ne neglexeris quod in te est donum, nempe quod tibi delegatum est ad prophe tandum, idque comprobatum manuum impositione a se nioribus. Verum nullus hinc alius colligetur sensus quam quem modo indicauimus.  \pend
\phantomsection
\addcontentsline{toc}{subsection}{\textit{Haec exerce, in his esto, ut tuus profectus manifestus sit in omni- bus. Attende tibi ipsi et doctrinae, persiste in his. Nam id si feceris te ipsum seruabis et eos qui te audi- erint. }}
\subsection*{\textit{Haec exerce, in his esto, ut tuus profectus manifestus sit in omni- bus. Attende tibi ipsi et doctrinae, persiste in his. Nam id si feceris te ipsum seruabis et eos qui te audi- erint. }}\pstart Epilogus est superiorum. ταυτα inquit μελὲτα, Haec meditare. Est autem meditari tam animo quam corpore adeoque totis uiribus in rem aliquam perficien dam incumbere. In his esto, hisce te exerceto iugiter, constanter, et uehementer, adeo ut in tota plebe ap- pareat tua sedulitas et fidelitas, hoc est, ut ex pia do- ctrinauitaque sancta omnes reddantur meliores. Finis enim huius negotii est ut prosimus non ut gloriam no  \pend
\section*{COMMENT. IN I. EPIST. }\pstart stram quaeramus. Imo non negligit sese qui gnauiter suum facit officium, negligit qui segniter suam prouin ciam administrat. Vnde protinus sequitur in Paulo, Attende tibi ipsi et doctrinae, persiste in his. Germa- nus uerteret, Mab gut adt sorg vnnd flyss vff oique selbe vno nff die leer, ut uidelicet doctrinae non repugnet uita. Addit utilitatem et praemium, Hinc e- nim inquit geminum capies fructum. Primum serua- bis teipsum, deinde eos qui te audierint. Atque haec qui dem tanti sunt momenti ut nihil supersit in orbe quod eum mouere possit quem ista non mouent.  \pend
\endnumbering\beginnumbering\section{CAP.V.}
\phantomsection
\addcontentsline{toc}{subsection}{\textit{Seniorem ne saeuius obiurges, sed adhortare ut patrem, iuniores ut fratres, mulieres natu grandio- res ut matres, iuniores ut sorores cum omni castitate. }}
\subsection*{\textit{Seniorem ne saeuius obiurges, sed adhortare ut patrem, iuniores ut fratres, mulieres natu grandio- res ut matres, iuniores ut sorores cum omni castitate. }}\pstart Iam praescribit Episcopo formulas quasdam gene- rales quomodo se gerat in docendo et exhortando erga cuiusuis aetatis et ordinis homines, item quomo- do se gerat erga uiduas presbyteros dominos ac ser- uos. Principio uult doctrinam exhortationem adeoque  obiurgationem episcopi ab omni tyrannidis specie quam alienissimam esse, cum quod iusta quaedam leni tas, quae uidelicet ad molliciem et sceleris impunita-  \pend
\vspace{0.5cm}\noindent
\vspace{0.2cm}\rule{1cm}{0.2pt}\\ 
\hspace{0.2cm}\textit{mg}
\footnotesize Opiurga- tio episcopi 
\normalsize| 
\section*{AD TIM. CAP. V. }
\marginpar{[ p.148 ]}\pstart tem non deflectit, deceat episcopum, tum quod inge- nium hominis ferme sic est ut malit duci quam cogi. Se niori ergo, ait, μη ἐπιπ λήξμς id est ne plagam infli- xeris, quod perinde est ac si dixisset Seniorem ne aspe re et saeuius obiurges. Est enim in uerbo metaphori- ca deflexio translatióue. Et Latini uerbera linguae usur pant pro maledicentia saeua. Reliqua uero plana sunt. Postremam particulam ἐμ πάση αγνεία In omni pu ritate, quidam referunt ad mox praecedens membrum, quasi uoluerit abesse omnem impudicitiae suspicionem: ut sit sensus, luniores castigabis sic ut sorores, idque pu- ritate tantaut impudicitiae nulla uel suspitio appare- at. Alii uero mea quidem sententia multo melius retu lere ad omnia praecedentia membra, ut sermo episco- pi siue cum senioribus agat, siue cum iunioribus, siue cum puellis purus sit ab omni specie uel saeuitiae uel contumatiae uel auaritiae uel odii.  \pend
\phantomsection
\addcontentsline{toc}{subsection}{\textit{Viduas honora, quae uerae ui- duae sunt. Quod si qua uidua libe- ros aut nepotes habet, discant pri- mum propriam domum pie tra- ctare et uicem rependere maiori- bus. Hoc enim est honestum et ac- ceptum coram deo. }}
\subsection*{\textit{Viduas honora, quae uerae ui- duae sunt. Quod si qua uidua libe- ros aut nepotes habet, discant pri- mum propriam domum pie tra- ctare et uicem rependere maiori- bus. Hoc enim est honestum et ac- ceptum coram deo. }}
\section*{COMMENT. IN I. EPIST. }\pstart Nunc praecipit de Viduis et copiosius quidem quod plurimum fraudis hic committeretur a quibusdam E- phesiorum uiduis. Constat autem cum ex σ. cap. Act. apost. tum ex praesenti loco ecclesiae munificentia et liberalitate in reliqua egenorum turba uiduas esse ali- tas. Quemadmodum uero in optimis quibusque institutis hominum auaritia et malitia suum quaerit compendi um, ita nimis impudenter in egenorum turbam immi- scebant se hic quoque citra delectum uiduae omnes, unde apostolus coactus est paulo circumspectius de his cu- randis et suscipiendis disserere. Principio in genere praecipit, Veras honora uiduas. Deinde uero per speci es quasdam exponit quae sint uiduae uerae, quae suscipi endae, quae relegandae. Poterit hinc quoque certum for- mari iudicium De alendis pauperibus, et qui digni qui item indigni sint alimonia et stipendio ecclesiasti co. Illud quoque hinc euidentissime probari potest quod ad episcopos cura pauperum pertineat, non quod o- rationem et praedicationem uerbi dominici iubean- tur deserere et mensis ministrare, sed quod officii eo- rum sit diligenter inuigilare ut illi qui pauperibus prae fecti sunt pauperes non negligant aut fraudent sed bo na fide curent: uerum redimus ad uerba apostoli. Vi- duas inquit honora, id est consilio auxilio et subsidio iuua ecclesiastico. Sic enim ipse apud Matth, dominus in cap. is5. uerbum honorare exposuit. Significanter  \pend
\vspace{0.5cm}\noindent
\vspace{0.2cm}\rule{1cm}{0.2pt}\\ 
\hspace{0.2cm}\textit{mg}
\footnotesize De uiduis. 
\normalsize| 
\hspace{0.2cm}\textit{mg}
\footnotesize Acto. 6. 
\normalsize| 
\section*{AD TIM. }
\marginpar{[ p.149 ]}\pstart uero et in loco adiicit Paulus, Quae uerae uiduae sunt hoc est orbatae desolatae et uiri uirium necessariorum et liberorum solatio destitutae. Immiscebant se enim, quod iam diximus, uiduarum turbae quae uerae non e- rant uiduae neque merebantur ecclesiastico subleuari sub sidio. Suppetebant enim aliis et uires et facultates, aliae uel necessarios habebant uel liberos. Atqui opes ecclesiasticae non alendis delitiis fouendoque otio, sed re creandae calamitati, praeterea egenis et pauperibus non opulentioribus et sua ope ualidis inseruire debent. Ideoque subiungit protinus, Quod si qua uidua liberos aut nepotes habet, ea uiduitatis praetextu nequaquam subterfugiet naturae debitum et pietatis officium, quin potius discant omnes uiduae imprimis domum pro priam id est liberos et nepotes ἐυσερωμ, id est nutri re sancte et pie instituere, atque ita uicem rependere maioribus, hoc est ea liberis et nepotibus impendere officia quae ipsis quoque impensa sunt a maioribus, qui si omnem liberorum curam abiecissent ait et otioso ali cui uiuendi generi sese dedidissent, quemadmodum uos in praesentiarum neglectis relictisque liberis ocio- se et secure cupitis ali ecclesiae stipendio, periissetis proculdubio. Cur igitur perdere uultis nunc quos do- minus praecepit seruare. Ad hunc uero modum et D. Ambrosius hanc Pauli sententiam retulit ad uiduas. Theophylactus existimauit referendam esse ad filios  \pend
\section*{COMMENT. IN I. EPIST. }\pstart et nepotes, in hanc nimirum sententiam, Quod si qua uidua est quae orbata quidem marito sed filios ad buc et nepotes habet, hi neutiquam committere de- bent ut illa ad patrocinium confugiat ecclesiae, sed ipsi potius huic ut necessariis et familiaribus reliquis pro piciant. Nam pietas hoc poscit ut non modo liberos et propriam curemus domum sed pietatem quoquam mutuam referamus iis a quibus accepimus primordia uitae. Vtraque sententia non male congruit Paulino in- flituto. Ad utramque sequitur causalis illa sententia, Hoc enim est honestum et acceptum coram deo. Nam si sacras uerses historias inuenies sanctissimos homines educandis liberis et nepotibus et incubuisse studiosissi me et maioribus mutuam pietatem quam officiosissi me retulisse. Notissimum est et Graecorum ἀντιπε- λαργο͂μ, de quo et Eras.Rot. in Chiliadum opere. Vbi ergo nunc sunt monachi et monachae, delicatu- lae, quae neque liberis educandis et pie instituendis inui gilant, neque uicem rependunt maioribus? Abdunt e- nim se in cauernas, imo in Adonidis illos hortos suos in quibus otiose et delicate uiuunt, neque ulla nepotum parentum aut necessariorum cura tanguntur. Vbi nunc sunt quibus infame et parum uidetur esse Chri stianum sed magis diffidentiae argumentum, si quis cu ret rem domesticam? Discant hinc quoq Catabaptistae quid per domum intelligat Paulus, liberos nimirum et  \pend
\section*{AD TIM. CAP. V. }
\marginpar{[ p.150 ]}\pstart nepotes, ne porro a scripturis alienum credant quod dicimus Apostolos baptisasse infantes. Constat enim domos ab istis esse baptisatas. Discant item ex hoc Pauli praecepto qui aerario ecclesiastico praefecti sunt quos recipiant ecclesiae stipendio alendos. Quod enim de uiduis dicitur id et de egenis in uniuersum, quod dispensandi rationem attinet, intelligitur.  \pend
\phantomsection
\addcontentsline{toc}{subsection}{\textit{Porron quae uere uidua est ac de- solata, sperat in Deo et perseuerat in obsecrationibus ac precationi- bus noctu dieque . Porro quae in de- litiis uersatur ea uiuens mortua est. Et haec praecipe ut irreprehensibi- les sint. Quod si quae suis et maxi- me familiaribus non prouidet, fidem abnegauit et est infideli deterior. }}
\subsection*{\textit{Porron quae uere uidua est ac de- solata, sperat in Deo et perseuerat in obsecrationibus ac precationi- bus noctu dieque . Porro quae in de- litiis uersatur ea uiuens mortua est. Et haec praecipe ut irreprehensibi- les sint. Quod si quae suis et maxi- me familiaribus non prouidet, fidem abnegauit et est infideli deterior. }}\pstart Haec ad expositionem eorum pertinent quae pro- posita sunt. Dixerat eas uiduas esse honorandas quae uere sint uiduae, nunc ergo demonstrare incipit quae nam sint uere uiduae. Illa, inquit, uere uidua est quae mariti, liberorum, nepotum, denique omni huius mun- di solatio, deftituta, a solo deo pendet, et in hunc om- nemspem fixit, ideoque non ociosa et delicatula rebus  \pend
\vspace{0.5cm}\noindent
\vspace{0.2cm}\rule{1cm}{0.2pt}\\ 
\hspace{0.2cm}\textit{mg}
\footnotesize Quae nam sint uere ui duae. 
\normalsize| 
\section*{COMMENT. IN I. EPIST. }\pstart intenta est uanis, sed seriis piis et sanctis, hoc ect, quae instar Annae (cuius in2. cap. euangel. histo. meminit Lu cas) noctes atque dies impendit precibus sacris sedula as siduaque est in omnibus officiis pietatis. Hic uero ora- tionis filum nonnihil interrumpens connectit uerae ui duae, fucatae exemplum siue notationem. Quae uero in delitiis uersatur, inquit, eauiuens mortua est. Sunt e- nim uiduae quae pietatem praetexentes nihil quam de- litias otium et licentiam carnis quaerunt, nempe ut hac ratione subterfugiant molestias quas uberrime pro fundit liberorum educatio et maiorum nutritio, ut subducant se a cura domestica et mariti imperio, utque ́ libere et sine ullius magisterio quod lubet uel faciant uel omittant. At tales, inquit apostolus, uiuentes mor- tuae sunt, Neque enim uiuere existimandi sunt homines qui carni non Deo uiuunt, cum extra pietatem non sit uita. De huiusmodi mortuis loquutus in euangelio do minus ad adolescentem quendam petentem ut ad suos redire et funera parentum curare liceret, dixit, Sine ut mortui sepeliant mortuos suos, tu uero sequere me. Sed et apudgentes sunt prouerbia quae plurimum huc faciunt. Sophocles enim prouerbiali figura dixit in Antigone, ut est in in Chiliadibus, oὺ τίθεμε ἐγὼ Σιῷ τοῦτορ, ἀλλἔα τυχομ ἠγοῦμααι νεκρόμ, id est, Hunc ego uiuere  \pend
\vspace{0.5cm}\noindent
\vspace{0.2cm}\rule{1cm}{0.2pt}\\ 
\hspace{0.2cm}\textit{mg}
\footnotesize Viduae fu catae. 
\normalsize| 
\hspace{0.2cm}\textit{mg}
\footnotesize Matth.8. 
\normalsize| 
\section*{AD TIM. CAP. V. }
\marginpar{[ p.151 ]}\pstart Hauld arbitror, uinum at cadauer inalco. Coepit itaque et uiuum cadauer, et uiuum sepulchram de iis dici qui sic uiuunt ut nihil uita dignum agant. Paulus ergo delicatam uiduam iudicauit esse mortuam, id est nihil minus quam uiduam, utpote in qua praeter inane nomen supersit nihih Subiungit his protinus ge- nerale praeceptum, Et haec praecipe ut irreprehensibi les sint, hoc est, Quandoquidem ergo non deerunt qui omnino uiduae et esse et uideri uolent, tu praeci- pies ut ita se gerant in uiduitate ne cui praebeant an- sam male suspicandi. Tandem uero in desertrices illas uiduas ut absterreat a uitae instituto inhonestissimo di stricte sic pronunciat, Quod si qua his praeceptionibu contemptis perrexerit liberorum aut nepotum aut qui alioqui ad eius familiam pertinent curam excute- re, ea merito non mortua tantum censebitur, sed per- fida quoque et infideli deterior. Vel quia hoc officii si bi naturae uinculo iunctis negat quod ethnici nulla le- gum necessitate sed naturae instinctu suis praestant. Vel quod grauius sit peccare eum qui sub lege uiuit, quam qui non subiectus est legi. At hodie summa habetur pi£ etas si quis domi relictis liberis et uxore aut etiam grandaeuis parentibus uel Hierosolymam abeat, uel in monasterium abdat sese sibi uicturus aut fortassis etiam uentri. Vide quo impudentiae uentum sit in ecclesia Chrisui Nec impari prorsus daementia desipiunt hodie.  \pend
\vspace{0.5cm}\noindent
\vspace{0.2cm}\rule{1cm}{0.2pt}\\ 
\hspace{0.2cm}\textit{mg}
\footnotesize Minae in ui duas deser- trices. 
\normalsize| 
\section*{COMMENT. IN I. EPIST. }\pstart et Catabaptistae. Relictis enim uxoribus liberis et fa miliis quibusdam sane non omnino male institutis hacte nus, tanquam attoniti et phanatici homines circum- cursitant atque ita se pomeria regni Christi propagatu ros deiurant, cum interim et rem familiarem conco- quant subuertantque totam, et ecclesiam Christi miris modis obturbent. Caeterum ut uno uerbo rem istam ab soluam, Viri boni honestaeque matrisfamilias est deum timere, familiae prospicere et quenque in sua uocatione manere. Ista qui negligit et curiositatem excolit, non est quod se fidelem esse uel credat uel iactitet, reuera enim impostor et contemptor est diuinae legis.  \pend
\phantomsection
\addcontentsline{toc}{subsection}{\textit{Vidua allegatur non minor an- nis sexaginta, quae fuerit unius ui- ri uxor, in operibus bonis homi- num testimonio comprobata, si fi- lios educauit, si fuit hospitalis, si sanctorum pedes lauit, si afflictis subministrauit, si in omni opere bono fuit assidua. }}
\subsection*{\textit{Vidua allegatur non minor an- nis sexaginta, quae fuerit unius ui- ri uxor, in operibus bonis homi- num testimonio comprobata, si fi- lios educauit, si fuit hospitalis, si sanctorum pedes lauit, si afflictis subministrauit, si in omni opere bono fuit assidua. }}\pstart Semel coeperat ueram describere Viduam, filum orationis interrupit interim, at nunc redit et exactis sime delineat quam uelit suscipi uiduam. Necesse enim  \pend
\vspace{0.5cm}\noindent
\vspace{0.2cm}\rule{1cm}{0.2pt}\\ 
\hspace{0.2cm}\textit{mg}
\footnotesize Anni ul- duae. 
\normalsize| 
\section*{AD TIM. }
\marginpar{[ p.152 ]}\pstart est ut in his recipiendis habeatur delectus, partim ne supra uires oneretur ecclesia, partim ne officium con feratur in indignas. Duo autem requirit inuera et Christiana uidua Paulus, aetatem uidelicet et uitam an teactam. Aetatem inquam grandaeuam uiribusque effoe tam, gignendis educandisque liberis uel incongruam uel diutina exercitatione fessam. Ideoque ne recipia- tur, inquit, uidua minor annis sexaginta. Atque carel haec aetas plerumque omni suspitione impudicitiae. Dein- de uitam in ea requirit pbam, nec illam modo qua nunc uiuit sed qua uixit olim. Siquidem ne huic quidem aetati pu- tat fidendum nisi eam uita anteacta commendet. Iam itaque species siue specimina quaedam purioris uitae re- censet, quae cum manifesta et exposita sint non est quod pluribus explicentur. Si inquit unico fuit mari- to contenta nec cupiditate uicta pluribus se prostituit, si benefactis sibi parauit laudem et honestam existi- mationem apud homines, uirtutesque eius praedicentur ab omnibus, si lib eros sancte in omni pietate educa- uit, si fuit hospitalis, aut si propter fortunae tenuita- tem hospitalitatem exercere nequiuit, sanctorum ta- men pedes lauit, hoc est fidelibus officiis sese submisit et charitate per omnia inseruiuit. Nam ablutio pedum abiectissimum quidem inter officia comprehendit sy- nel docha quadam reliqua et maiora omnia, id quod ui pIoan cap, colliquescit. Rursus si, inquit, egestate  \pend
\vspace{0.5cm}\noindent
\vspace{0.2cm}\rule{1cm}{0.2pt}\\ 
\hspace{0.2cm}\textit{mg}
\footnotesize Vita ante- acta uidux 
\normalsize| 
\section*{COMMENT. IN I. EPIST. }\pstart laborantibus et afflictis de suis subministrauit facul- tatibus, breuiter si nullum opus bonum praetermisit, sed potius in omni iustitia facienda fuit assidua: si in- quam talem inueneris uiduam non dubitabis inopiam eius facultatibus subleuare ecclesiasticis. Hac uero lan ce hac regula metietur expendetque uir ecclesiasticus quos e pauperum turba recipiat in eorum numerum qui ecclesiastico aluntur stipendio. Si uero prexerimus mimos, scurras, scorta, lenas, ambubaias, decoctores, milites, nebulones, et hoc genus omne ecclesiasticis nutrire opibus, sentiemus aliquando potentem illam dei manum quam monachis et sacrificis pontifitiis ab utentibus bonis ecclesiasticis uidimus inflictam gra- uiter. Confutatur item ex praesenti loco Pauli quicquid de anno probationis nugati sunt impurissimi mona- chi. Equidem insigne est humanae infirmitatis argu- mentum, quod apostolus tantus noluit uiduam recipi minorem annis sexaginta. Arbitror profecto hoc apo stoli iudiciumn egregie nobis obseruandum tenendumque ́ ne praecipitanter et inconsulte nimis agentes adul- terii, fornicationis et scelerum nephandissimorum authores fiamus.  \pend
\phantomsection
\addcontentsline{toc}{subsection}{\textit{Porro iuniores uiduas reiice, cum enim lasciuire coeperint aduer- sus Christum, nubere uolunt, ha- }}
\subsection*{\textit{Porro iuniores uiduas reiice, cum enim lasciuire coeperint aduer- sus Christum, nubere uolunt, ha- }}
\section*{AD TIM. CAP. V. }
\marginpar{[ p.153 ]}
\phantomsection
\addcontentsline{toc}{subsection}{\textit{bentes condemnationem quod pri- mam fidem reiecerint, simul autem et otiosae discunt circumire domos, imon non solum otiosae uerume- tiam garrulae et curiosae, loquen- tes quae non oportet. }}
\subsection*{\textit{bentes condemnationem quod pri- mam fidem reiecerint, simul autem et otiosae discunt circumire domos, imon non solum otiosae uerume- tiam garrulae et curiosae, loquen- tes quae non oportet. }}\pstart Adhuc exprimit planioribus et pluribus uerbis quod quidem ex disputatis colligi potuisset sed mino- re cum authoritate. Ergo veωτέρας id est iuniores in quit et minores annis sexaginta uiduas a stipendio uiduarumque collegio procul arcebis ecclesiastico. Atque  hic causas recenset cum uarias tum grauissimas. Prima huiusmodi est, Quae recipiuntur uiduae quasi carni et mundo iam mortuae ueluti Christo nubunt et spem de se praebent nunquam se ad coniugia mundi aspiratu- ras. Caeterum cum uires sint adhuc integrae et anni ad nuptias congrui, quibus iam accessit otium potissi- mum et efficacissimum ad excitandam alendamque li- bidinem incitabulum, fit ut uictae iam cupiditatibus of ficii et instituti sui immemores ferocire lasciuireque ́ (Sρηνῖαμ enim ferocire significat siue lasciuire, quod deductum est a uerbo seρeμ ceu metaphora sumpta a iumentis quae cum pabulo ferociunt, auel- lunt habenas et suopte arbitrio feruntur) incipiant  \pend
\vspace{0.5cm}\noindent
\vspace{0.2cm}\rule{1cm}{0.2pt}\\ 
\hspace{0.2cm}\textit{mg}
\footnotesize De uiduis iunioribus. 
\normalsize| 
\hspace{0.2cm}\textit{mg}
\footnotesize uαrαςρu ιἁσωυσι. 
\normalsize| 
\section*{COMMENT. IN I. EPIST. }\pstart et omnino nubere uelint. Et sane nubere per se qui- dem uitio caret. Honorabile enim est connubium a- pud omnes. Quia uero fidem et Christo sponso et ec- clesiae dederunt semel, et sese sua sponte abdicarunt connubio, hinc lasciuia et nuptiae earum in ignomini am uergunt Christi, id quod Paulus dixit Lasciuire ad uersus Christum. Et prophani homines culpam istam iis non imputant quorum uitio fit, sed religioni omnem et Christo impingunt, clamantes, Haec illa nimirum egregia uiduarum Christianarum castimonia est, haec sacrosancta Christianorum fides, haec denique illa ec- clesia est quae libidini nutrimenta suppeditat, seilicet. Atque boc quidem alterum uitium est quo et Chri- stum et ecclesiam totam denique reli gionem infamant. Ad haeret huic et alterum, nempe quod uiduae sibiipsis perpetuam inurunt ignominiae maculam, dum scilicet fidem pactam soluunt. Paulus enim per κρῖμα id est iudicium, quod interpres noster non male uertit con- demnationem non intellexit animae perditionem sed maculam infamiae ein nadred unno furzugldao man sy diumb vssryot vno bsoillt. Verum at tigimus hunc locum et in Comment. nostris in cap. 23.Act. apost. Et Paraphrastes eodem modo hunc lo- cum explicans, Atque interim hanc infamiae maculam, inquit, sibi contrahunt, quod fidem quam Christo pa cti uidentur fecerint irritam, duplici nomine notandae  \pend
\section*{AD TIM. CAP. V. }
\marginpar{[ p.154 ]}\pstart et quod temere susceperint professionem castitatis non satis exploratis uiribus, et quod a suscepta resi- lierint. Impugnat hic locus Vota et uitam monasti- cam in uniuersum omnem quam fortissime. Idem enim iudicium est de Monachis et Nonnis quod de Viduis. Tam aperta et firma sunt hic omniaut solus hic lo- cus Vota prohibuisset monastica, nisi peccata nostra meruissent aliud. Supersunt et alia incommoda istis haud meliora. Tametsi enim non lasciuiant, inquit, ad uersus Christum neque uiolent pactam fidem, attamen conuersationem uitae ita instituunt ut plane dedecori sint ecclesiae. Otium enim omnium uitiorum suscita- bulum et alimentum docet illas oberrare per aedes a- lienas, atque in his nugari et garrire quae uiduas imo honestas mulieres non decet. Accedit his et curiositas quae et haec et aliain ipsis pariet dedecora. Et curio sitas quidem foedum in omnibus hominibus, praecipue autem in foeminis uitium est. Atqui hisce uitiis omni bus certo ita inuoluentur scio ut uiduitas ipsarum plus ignominiae quam gloriae habitura sit in ecclesia. Quid ergo fiet?  \pend
\phantomsection
\addcontentsline{toc}{subsection}{\textit{Volo igitur iuniores nubere, liberos gignere, domum admini- strare, nullam occasionem dare ad- uersario ut habeat maledicendi cau- }}
\subsection*{\textit{Volo igitur iuniores nubere, liberos gignere, domum admini- strare, nullam occasionem dare ad- uersario ut habeat maledicendi cau- }}
\vspace{0.5cm}\noindent
\vspace{0.2cm}\rule{1cm}{0.2pt}\\ 
\hspace{0.2cm}\textit{mg}
\footnotesize Curiosstas mater est garrulita- tis. 
\normalsize| 
\section*{COMMENT. IN I. EPIST. }
\phantomsection
\addcontentsline{toc}{subsection}{\textit{sam. lam enim nonnullae deflexe- runt sequutae Sathanam. }}
\subsection*{\textit{sam. lam enim nonnullae deflexe- runt sequutae Sathanam. }}\pstart Planissima sunt quae dicit, Cum ergo ait haec ita habeant tutius multo est ut uiduae iuniores nubant maritis quorum authoritate coerceatur ipsarum li- centia, ut gignant liberos et administrent familiam, qua re fiet ut ab otio abstractae et garrulae et curio- sae esse desinant, sed ea potius curent quae familiae sunt. Et haec plane optima est omnium mulierum discipli- na, quae et coelitus ad nos delata et uoce propheta- rum ac apostolorum nobis proposita, sanctissimorum quoque hominum exemplis comprobata est, ut nunc con scientia bona corruptissimam illam disciplinam mo- nasticam diuersum praecipientem contemnamus et damnemus. Addit his aliud quoque praeceptum aposto- lus, Sic uidelicet uitam instituant mulierculae ne ullam dent calumniatoribus maledicendi ansam. In quo bre- ui accinctoque praecepto multa simul monita continen tur. Ne ipsam exquisitior cultus infamet, ne oculorum nutibus et hilaritate uultus uirorum post se greges trahat et conueniat ei uersiculus ille uulgatus, Risit et arguto quiddam promisit ocello: breuiter ne quid uel dictis uel factis aut denique toto corporis habitu de- signet quod in impudentiae aut sinistram inhonestamue alians suspicionem rapi possit. lam cum praeceptum de  \pend
\vspace{0.5cm}\noindent
\vspace{0.2cm}\rule{1cm}{0.2pt}\\ 
\hspace{0.2cm}\textit{mg}
\footnotesize Officia iu- niorum ui duarum. 
\normalsize| 
\section*{AD TIM. CAP. V. }
\marginpar{[ p.155 ]}\pstart reiiciendis iunioribus uiduis matrimonioque subiugan- dis uideri poterat alicui nimis seuerum, sub finem ap- pendit praecepti sui rationem siue defensionem com- prehensam parcissimis hisce uerbis, Iam enim nonnul lae deflexerunt sequutae Sathanam. Argumentatur e- nim ab exemplo et ait, Non temere praecipio iunio- res reiicere uiduas. Iam enim quarundam perfidia do cuit quid ipsis fidamus, nihil plane. Nam semel susce- ptae in curam ecclesiae mox atque otii sensere blandimen ta simul et lasciuire coeperunt, et ruptis pudicitiae uinculis scortari quoque , circumcursitare, curiose age- re, cornicari, calumniari, breuiter sathanam imitari exintegro. Ista porro non fient cum maximo ecclesiae et religionis detrimento, si retineantur in officio et coniugio sancto eaque ratione abarceantur ab otio a hypocrisi et sanctimonia siue castitate ficta. Haec per omnia quadrant et monasticis.  \pend
\phantomsection
\addcontentsline{toc}{subsection}{\textit{Quod si quis fidelis aut si qua fidelis habet uiduas, suppeditet il- lis et non oneretur ecclesia, ut his quae uerae uiduae sunt suppetat. }}
\subsection*{\textit{Quod si quis fidelis aut si qua fidelis habet uiduas, suppeditet il- lis et non oneretur ecclesia, ut his quae uerae uiduae sunt suppetat. }}\pstart Verum opponit hic aliquis, Tantae uero opes quan tae uiduae meae uel ad uitam sustent andam uel maritum aliquem conciliandum sufficerent non suppetunt, ic-  \pend
\section*{COMMENT. IN I. EPIST. }\pstart circo ad ecclesiae cogor subsidium confugere. Respon det apostolus, Plane non est necesse ut egestatis gratia castitatis professionem ulla suscipiat uidua. Debet enim horum inopia cognatorum subleuari beneficentia. Non enim par est ut Christianus aut Christiana cuius pietas debebat etiam alienis opitulari sibi cogna- tam uiduam destituat et alendam obtrudat ecclesiae. Hinc enim sequitur incommodum duplex. Prius one- ratur ecclesia, posterius exhaustis demum facultatibus ecclesiasticis non supererit quo alantur uerae uiduae. At hodie neglectis istis ueluti nunquam uel dictis uel scriptis a Christi domini apostolis audent opulentio- res magnates et nobiles liberos suos ecclesiae alendos obtrudere. Quibus enim numerosa obtingit soboles ii ex omnibus unum aliquem ueluti familiae columen deligunt, huic relinquunt haereditatem et facultates omnes, caeteros uero uel fatuos uel non prorsus qui- dem insipidos sed infoeliciores tamen ueluti spurios abdicant maximeque reluctantes curae committunt ec- clesiasticae. Videmus enim maximam ecclesiasticarum facultatum portionem in huiusmodi hominum pote- statem esse redactam. Atque hi quidem in ecclesiasticis opibus luxuriant hodie. Multi his ueritatem euange- licam renascentem oppugnant. Hi cum magnatum sanguine sint creti proceres et uiros in republica pri marios habent et fautores instituti sui et defensores,  \pend
\section*{AD TIM. CAP. V. }
\marginpar{[ p.156 ]}\pstart mutuasque utrinque sibi conferunt operas, uerum uincet tandem ueritas, restituetque usui genuino ecclesiae suae thesauros Christus dominus. Cui gloria in saecula se- culorum, Amen. Haec uero de Viduis hactenus.  \pend
\phantomsection
\addcontentsline{toc}{subsection}{\textit{Qui bene praesunt presbyteri duplici honore digni habeantur, maxime ii qui laborant in sermo- ne et doctrina. Dicit enim scriptu- ra, Boui trituranti non obligabis os. Et, Dignus est operarius mer- cede sua. }}
\subsection*{\textit{Qui bene praesunt presbyteri duplici honore digni habeantur, maxime ii qui laborant in sermo- ne et doctrina. Dicit enim scriptu- ra, Boui trituranti non obligabis os. Et, Dignus est operarius mer- cede sua. }}\pstart Ad disputationem de Alimonia Viduarum in qua comprehendimus omnem pauperum curam pertinet et tractatio de stipendio episcoporum. Namut paupe res ita presbyteri quoque ecclesiae toti seruientes ex ec- clesiastico alendi sunt aerario. Sed et Christus ipse et apostoli ex munificentia uixerunt sponte offerentium. Et in Corinthiis scripsit Apostolus, An nescitis quod hi qui in sacris operantur ex sacrificio uiuunt? Qui sacrario assistunt una cum sacrario partem accipi- unt? Sic et dominus ordinauit ut qui euangelium an- nunciant de euangelio uiuant. Legimus praeterea in Actis apostolorum uiros cordatiores facultates suas rd pedes apportasse apostolorum, quibus et paupe-  \pend
\vspace{0.5cm}\noindent
\vspace{0.2cm}\rule{1cm}{0.2pt}\\ 
\hspace{0.2cm}\textit{mg}
\footnotesize De faculta tibus eccle ssasticis et usu earum. 
\normalsize| 
\section*{COMMENT. IN I. EPIST. }\pstart peribus et uiris subueniretur ecclesiasticis. Eius rei gratia et hoc exemplum imitantes maiores nostri tam primarii quam plebei uiri facultates suas donarunt ecclesiis. Quo ego pertinere puto regum donationes, nobilium fundationes, decimas, possessionum ususfru- ctus, ipsas possessiones, omnes denique opes ecclesiasti- cas. Has uero prisci illi patres nostri et religionis ue- rae Mecoenates proculdubio studiosis et pauperibus, non otiosis et uentribus desti narunt alendis. Neque du bium est quin ueteres illi sancti et docti episcopi istis sancte sint usi. Certe Gregorius ille magnus episcopus Roma ad Cantuarien. Episcopum August. disertis uer bis, Mos, inquit, sedis est apostolicae ordinatis episco- pis praecepta tradere, ut ex omni stipendio quod acce dit quatuor fieri debeant portiones, una uidelicet e- piscopo et familiae propter hospitalitatem atque susce ptionem, alia clero, tertia pauperibus, quarta eccle- siis reparandis. Intellexit autem per Clerum diacono- rum siue studiosorum collegia. Nam et ueterum coe- nobia aliud non fuerunt quam studiosorum collegia et pauperum xenodochia. Id quod liquet uel ex Au- gustini uita per Posidonium conscripta, in cuius 5.11. 23. et 24. cap. plurima quae huc faciunt inuenies. Ha- ctenus ergo in dispensandis facultatibus ecclesiasticis haud uehementer peccatum est uel a mediae aetatis ac posterioris saeculi episcopis. Tandem uero ubi tempo-  \pend
\vspace{0.5cm}\noindent
\vspace{0.2cm}\rule{1cm}{0.2pt}\\ 
\hspace{0.2cm}\textit{mg}
\footnotesize Monaste- ria uetusta. 
\normalsize| 
\section*{AD TIM. }
\marginpar{[ p.157 ]}\pstart rum iniuria et supina episcoporum postremi saeculi negligentia pro religione obrepsisset superstitio, ui- rique ecclesiastici dominari non ministrare et princi- pes non dispensatores rerum ecclesiasticarum esseuel lent, factum est ut noua religionis facie constituta in a- lium quoque usum conuersae sint opes ecelesiasticae, et nunc cedant non studiosis, pauperibus, ueris fidei san- ctae praedicatoribus et iis quibus destinatae sunt ab ini tio usibus, sed Pontificibus episcopisque hoc est regibus et principibus, monachis et sacrificis hoc est super stitionum ministris, quibus interim ex superstitioso- rum hominum obsequiis haud parum accessit diuitia rum, adeo ut in horum potestate sint hodie cultissima et opima omnium regnorum et nationum loca, Cu- ius quidem rei gratia mira hodie oboritur omnium commixtio rerum, tam sane periculosa ut nisi huic ma turo aliquo subueniatur consilio damnum aliquod ma ximum toti Christiano orbi datura sit aliquando. Au- rum quidem hoc plane Tolosanum est, neque iniustum facile uel usum uel possessorem patitur. Proinde non est quod huic negotio consultum putent principes et magistratus Christiani si ecclesiis direptis omnes facul tates ecclesiasticas conuertant in usus prophanos. Non enim in hoc primitus destinatae sunt a testatoribus piis sed in usus sacros. Habent reges et magistratus tribu- ta et uectigalia sua, habent aeraria publica, habent  \pend
\vspace{0.5cm}\noindent
\vspace{0.2cm}\rule{1cm}{0.2pt}\\ 
\hspace{0.2cm}\textit{mg}
\footnotesize Religio pe perit diuiti as et filia deuorauit matrem. 
\normalsize| 
\section*{COMMENT. IN I. EPIST. }\pstart donationes et reditus suos iis utantur ad res prophas nas et non depraedentur temere ecclesias Christi. Quod si omnino satisfacere suo uolunt officio, aboleant potius impios bonorum ecclesiasticorum abusus, imo si epi- scopi sua sponte iugum Christi recipere nolint reuo- cent ipsi pietatem auitam, restituant ecclesiae religio- nem ueram quantumque conseruandae et alendae huic sufficiat facultates ex omnibus ecclesiasticis opibus praecipue autem e priscis istis donationibus seligant priscumque ueteris ecclesiae reuocent institutum. Pri- sco autem instituto liberaliter et quantum cuique sa- tis erat ipsumque deposcebat decorum prospiciebatur ex aerario ecclesiastico studiis honestioribus et sacris, studiosae iuuentuti, episcopis deinde et familiae ipso- rum, praecipue uero pauperibus, uiduis, peregrinis, ca ptiuis et necessitate in ecclesia laborantibus, quin et ad sarcta tecta portio opum seruabatur certa. Haec in quam omnia et si qua alia ecclesiae necessaria sunt re stituant prius quam ecclesiasticarum opum uel obolum tollant, aut grauissimam supremi numinis ultionem expectent. Caeterum istis quidem dispositis rite si quid supfuerit id nequaquam in priuatos insumatur usus sed cedat magis reipub. uel distribuatur iustis haeredibus, aut seruetur subleuandae calamitati fu urae. Certe si a- nimus fuerit simplex sacer et probus, dispensatio quo que iusta erit. Sacrilegi et auari perperam agent o-  \pend
\section*{158 AD TIM. CAP. V }
\marginpar{[ p.158 ]}\pstart mnia, idque in malum suum maximum. Paulus in prae- senti, Qui bene praesunt presbyteri ait duplici hono- re digni habeantur. Intelligit autem omnes in ecclesia seniores quibus habenae administrationis commissae sunt. Duplex autem honor duplicatum est subsidium, aut quod putauit Theophylactus abundans precium stipendiumue. Hoc uero iis maxime impartiri iubet qui laborant in sermone et doctrina. Neque enim om- nes presbyteri in ecclesia laborant in sermone et do- ctrina. Cum enim uaria sint in ecclesia munia, non unius quoque generis ministri aut presbyteri sunt. Ve rum quotquot ecclesiae inseruiunt suo ministerio, sti- pendio quoq uiuere debent ecclesiastico. Id quod pro nunciatis quibusdam ex lege et euangelio subiectis demonstrat atque confirmat. Prius sumptum est ex Deu. 25.cap. Posterius ex Matth.1o.cap. De utroque annota- uimus quaedam in cap.9. prioris ad Corinth. ut et de stipendio ministris uerbi soluendo disputauimus in s. cap.ad Gal. et in5. cap.1. ad Thessalon. Iterum uero monet honore dignos esse habendos qui bene praesunt et laborant in uerbo et doctrina, nullam mentionem facit cantorum, sacrificorum, aut otiosorum hominum. De his alibi scripsit, Qui non laborat non edat. Quod interim quidam laborant strenue satis praestaret non laborare.  \pend
\section*{COMMENT. IN I. EPIST. }
\phantomsection
\addcontentsline{toc}{subsection}{\textit{Aduersus presbyterum accusa- tionem ne admiseris nisi sub duo- bus aut tribus testibus. Eos ꝗ pec- cant coram omnibus ardue, ut et caeteri timorem habeant. }}
\subsection*{\textit{Aduersus presbyterum accusa- tionem ne admiseris nisi sub duo- bus aut tribus testibus. Eos ꝗ pec- cant coram omnibus ardue, ut et caeteri timorem habeant. }}\pstart Quemadmodum ministro Christi sedulo et probo debetur alimonia ita et defensio, at huiusmodi ne pa- trocinium inueniat aut sperare possit in scelere, sed hic quoque habeat quod metuat. Iccirco apostolus, Ad uersus presbyterum ait accusationem ne admiseris ni- si sub duobus aut tribus testibus. Ratio uero huius le- gis haec est. Presbyter ueritatis est minister. At obse- quium, iuxta Comici sententiam, amicos ueritas odium parit: quocirca metuendum est perpetuo ne quid for san tentent ex odio et inuidia qui accusant presbyte rum quo licentius peccent porro. Deinde nisi integra et sacrosancta sit sacerdotis authoritas profecto per quam uilis erit eiusdem doctrina. Iam si liberum per- mittatur omnibus quoties libet sacerdotis uel uitam uel doctrinam incessere, plane concidit istius authori tas. Vt itaque et authoritas sacerdotis esset illibata et doctrina eius nulla notaretur infamia praecipit aposto lus ne qua aduersus presbyterum accusatio recipiatur nisi eadem legitime duorum aut trium testimonio fue-  \pend
\vspace{0.5cm}\noindent
\vspace{0.2cm}\rule{1cm}{0.2pt}\\ 
\hspace{0.2cm}\textit{mg}
\footnotesize Priuilegia et iudicia e- piscoporum 
\normalsize| 
\hspace{0.2cm}\textit{mg}
\footnotesize Deut.17. 
\normalsize| 
\section*{AD TIM. CAP. V. }
\marginpar{[ p.159 ]}\pstart rit comprobata. Addit huic legi nunc aliam, Quod si accusatio isti intentata rite fuerit comprobata, certe tantum non obstabit illius authoritas aut doctrinae di- gnitas quominus arguatur grauissime, idque coram om- nibus. Causa uero huius legis est, ne sibi ob sacer dotii functionem placeat quisquam et mox sibi pu- tet licere quod libet, uerum cogitet magis quod pu- blico fungatur munere et quod peccata ipsius iam non priuata sint sed publica quodque contagio infici- ant reliquos etiam. Ergo ne episcoporum impunitas pernitiosum exemplum ministret multitudini, quin potius uniuersa ecclesia, uidens ne episcopo quidem impune esse peccare, metuat certoque sibi persuadeat peccata sua non fore impunita, tam districte praecepit apostolus ut presbyterum palam arguat episcopus syn cerior. Nemini autem hae leges nouae uideri debent. Apud omnes enim gentes concessa sunt iusta et aequa, quamuis apud alias et nimia nonnunquam ministris sacrorum priuilegia. His interim adiectae sunt et pri uatae quaedam leges quae acerrime iuberent in impu- ros aut negligentiores animaduertere. Haec ego omnia quae scilicet intra modum aequum et iustum consistunt hodie uellem reuocata a primariis uiris religionis no- strae, quod certo habeam exploratum ne rempubli- cam quidem sine istis foeliciter habituram diu. Immu- nitatem Papisticam et licentiam istam nimiam non probo interim,  \pend
\section*{COMMENT. IN I. EPIST. }
\phantomsection
\addcontentsline{toc}{subsection}{\textit{Obtestor in conspectu dei et domini lesu Christi et electorum angelorum ut haec serues sine prae- cipitatione iudicii, nihil faciens iu- xta propensionem animi, }}
\subsection*{\textit{Obtestor in conspectu dei et domini lesu Christi et electorum angelorum ut haec serues sine prae- cipitatione iudicii, nihil faciens iu- xta propensionem animi, }}\pstart Solent autem ad munia gnauiter et fideliter o- beunda sacramento astringi quibus aliqua in repub. functio obuenit. Ad quem morem apostolus Timothe- um nunc obtestatur et ait, Obtestor te Timothee, per Deum patrem, quo teste et authore geritur haec res, perque Iesum Christum cuius ministri sumus, perque ele ctos angelos dei ministros, et istarum rerum specta- tores ut ea quae tibi praescripsi hactenus praecipue au- tem de iudiciis obeundis serues χωρις προκριματος, id est sine praeiudicio. Est autem praeiudicium cum ex opinione concepta temere quidpiam agitur, et iudi- camus ipsi apud nos antequam compertum habemus, cum magis ex ipsa re cognita colligenda sit sententia. etηδιε μ ττοισομματα πρόσηλισιμ, id est nihil faciens in aliam partem inclinans. Nullus enim hic affectus uo candus est in consilium, sed quod ipsa res poscit quod denique aequum et iustum iubet faciendum est. Haec ue ro memori mente reponenda sunt ut perpetuo oculis obuersentur nostris utque haec primum cogitemus in  \pend
\vspace{0.5cm}\noindent
\vspace{0.2cm}\rule{1cm}{0.2pt}\\ 
\hspace{0.2cm}\textit{mg}
\footnotesize Sancto et integro ani mo fiant o- mnia. 
\normalsize| 
\section*{AD TIM. CAP. V. }
\marginpar{[ p.160 ]}\pstart omnibus nostris functionibus, praecipue iis quibus uel iudicia agitamus uel ministros ecclesiis ordinamus. Se quitur enim quod ad haec pertinet.  \pend
\phantomsection
\addcontentsline{toc}{subsection}{\textit{Manus cito ne cui imponas, ne que communices peccatis alienis. Temetipsum purum serua. }}
\subsection*{\textit{Manus cito ne cui imponas, ne que communices peccatis alienis. Temetipsum purum serua. }}\pstart Manus enim imponere aliud non est quam ecclesiae aliquem praeficere et ordinare. Caeterum hic magno et acri opus est iudicio cauendumque ne quid fiat uel ex praeiudicio uel in hanc uel illam partem. Si enim in dignus aliquis hoc est uel indoctus uel impurus ordina tur, iam scelera quoque impuri huius hominis eum pol- luunt qui inexplorato aut indigno sanctam illam de- legauit functionem. Manuum enim impositio commu nionis signum est. Praeterea uidetur uitiis fauere quis- quis commaculatum uitiis subuehit ad honores. Hic ue£ ro Paulus, Nolo te graues, inquit, alienis peccatis, uo- lo ut purum et illibatum te custodias. Aliena enim peccata sunt quae licet ipsi opere non faciamus sine no bis tamen hoc est sine nostro consilio impulsu admini- culo et suffragio non fiunt, quae quidem patent latissi me. Ex quibus liquido colligi potest quantis se pecca- tis obstrinxerint onerarintque quibus hactenus fuit or- dinandi potestas. Experientia enim elamitat quales or dinauerint. Videant hodie quoque quid agant quibus  \pend
\vspace{0.5cm}\noindent
\vspace{0.2cm}\rule{1cm}{0.2pt}\\ 
\hspace{0.2cm}\textit{mg}
\footnotesize Aliena pec cata. 
\normalsize| 
\section*{COMMENT. IN I. EPIST. }\pstart haec functio delegata est in ecclesia.  \pend
\phantomsection
\addcontentsline{toc}{subsection}{\textit{Ne posthac bibas aquam, sed uino paululo utere propter stoma- chum tuum et crebras tuas infir- mitates. }}
\subsection*{\textit{Ne posthac bibas aquam, sed uino paululo utere propter stoma- chum tuum et crebras tuas infir- mitates. }}\pstart Cohaeret hoc praeceptum cum superioribus quod in his meminisset puritatis, deinde praecepisset ut sibi ab alienis caueret peccatis, nunc ergo praecipit de u- su uini ne se nimia conficeret abstinentia, hoc est ne pu taret se puritati studere, interim uero sese propria pe rimeret culpa. Nam crimen et infamiam homicidii la trociniique subterfugere ac declinare non potest quis- quis nimia inedia aut nimio corporis neglectu se con- ficit et officiis pietatis ante tempus inutilem reddit. Huc refero immoderata studia, uigilias nimias, labo- res ultra uires susceptos, et quaecunque alia huius gene ris sunt. Prudenter enim deus ( quod ait Ambro.) sibi seruiri uult, non ut serui eius nimietate sua debiles fi- ant et postea medicorum suffragia requirant. Ten- perandum est enim ut si fieri potest coeptum obsequi um gradatim prouehatur quam per inconsiderantiam minuatur. Intemperantia ipsam animam inquietam fa£ cit ut cum de infirmitate sollicita est non tantum dedi ta sit diuinis seruitiis. Peccant ergo grauiter quotquot abstinentia nimia et laboribus inexhaustis uires ani-  \pend
\vspace{0.5cm}\noindent
\vspace{0.2cm}\rule{1cm}{0.2pt}\\ 
\hspace{0.2cm}\textit{mg}
\footnotesize Abstinentia episcopum decet sed non nimia. 
\normalsize| 
\section*{AD TIM. CAP. V. }
\marginpar{[ p.161 ]}\pstart mi et corporis exhauriunt itaque praematuro senio con senescunt, adeout cum officiis pietatis ecclesiae debe- bant maxime utiles esse olim sese fecerunt inutiles. Colligimus autem ex hoc loco Timotheum fuisse ab- stemium. Certe Plato de legibus iubet pueros usque ad etatem XVIII. annorum a uino abstinere quod non sit necesse ignem igni addere. Inde uero usque ad XL. modice uti. Post eam aetatem se plusculum inuitare ad uitae taedium leuandum. Itaque Paulus, Paululo inquit ui no utere, idque propter stomachum tuum et crebras tu as infirmitates. Inest enim uino temperate sumpto mi ra medicandi uis. De qua re qui uult legat 2.ca.23. lib.  Pli. Hoc ipsum si sumatur immodice maximorum quoquam causa est morborum, id quod idem ille Plinius testatur lib.14.cap.22. Nostri episcopi nulla praeceptione ut ui num bibant opus habent cum ultro bibant copiosius quam deceat. Freno ergo cohibendi calcaribus nequa- quam exagitandi sunt.  \pend
\phantomsection
\addcontentsline{toc}{subsection}{\textit{Quorundam hominum pecca- ta ante manifesta sunt praeceden- tia ad iudicium, quosdam uero et subsequuntur. Consimiliter et bo- na opera prius manifesta sunt, et ea quae secus habent occultari non possunt. }}
\subsection*{\textit{Quorundam hominum pecca- ta ante manifesta sunt praeceden- tia ad iudicium, quosdam uero et subsequuntur. Consimiliter et bo- na opera prius manifesta sunt, et ea quae secus habent occultari non possunt. }}
\section*{COMMENT. IN I. EPIST. }\pstart Epanaphora est. Redit enim ad id quod praecepe rat ne communicaret peccatis alienis, et ait, Nolim autem existimes omnia tuorum peccata tibi esse impi tanda, sed ea tantum quae manifesta sunt. Verbi gra- tia, Vixit aliquis inter uos adeo perdite ut nemo hacte nus quicquam sibi de hoc homine bonae frugis potuit polliceri, huic uero communicas tu et hunc subuehis ad honores et officia publica, in quibus administran- dis sui per omnia similis nihil facit laude dignum sed omnia corruptissime, istius certe corruptio et pecca- ta merito tibi quoque imputantur, qui non nescieris cui communicaueris et quali imposueris manus. Rursus insinuat se amicitiae tuae alius quispiam et ipse corru- ptus et nepharius, uerum nouit interim improbitatem suam dissimulare ut humano iudicio deprehendi non possit. Tali si communicaueris aut functionem commi seris publicam in qua gerenda boni uiri officio non fun gatur sed erumpens palam factitet quod clam hacte- nus gessit ad occasionem in pectore, nequaquam istius improbitatis particeps eris, cum talem esse nescieris, Facit ergo apostolus duo peccati genera. In primo ge- nere sunt manifesta illa praecedentia ad iudicium. In secundo genere sunt occulta subsequentia iudicium. Per iudicium autem intelligunt alii extremum iudici- um, ego uero hominis iudicium et cognitionem intel ligo, ut praecedentia iudicium ea sint peccata qua  \pend
\vspace{0.5cm}\noindent
\vspace{0.2cm}\rule{1cm}{0.2pt}\\ 
\hspace{0.2cm}\textit{mg}
\footnotesize Alia pecca ta manife, sta, occulta sunt alia. 
\normalsize| 
\section*{AD TIM. CAP. VI. }
\marginpar{[ p.162 ]}\pstart pracesserunt et cognitionem nostram non subter- fugerunt, quippe cum prius sint facta quam nos de homine formauerimus iudicium. Subsequentia er- go erunt quae postea facta sunt aut fiunt, et de qui- bus iudicare prius non licuit, utpote nondum factis. Consimilem esse rationem bonorum operum nemo iam est qui non intelligat. Ista uero si obseruemus fideliter plurimum commodi afferent nobis in omni uita no- stra. Latissime enim patet istarum sanctissimarum le- gum saluberrimus usus. Nos potissima perstringimus tantum, ne quem oneremus disputatione nimis proli- xa. Pii uero lectoris fuerit ista pro rerum, locorum, personarum et temporum ratione copiosius et pru dentius diducere. Commentatoris uice fungor non o- ratoris. Indico non explico omnia. Sententiam uero uer borum Pauli quanta possum diligentia et breuitate expono.  \pend
\endnumbering\beginnumbering\section{CAP. VI.}
\phantomsection
\addcontentsline{toc}{subsection}{\textit{Quicunque sub iugo sunt serui suos dominos omni honore di- gnos ducant, ne nomen dei et do- ctrina male audiat. Qui uero fide- les habent dominos ne contemnant, sed magis inseruiant quod fideles sint ac dilecti qui beneficentiae par- }}
\subsection*{\textit{Quicunque sub iugo sunt serui suos dominos omni honore di- gnos ducant, ne nomen dei et do- ctrina male audiat. Qui uero fide- les habent dominos ne contemnant, sed magis inseruiant quod fideles sint ac dilecti qui beneficentiae par- }}
\section*{COMMENT. IN I. EPIST. }
\phantomsection
\addcontentsline{toc}{subsection}{\textit{ticipes sunt. }}
\subsection*{\textit{ticipes sunt. }}\pstart Praescribit formulam qua utantur episcopi in insti tuendis Seruis informandisque officiis eorum. Loquitur autem de Christianis Seruis non ministris quidem sed seruitutem seruientibus, quale seruorum genus olim cum apud Hebraeos et Gentes reliquas, tum hodie a- pud Turcas frequentissimum est. Hinc uero periphra stica notatione seruitutis ingenium notans significan- ter, Quicunque ait sub iugo serui sunt, id est, Quicunque  libertate amissa colla sua submiserunt alieno imperio, ut non suo sed alieno uiuant arbitrio, obediant domi- nis suis. Sumpta est enim a bubus metaphora quorum libertas iugo cornibus imposito alligatoque opprimi- tur: ita uero et seruitus hominem liberum alieno sub- iicit imperio ut iam non sibi sed domino suo uiuat ab huius nutu et imperio pendeat. Horum ponit genera duo. Alii enim ethnicis parent dominis, fidelibus siue Christianis alii. Iam erant in seruorum turba non pau ci qui crederent Christiana libertate se quoque a Domi norum suorum liberatos esse imperio. Nam ad hune modum argutabantur, Si dominus impius est indignum plane cui seruiat sanguine Christi redemptus. Si uero pius est dominus, iam maxime iniquum est ut a seruo officium petat quem in fide agnoscit esse fratrem. At Paulus longe aliud sentiens praecipit episcopo ut do- ceat seruos obedire dominis suis. Et principio quidem  \pend
\vspace{0.5cm}\noindent
\vspace{0.2cm}\rule{1cm}{0.2pt}\\ 
\hspace{0.2cm}\textit{mg}
\footnotesize De seruis praeceptio. 
\normalsize| 
\hspace{0.2cm}\textit{mg}
\footnotesize Seruotum uenera. 
\normalsize| 
\section*{AD TIM. CAP. VI. }
\marginpar{[ p.163 ]}\pstart loquitur de iis quibus obtigere domini infideles. Hi in quit, τούσἰδείους δεσπόταρ, id est proprios illos dominos suos non ideo contemnant quia infideles sunt, sed omni honore dignos ducant, hoc est ament et re- uereantur et quaecunque ipsis debent officia praestent fidelissime. Addit stimulum, Ne nomen et doctrina dei male audiat. Siquidem non poterit non Christi et nomen et doctrina apud Gentes fieri infamis et inui diosa si conspexerint uos per illam plane mutatos nunc procaciores incommodiores minusque fidos esse factos quam eratis priusquam nomen Christi profiteremini. Nil plane nouum si assidue querantur doctrinam quam uocant nouam seditionis esse seminarium. Oportet i- gitur per omnia ita nos erga ethnicos gerere ut senti- ant per religionem factos esse commodiores iustiores et officiosiores. Atque hac ratione si qua alia pertrahe mus istos in religionem nostram, cum morum asperi- tate et conuitiis illos tam a nobis quam a religione prorsum alienemus. Deinde hoc ipsum praescribit e- tiam Seruis iis qui dominos nacti sunt Christianos. Qui uero fideles id est Christianos habent dominos, inquit, non ideo contemnaut quod fratres facti sunt hoc est quia fidem receperunt et in consortium uene re fidelium. Tamet si enim fides Christiana quod inte- riorem attinet hominem omnes aequet, in Christo enim Iesu non est masculus aut foemina, dominus aut seruus,  \pend
\vspace{0.5cm}\noindent
\vspace{0.2cm}\rule{1cm}{0.2pt}\\ 
\hspace{0.2cm}\textit{mg}
\footnotesize Serui Chri stianorum. 
\normalsize| 
\section*{COMMENT. IN I. EPIST. }\pstart tobilis aut plebeius, sed omnes unum sumus in Chri- sto, attamen quod exteriorem hominem et politiam attinet rerum humanarum, praeest dominus, et subest Jeruus. Vnde alibi quoque praecepit Paulus, Reddite ti- morem et honorem iis quibus debetis, et quilibet in sua uocatione permaneat. Et in praesenti uerbis clarissi mis idem praecipiens, Tantum abest inquit ut domini receptae fidei Christianae gratia aspernandi deserendi ue uideantur ut hoc nomine magis sint obseruandi. At- que huius praecepti sui causam eandem per antistre- phon subiungit ob quam seruis domini Christiani uide bantur negligendi, óτι πις οί ἐισι κα) ἀγαπητοί, οι τῆς ἐνεργεσίας ἀντιλαμβαυόμενοι, id est, Quia fideles sunt et delecti siue contenti et beneficen tiae recompensatores. Quia, inquit, iam non modo do- mini sed etiam fratres id est fideles sunt, ideo magis nunc quam olim obedire illis et inseruire debetis. De inde quia iam non morosi et difficiles sunt amplius sed per fidem facti ἀγαπητoι, qui officia uestra boni consulere norunt et aequis contenti sunt seruitiis. Nam αγαττάω significat et diligo et boni consulo siue contentus sum. His accedit tertium, Quia beneficen- tiae ipsorum participes estis, et nunc quidem liberali- us quam olim. Cum enim fidem Christi receperunt do mini uestri nunc quidem magis benefici per fidem fa- cti officia uestra in se collata munificentius recompen  \pend
\section*{AD TIM. }
\marginpar{[ p.164 ]}\pstart fabunt. Ad hunc modum exposuit haec Pauli uerba et Theophyl, cuius uerba si quis requirat haec sunt, Do- mino fideli magis inseruies, praesertim cum inte sit i- dem beneficus, quippe qui ut te alat et uestem suppe- ditet et bona reliqua impartiatur curis maximis affli ctetur. Haec Theophy. Monent haec quoque dominos of ficii sui. Verum qui plura de hac re cupit legat quae annotaui in s. cap. ad Ephes.Porro si Christiani sunt qui seruorum utuntur opera, profecto Christianorum nomine non exciderunt qui ex pecuniis uel possessio- nibus usumfructum non iniquum percipiunt. Refelli- tur ergo hoc Pauli loco peruersus Catabaptistarum error, qui Christianos esse negant quicunque uel seruo- rum opera utuntur uel ex usufructu non iniquouicti- tant. Fatemur sane perfectiorem esse qui seruis suis manumissis omnes libertate donat, aut qui mutuum dat et inde nihil sperat, uerum nullo modo patietur ueritas ut Christiano priuetur consortio, qui seruorum opera moderate utitur et de elocatis facultatibus fru ctum percipit aequum. De qua re copiosius disputaui in lib. olim contra Catabaptistas aeditis, Certe episco porum est haec contra furaces et seditiosos tradere fideliter ne quis praetexens pietatem ueritatem bla- sphemet.  \pend
\phantomsection
\addcontentsline{toc}{subsection}{\textit{Haecdoce et exhortare. Si quis }}
\subsection*{\textit{Haecdoce et exhortare. Si quis }}
\section*{COMMENT. IN I. EPIST. }
\phantomsection
\addcontentsline{toc}{subsection}{\textit{diuersam sequitur doctrinam et non accedit sanis sermonibus do- mini nostri lesu Christi et ei quae secundum pietatem est doctrinae, is inflatus est, nihil sciens, sed insa- niens circa quaestiones et disputa- tionum pugnas. }}
\subsection*{\textit{diuersam sequitur doctrinam et non accedit sanis sermonibus do- mini nostri lesu Christi et ei quae secundum pietatem est doctrinae, is inflatus est, nihil sciens, sed insa- niens circa quaestiones et disputa- tionum pugnas. }}\pstart Colligit summam rei et infert nune quod potissi- mum hac epistola intendit, ut episcopus ecclesiis sa- nam doctrinam tradat. Haec, inquit, quae praescripsi hic et docui alibi doceto fideliter, ad ea hominum a- nimos exhortato et exerceto. Verum obiicit hic ali- quis, Atqui nolent multi hoc doctrinae genus, uolent a- liud sibi tradi, adeoque ipsimet aliud docebunt. Hic er- go apostolus, Si quis, inquit, diuersam sequitur doctri- nam et c. Perstringit autem ueritatis et simplicita- tis Christianae aduersarios ac corruptores, docens ipsos omni modo fugiendos esse. Duo uero in his praecipue accusat, studium conflictandi et cumulandi opum si- ue quaestus. Quae quidem episcoporum praecipue pe- stes sunt, quamobrem a Paulo in praesenti suis plane coloribus distinguuntur. Si quis inquit, ἐτεροδιιδα- σκαλe͂, id est uel diuersa docet uel diuersam seqtur doctrinam, et non accedit sanis sermonibus donini  \pend
\vspace{0.5cm}\noindent
\vspace{0.2cm}\rule{1cm}{0.2pt}\\ 
\hspace{0.2cm}\textit{mg}
\footnotesize In disputa- tores et de ctrinae di- uersae do- ctores. 
\normalsize| 
\section*{AD TIM. CAP. VI. }
\marginpar{[ p.165 ]}\pstart nostri lesu Christi, hoc est ei quae secundum pietatem est doctrinae, fidem nimirum et charitatem tradenti, is inflatus est, nihil sciens, et iccirco negligendus. Nam 1.ad Cor.8.cap. idem, Scientia inflat inquit, charitas aedificat. Quod si quis uidetur sibi aliquid scire, non- dum ita doctus est ut doctum esse oportebat. At qui Deum diligit hunc docuit atque eruditum reddidit de- us. Et ad Rom.12.cap. Dico inquit per gratiam quae data est mihi cuilibet uersanti inter uos ne quis arro- ganter de se sentiat sed ita sentiat ut modestus et so- brius sit ut cuique deus partitus est mensuram fidei. Pro inde uera et iusta scientia intra metas fidei et chari- tatis manens humilis et sobria est. Vbi uero superbia et arrogantia est, ubi fidei et charitatis regulae ne- gliguntur, ibi nihil solidae est sapientiae utut uerba ap£ paratus denique totus uitae habitus splendeat. Recte er- go Paulus dixit, Qui non accedit sanis sermonibus Christi is inflatus est nihil sciens, ἀ λα νοσῶμ sed ae- grotans circa quaestiones uα) λoγομαχιας, et uer borum siue disputationum pugnas. Id quod mire in eos quadrat qui omnibus saeculis pura et simplici reli cta ueritate de rebus nihili disputare coeperunt. Cuius quidem generis sunt Valentinianorum apud Iraeneum et Tertullianum deliria, et caeterae haereticorum di sputationes prodigiosae. Nostro saeculo quaestionum la birynthis et uerborum pugnis inuoluerunt religio-  \pend
\section*{COMMENT. IN I. EPIST. }\pstart nem circa uerba digladiantes ipsamque rem negligentes, Scholasticae Theologiae authores, ut parum in his liceat inuenire simplicitatis et puritatis apostolicae, quae qui dem ut simplicissima est ita abhorret ab omni uerbo- rum pugna et quaestionum tricis. Docet autem quam simplicissime fidem innocentiam et charitatem. Ergo de pietate uera in hac magna est cura, de uerbis aut quaestionibus nulla concertatio. Vera enim religio non poterit non negligi cum toti ad disputandum et rixan dum conuertimur. Caeterum hie non omnes damna- mus disputationes sed eas tantum quae huiusmodi sunt quales iam describit apostolus.  \pend
\phantomsection
\addcontentsline{toc}{subsection}{\textit{Ex quibus nascitur inuidia, con- tentio, maledicentia, suspitiones malae, superuacaneae conflictatio- nes hominum mente corruptorum, et quibus adempta est ueritas, qui existimant quaestum esse pietatem. Se iungere ab his. }}
\subsection*{\textit{Ex quibus nascitur inuidia, con- tentio, maledicentia, suspitiones malae, superuacaneae conflictatio- nes hominum mente corruptorum, et quibus adempta est ueritas, qui existimant quaestum esse pietatem. Se iungere ab his. }}\pstart Alterum hoc est malum ex duobus quibus diximus episcopos ceu peste infectos interire, studium nempe conflictandi siue disputandi. Illud iam uiribus et fru- ctibus suis enumeratis graphice depingit. Ex huiusmo di rixosis inquit concertationibus nascitur inuidia, pri  \pend
\vspace{0.5cm}\noindent
\vspace{0.2cm}\rule{1cm}{0.2pt}\\ 
\hspace{0.2cm}\textit{mg}
\footnotesize Disputatio num et quae stionum fru- ctus et uti- litas. 
\normalsize| 
\section*{AD TIM. CAP. VI. }
\marginpar{[ p.166 ]}\pstart mum dum ex aliorum authoritate eleuata nobis para mus authoritatem, nec gloriam dei aut ueritatem sed propriam inquirimus laudem et popularium applau sum. Deinde nascuntur contentiones dum gliscente ri£ xarum calore nullus alii uult concedere. Hinc ubi res exiit in rabiem furor disputantium uberrime profun- dit exquisitissima conuitia, et nonnunquam a uerbis ad uerbera descenditur. Praeterea nascuntur et im- piae suspitiones in Deum et omnem pietatem dum nul lum facimus inquirendi dubitandi et ratiocinandi fi nem, dum ea quae simplicissima alias sunt uocamus in quaestionem. Quo mihi egregia certe Tertulliani uer- ba pertinere uidentur quae ad hunc modum leguntur in libello de praescription. haereticorum, Quid Athe- nis et Hierosolymis? Quid Academiae et Ecclesiae? Quid haereticis et Christianis? Nostra institutio de porticu Solomonis est, qui et ipse tradiderat dominum in simplicitate cordis esse quaerendum. Viderint qui Stoicum et Platonicum et Dialecticum Christianis- mo praetulerunt. Nobis curiositate opus non est post Christum Iesum, nec inquisitione post Euangeliam. Cum credimus nihil desyderamus ultra credere. Et paulo post, At enim si quod debui credere credidi, et aliud denuo puto requirendum, spero utique aliud esse inueniendum, nullo modo speraturus istud nisi quia aut non credideram qui uidebar credidisse aut desii  \pend
\section*{COMMENT. IN I. EPIST. }\pstart eredidisse. Ita fidem meam deserens negator inuenior. Semel dixerim, Nemo quaerit nisi qui aut non habuit aut perdidit. Haec Tertull. Nascuntur item ex disputa- tionibus atque quaestionibus illis παραδ'ίατριssαι. Di- cuntur autem authore Eras. congressus et disputatio nes philosophorum Diatribae Graecis, quas hic addita praepositione inutiles facit et nihil ad rem pertinen- tes. Addit idem Graecanica scholia admonere metapho ram sumptam a scabiosis qui mutuo confrictu conta- ctuque morbum communicant, ut sensus sit hoc rixo- sum genus hominum scabiem suam multis ueluti con- tagio affricare, ut et ipsi contendere, male suspicari, maledicere et inuidere incipiant. Et hactenus in ipsam rem strinxit ealamum, iam conuertitur in per- sonas, Iam qui hisce rebus student homines, ait, sunt mente corrupti, id est quibus ratio turpibus obtene- brata affectibus, quod uerum est amplius non uidet, sed gloriam sectatur et quaestum. Sequitur enim qui- bus adempta est ueritas et qui existimant quaestum es se pietatem, hoc est qui utilitate metiuntur religio- nem. Non enim spectant quantum Christi gloriae et publicae accedat utilitati, sed quantum ipsis propter re£ ligionem numeretur. His nunc infert quod intendit, Se iungere ab his. Nam et sapiens praecepit, Contra uer bosum noli contendere uerbis. Et Paulus ad Titum, Artrcuunt Oouuuim post unam atque alteram admoni  \pend
\section*{AD TIM. CAP. VI. }
\marginpar{[ p.167 ]}\pstart tionem deuita. Pertinax est hypocrisis et inscitia. Proinde pietas uerum constanter et simpliciter pro- ponit et asserit, non indecore contendit. Et multas qui dem huiusmodi disputationes moderatas et pias lege re est apud priscos ecclesiae doctores, quorum alii in- terim sui plane obliti nimium indulsere affectibus. Scho lasticae autem disputationes et quaedam nostri saeculi acerbae concertationes parum sane aut nihil plane laudis obtinebunt apud posteros, quin et nuper ab il lis abhorrere coeperunt quicunque pietati euangelicae operam hactenus non infoeliciter nauarunt. Quor- sum uero attinet, charissimi fratres, nouo quodam do cendi genere pietatem euangelicam propagare uelle? Si tenet nos nominis Christi gloria et religio, si hanc per totum terrae orbem propagare cupimus, cur non potius iis consiliis in ea re perficienda utimur qua ui- demus usos domini apostolos, qui nullis contentioni- bus disputationibus aut curiosis quaestionibus inuolue runt mysteria euangelica sed pura potius et simplici ueritate euoluerunt, docentes ea quae pietatis sunt pure et simpliciter. Quod si cui omnino collubitum est ri- xari disputare et quaestiones nectere, nectat sane et rixetur sed deuitetur interim a fidelibus. Est enim ho- mo mente corruptus et a ueritate alienus ꝗ religionem non uero iusto aequo et bono, sed utilitate metitur. Si enim gloriam dei et pictatis profectum quaerit cur sa  \pend
\section*{COMMENT. IN I. EPIST. }\pstart nis sermonibus Iesu Christi non obedit?  \pend
\phantomsection
\addcontentsline{toc}{subsection}{\textit{Est autem quaestus magnus pi- etas cum animo sua sorte conten- to. Nihil enim intulimus in mun- dum, uidelicet nec efferre quicquam possumus, sed habentes alimenta et quibus tegamur his contenti e- rimus. }}
\subsection*{\textit{Est autem quaestus magnus pi- etas cum animo sua sorte conten- to. Nihil enim intulimus in mun- dum, uidelicet nec efferre quicquam possumus, sed habentes alimenta et quibus tegamur his contenti e- rimus. }}\pstart Descendit iam ad aliam episcoporum pestem, studi- um uidelicet pecuniae siue quaestum, quem paulo post accusabit grauiter, in praesenti ἀυτάρκσαμ laudat ut a lucri studio abstrahat melius. Alii, inquit, quibus ue- ritas et religio uenalis est summam foelicitatem exi- stimant si opibus honoribus et uoluptatibus affluant, uerum consensus omnium sapientum credit quaestum uere maximum et beatitudinem in terris absolutissi- mam esse si quis sit pius, deinde sua sorte praesenti- busque rebus contentus. Hinc enim disputatum est a- pud philosophos solum sapientem esse diuitem, eo quod nihil concupiscat. Qui enim concupiscit adhuc, nondum satis habet: qui uero satis non habet, reuera di ues non est. Huc pertinet Psal.111.cuius initium, Beatus uir qui timet dominum etc. Et Seneca in epistola ad Lucilium egregia de hac re disputauit praecipue epi-  \pend
\vspace{0.5cm}\noindent
\vspace{0.2cm}\rule{1cm}{0.2pt}\\ 
\hspace{0.2cm}\textit{mg}
\footnotesize Sua sorte contentum et pium esse summus quaestus est 
\normalsize| 
\section*{AD TIM. }
\marginpar{[ p.168 ]}\pstart stola 12o. Bias quoque apud Ausonium breuiter et sa- pienter hoc ipsum docuit quod hic apostolus tradit, Ait enim, Quae nam summa boni? mens quae sibi con- scia recti. Quis diues? qui nihil cupiat. Quis pauper? auarus. Congruit his etiam primus ille Horatii Sermo. liber mire suggillans cupidos illos et nunquam sua sorte contentos. Instabilitatem item et incertitudinem fluxarum opum notat apostolus cum subiicit, Nihil in tulimus in mundum uidelicet nec efferre quicquam possumus, Cur ergo diuitiis ut propriis inhaeremus, cum nostris facultatibus istas neque parauerimus neque  nobiscum intulerimus in mundum, quas denique aliis si mul atque possidere incipimus relinquere cogimur. Nul lus igitur ne usus opum, dicat aliquis? Respondet, Opum usus est ut hominibus suppetant quo uestiantur et a- lantur. Haec ubi contigerunt homini non est quod frau datis reliquis sibi eumulet, aut unus omnia occupare cogitet. Pulchre hunc locum exponens Theoph. De- finit tandem, inquit, quid ista sufficientia sit, aitque ; suffi£ cientiam esse tantum rerum potiri quantum nobis ad uictum sat est non ad luxum, et talibus operiri qui- bus tegatur non ornetur hoc corpus, sitque in uestem quodcunque offertur. Prudentius quoque doctissime haec Pauli uerba carmine reddens, canit.  \pend\pstart Summa quies nil uelle super quam postulet usus Debitus, ut simplex alimonia uestis et una  \pend
\section*{COMMENT. IN I. EPIST. }\pstart Infirmos tegat ac recreet mediocriter artus Expletumque modum naturae non trahat extra.  \pend
\phantomsection
\addcontentsline{toc}{subsection}{\textit{Caeterum qui uolunt ditescere incidunt in tentationem et laque- um, et cupiditates multas, stultas ac noxias, quae demergunt homi- nes in exitium et interitum. Siqui- dem radix omnium malorum est studium pecuniae, quam quidam dum appetunt aberrauerunt a fi- de et seipsos implicuerunt dolori- bus multis. }}
\subsection*{\textit{Caeterum qui uolunt ditescere incidunt in tentationem et laque- um, et cupiditates multas, stultas ac noxias, quae demergunt homi- nes in exitium et interitum. Siqui- dem radix omnium malorum est studium pecuniae, quam quidam dum appetunt aberrauerunt a fi- de et seipsos implicuerunt dolori- bus multis. }}\pstart Exponit quantis cum periculis coniunctum sit stu- dium pecuniae et quantum malorum auaritia episco porum pestis maxima post se trahat. Qui ditescere uo lunt, inquit, non qui diuites sunt. Fuerunt enim sanctis simi uiri diuites, nullus autem sanctorum ditescere uo luit. Nam ditescere uolunt quotquot inuita quod dici solet Minerua per phas et nephas opes colligere desti narunt. Qui ergo semel apud sese statuerunt diuitias corradere incidunt in tentationem et laqueum, hoc est in cupiditates multas easque stultas et noxias. Tan- demenim onerati deprimuntur in exitium et interi-  \pend
\vspace{0.5cm}\noindent
\vspace{0.2cm}\rule{1cm}{0.2pt}\\ 
\hspace{0.2cm}\textit{mg}
\footnotesize Pestis ma- xima est a- uaritia. 
\normalsize| 
\section*{AD TIM. CAP. VI. }
\marginpar{[ p.169 ]}\pstart tum ut neque praesentibus bonis neque futuris fruantur. Id quod Luc.12, et 16. cap. apertissimis paradigmatis Pposuit dominus. Aristonymus quoque in historia Auari uitam pollincturae id est coenae mortui comparat ubi cum adsint caetera, laetitia deest. Quid uero uanius et stultius diuitum consiliis uel cogitari potest? Habent in his primas et postremas ambitio luxus et super- bia. Atqui istis polluitur ineptiis mens ad contempla- tionem rerum coelestium et aeternarum condita, de- peritque adeó homo totus tam corpore quam animo. Se quitur enim Sententia quae haec confirmet et quam cum Bione sophista communem habet apostolus, Aua ritiam esse radicem omnium malorum. Cui per omnia congruit quod Graecorum quidam dixere τump φiλαρ γυρίαμ μητρότολιρ ἐιναε πἀσησκακίασ, id est, auaritiam esse totius uitiositatis metropolim. Vnde e- tiam Cicero Officio. lib.2. Nullum, inquit, est uitium tetrius quam auaritia. Et Philostratus, Non sunt tracta biles diuitiae, inquit, debentque inter ardua et difficillima quaeque comnumerari. Qui enim diuitiis inebriantur sen sum ac rationem amittunt affectumque sequuntur impel lentem ad inhonestissima quaeque . Vere enim Prudentius in Psychomachia canit,  \pend\pstart Neque enim est uiolentius ullum Terrarum uitium quod tantis cladibus aeuum Mundani inuoluat populi damnetque gehennae.  \pend
\section*{COMMENT. IN I. EPIST. }\pstart Cura, fames, metus, anxietas, periuria, pallor, Corruptela, dolus, commenta, insomnia, sordes Eumenides furiae, monstri comitatus aguntur. Buius generis et alia quaedam ex Sallustio annotaui in 2.cap.1.ad Thess. Sequitur ad haec, Quam quidam dum appetunt ait aberrauerunt a fide. Alibi enim auari- tiam appellat idololatriam. Et sane fides scopum pro- ponit Deum, studium pecuniae cupiditatem. At domi- nus dicit, Vbi thesaurus tuus ibi et cor tuum. Corrum puntur item animi auaritia soluiturque auro fides. Ho- die etiam a uera fide longissime aberrant quidam qui pluris faciunt Mamona quam Christum. Interim dum studiose sectantur opes ut has insumant uoluptatibus et largitionibus ad honores penetrent, multis sese im plicant doloribus, partim quod non sine magnis mo- lestiis facultates sibi parant, partim uero quod illas maiori cum animi cruciatu custodiant, denique maxima cum poenitudine tandem amittant, ut nunc non refe- ram honorum pericula et uoluptatum discrimina.  \pend
\phantomsection
\addcontentsline{toc}{subsection}{\textit{Tu uero homo dei ista fuge, se- ctare autem iustitiam, pietatem, fi- dem, charitatem, patientiam, man- suetudinem. Certa bonum certa- men fidei, apprehende uitam aeter- }}
\subsection*{\textit{Tu uero homo dei ista fuge, se- ctare autem iustitiam, pietatem, fi- dem, charitatem, patientiam, man- suetudinem. Certa bonum certa- men fidei, apprehende uitam aeter- }}
\section*{AD TIM. }
\marginpar{[ p.170 ]}
\phantomsection
\addcontentsline{toc}{subsection}{\textit{nam, ad quam et uocatus es et pro- fessus es bonam professionem co- ram multis testibus. }}
\subsection*{\textit{nam, ad quam et uocatus es et pro- fessus es bonam professionem co- ram multis testibus. }}\pstart Conuersus iam ad Timotheum docet quae hic dece ant episcopum imo Christianum quemlibet, ut scilicet auaritiam fugiat et quae ex ipsa nascuntur, cogitet autem sese totum Deo esse consecratum non cupidita ti et Pluto. Tu uero, inquit, qui Deo consecratus es haec indigna tua professione fugito, ueras opes sequens iustitiam, pietatem, fidem, charitatem et c. In his enim qui diues est uere foelix ueras opes assequutus est. Qui uero hisce caret, opes uel amplissimas possidet, onus habet non opes. Iustitia unicuique suum reddit, ne minem fraudat, neminem afficit iniuria. Pietas assidu- um est dei studium cultus et religio uera. Fides niti- tur bonis aeternis, aspernatur peritura. Charitas mise- ricors est, condonat et libenter communicat omni- busque inseruit fideliter. Patientia spe immortalitatis et amore iusti perdurat in aduersis, fert aequo animo fa- mem, nuditatem et omnium rerum inopiam. Nouit enim meliorem esse pauperem iustum quam diuitem iniquum. Mansuetudo animi tranquillitas est quae iram, uindictam, rixas et perturbatiores animi affectus con pescit. Haec uero decent episcopum Christianum, hi- sce assequendis atque cumulandis opibus inuigilet, istas  \pend
\vspace{0.5cm}\noindent
\vspace{0.2cm}\rule{1cm}{0.2pt}\\ 
\hspace{0.2cm}\textit{mg}
\footnotesize Quid dece- at Christia num. 
\normalsize| 
\section*{COMMENT. IN I. EPIST. }\pstart doceat ueras esse opes. Nam quanto animus corpore praestantior est tanto excellentiora sunt animi bona. Certa bonum certamen fidei. Hoc autem opus hic la- bor est. Conanti enim ueras animi opes assequi con- fertim occursitant hostes exercitatissimi, cupiditates uidelicet creberrimae quibus infracto et audaci ani- mo resistendum fideque est repugnandum. Atque in hoc conflictu subinde in mentem reuocandum ad quid nos uocarit dominus, non ad uoluptates et opes parandas amplas, sed ad aeterna consequenda gaudia. Difficile autem est iuxta uerbum domini diuitem ingredi in re gnum dei. Ista uero ut quam strenue et diligenter fa- cias inuitat te tua professio quam coram multis testi- bus professus es. Sic enim ab anteactis et pristino a- nimi rerumque agendarum ardore ad maiora incitat. Quasi dicat, Hactenus magna constantia ueritatem praedicasti sed et uita pia fidem tuam attestatus es a- deo ut multi sint qui praedicent uirtutes tuas, perges itaque colophonem rebus optime coeptis imponere. A- lii exponunt de professione baptismatis, qua sese ob- uinxerit Christo, mundo autem et pompis eius re- nunciarit. In qua sententia uideo esse et Ambrosium cuius uerba haec sunt, Fidei uictoria est cum omnia ui tia et crimina subiugantur, ut ad aeternae uitae praemia ueniatur, cuius confessio inter ipsa rudimenta fidei te ste interrogante et respondente monimentis ecclesia  \pend
\section*{AD TIM. }
\marginpar{[ p.171 ]}\pstart flicis continetur. Et sane non raro trahuntur a sacra- mentis Christianorum argumenta in exhortationibus.  \pend
\phantomsection
\addcontentsline{toc}{subsection}{\textit{Praecipio tibi in conspectu dei qui uiuificat omnia et lesu Christi qui testatam fecit sub Pontio Pila- to bonam professionem, ut serues praeceptum immaculatus irrepre- hensibilis, usque ad apparitionem domini nostri lesu Christi, quam temporibus suis ostensurus est il- le beatus et solus princeps, rex re- gnantium et dominus dominan- tium, qui solus habet immortalita- tem, lucem habitans inaccessam, quem uidit nemo hominum, neque  uidere potest, cui honor et imperi- um sempiternum. Amen. }}
\subsection*{\textit{Praecipio tibi in conspectu dei qui uiuificat omnia et lesu Christi qui testatam fecit sub Pontio Pila- to bonam professionem, ut serues praeceptum immaculatus irrepre- hensibilis, usque ad apparitionem domini nostri lesu Christi, quam temporibus suis ostensurus est il- le beatus et solus princeps, rex re- gnantium et dominus dominan- tium, qui solus habet immortalita- tem, lucem habitans inaccessam, quem uidit nemo hominum, neque  uidere potest, cui honor et imperi- um sempiternum. Amen. }}\pstart Astringit his quoque ueluti sacramento fidelem dei ministrum quemamodum in cap.5.ut praescriptas prae ceptiones et institutiones fideliter custodiat et offi- cio suo in omnibus satisfaciat. Caeterum quo obstrin- geret episcopum sacramento religiosiore hunc ueluti  \pend
\vspace{0.5cm}\noindent
\vspace{0.2cm}\rule{1cm}{0.2pt}\\ 
\hspace{0.2cm}\textit{mg}
\footnotesize Sacramen- tum episcos porum. 
\normalsize| 
\section*{COMMENT. IN I. EPIST. }\pstart in conspectum cuncta uiuificantis dei producit, eiusque ́ pfessionis meminit quam dominus noster fecit sub Pon tio Pilato, hoc est commonet ipsum passionis et mor tis dominieae, quam interim prodere uidetur quisquis haec instituta uel negligit uel relinquit. His addit spem et amplissimum sanctorum praemium. In apparitio- ne enim domini nostri lesu Christt rependet dominus unicuique secundum facta sua. Per apparitionem uero domini intelligimus uel fatalem cuiusuis hominis di- em, uel iudicium illud extremum quo omnis caro iu- dicabitur. Hie interim obiter describit gloriam et maiestatem aeterni numinis, in hoc ut sibi metueret et principi tanto adhaereret fidelius, et de auxilio et pollicitationibus eius minus dubitaret quisquis ei ser- uit. Porro non est quod hic quis a me prolixiorem expectet disputationem de eo quomodo pater solus beatus sit et immortalitatem habeat, aut quomodo lucemis habitet inacessam inuisibilis, quem nemo ui- derit hominum imo ne uidere quidem possit. Ista e- nim doctissime explicata sunt a sanctis patribus Au- gust. imprimis in lib. de Trinit. dei, item ab Ambro- sio contra Arrianos, et a Tertull. contra Praxeam. Id uero huc maxime mihi pertinere uidetur quod in haec Pauli uerba scripsit Ambrosius, Non sollicitus, inquit, de cura Timothei tam circumspectus est, sed propter successores eius ut exemplo Timothei eccle-  \pend
\section*{AD TIM. }
\marginpar{[ p.172 ]}\pstart siae ordinationem custodirent ipsi quoque futuris for- mam tradentes a semetipsis inciperent. Iam cum ne mi nimum quidem ex his praeceptionibus omnibus custo- ditum sit nil mirum si hodie conflictetur pericliteturque ́ ecclesia.  \pend
\phantomsection
\addcontentsline{toc}{subsection}{\textit{His qui diuites sunt in praesen- senti saeculo praecipe, ne elato sint animo, neque spem ponant in diui- tiis incertis, sed in deo uiuente, qui praebet nobis omnia affatim ad fruendum, ut bene faciant, ut diui- tes sint operibus bonis, ut faciles sint ad impartiendum, libenter con- municantes, recondentes sibiipsis fundamentum bonum in posterum ut apprehendant aeternam uitam. }}
\subsection*{\textit{His qui diuites sunt in praesen- senti saeculo praecipe, ne elato sint animo, neque spem ponant in diui- tiis incertis, sed in deo uiuente, qui praebet nobis omnia affatim ad fruendum, ut bene faciant, ut diui- tes sint operibus bonis, ut faciles sint ad impartiendum, libenter con- municantes, recondentes sibiipsis fundamentum bonum in posterum ut apprehendant aeternam uitam. }}\pstart Multis sane exposuit hactenus quanto cum pericu lo coniunctum sit uelle ditescere, iam ne quis hinc col ligeret diuitem esse aut uti diuitiis licere nemini, prae- scribit diuitibus Christianis quis usus opum et quae sint diuitum officia, adeoque ipsi episcopo leges fert quas demum tradat opulentioribus, Plurima uero praece- pta dat. Primum, Ne diuites sint elato animo. Nam iu  \pend
\vspace{0.5cm}\noindent
\vspace{0.2cm}\rule{1cm}{0.2pt}\\ 
\hspace{0.2cm}\textit{mg}
\footnotesize De diuiti- bus et offi ciis corum. 
\normalsize| 
\section*{COMMENT. IN I. EPIST. }\pstart xta Diuerbium, Satietas parit ferociam: siue quod So- lon dixisse fertur, Satietas nascitur ex opulentia, ex sa tietate ferocia. Hinc uero regum sapientissimus So- lomon orauit dominum diuitias ne daret sibi, ne forte exaturatus dicat quis est dominus? Secundum praece- ptum est, Ne spem ponant in diuitiis. Cui non absimile est illud Dauidis, Diuitiae si affluant nolite cor appo- nere. Sunt enim incertae et periturae. Deus aeternus et Schadai siue Saturnus est qui exaturat implet uiuifi- cat fouet et alit omnia. Huius uirtus et beneficentia affatim suppeditat omnibus animantibus unde uiuant. Id quod Dauid pluribus exponit Psal.103.Huc perti- net bona pars s. capit. apud Matth. et 12. apud Luc. Qui uero neglecto deo uiuo, mamonae adhaerent, alibi ab apostolo idololatrae appellantur. Tertium praece- ptum est, Vt sint benefici. In hoc enim conferuntur a deo opes ut iuste et largiter his utamur aliisque dispen semus non autem ut nos omnia occupemus soli. Hinc prouerbio recte iactatum et in sententias sapientum est relatum, Auarum nisi cum moritur nil recte face- re. Exempla quoque diuitum sanctorum patrum docent beneficentiam. Et diues epulo apud Lucam non ob di uitias damnatur amplas, sed ob ferociam et auaritiam uitae quoque luxum et quod pauperes negligeret nec esset beneficus. Quod si obiicit diues aliquis, Verum euacuantur scrinia mea Christiana illa beneficentia  \pend
\vspace{0.5cm}\noindent
\vspace{0.2cm}\rule{1cm}{0.2pt}\\ 
\hspace{0.2cm}\textit{mg}
\footnotesize Gene.17. Exod.34. 
\normalsize| 
\section*{AD TIM. }
\marginpar{[ p.173 ]}\pstart plurimumque decedit censibus. Hic Paulus, Diuites sint, inquit, bonis operibus. Nam bona opera demum sunt uerae opes ut nunc animi dotibus tanto plus accedat quanto magis decedit praediis, censibus, scrinio. Quar tum praeceptum est, Vt sint faciles ad impartiendum, hoc est ut libenter communicent. Solet enim diuitum uulgus grauatim impartiri si quid impartitur. At de- us hilarem et uolentem datorem requirit. Et mune- rum, iuxta prouerbium, animus optimus est. Sequitur ultimo loco harum praeceptionum fructus, Si inquit di uites non sint elato animo, non fidant incertis diuitiis, non cumulent auare opes, sed benefici potius sint et libenter impartiantur, reponent sibi thesaurum in coe lis, quo perfruantur ubi hinc fuerit excessum: in aeter na enim uita amplissimas reperient opes repositas, ibi inuenient magnifice extructa palatia, in quibus aeter- num gaudebunt, id quod apostolus occulta proculdu- bio antithesi desyderiis et consiliis diuitum huius sae culi opposuit, qui toti in hoc sunt ut extruant aedes am plas et congerant unde aliquando splendide uiuant. Verum citius quam uellent efferuntur ex aedibus suis et quod magno studio diu congesserunt cedit plerumque  quibus maxime noluerunt, id quod carmine uenusto cecinit poëta,  \pend\pstart Diuitias animi solas ego iudico ueras Qui rebus pluris se facit ipse suis  \pend
\section*{COMMENT. IN I. EPIST. }\pstart Hunc adeo ditem hunc opulentum rite uocamus Magnarum quis sit qui uidet usus opum Calculus at si quem misere numerandus adurit Qui misere semper diuitias cumulet Hic ut apes paruo crebroque foramine fosso Sudat in alueolo mel alii comedunt.  \pend\pstart Hactenus uero de iusto opum usu et quanta humans generis pestis sit auaritia.  \pend
\phantomsection
\addcontentsline{toc}{subsection}{\textit{O Timothee depositum serua, deuitans prophanas uocum inani- tates et oppositiones falso nomi- natae scientiae, quam nonnulli pro- fitentes circa fidem aberrauerunt. Gratia sit tecum. Amen. }}
\subsection*{\textit{O Timothee depositum serua, deuitans prophanas uocum inani- tates et oppositiones falso nomi- natae scientiae, quam nonnulli pro- fitentes circa fidem aberrauerunt. Gratia sit tecum. Amen. }}\pstart Breuissimum subtexit Epilogum, quo paterna et fingulari cura omni episcopo sanam commendat do- ctrinam, facilem et inanem quorundam loquacita- tem damnat. Atqui hoc potissimum huius epistolae ca- put fuisse iam saepius monui. Caeterum doctrinam sa- nam eleganti tropo ταρακαταθήκιεν id est commis- sum siue depositum et quasi quis dicat thesaurum uo cauit. Nam qua fide uir bonus depositum seruat ea- dem et maiore episcopus uerbum ueritatis custodit et cauet ne qua fraude uel mutetur uel corrumpatur  \pend
\section*{AD TIM. CAP. VI. }
\marginpar{[ p.174 ]}\pstart ctel prorsus surripiatur. Porro ex hoc thesauro demen sum cibum ecclesiae ut fidelis et prudens dispensator diuidet in tempore suo. Sana uero doctrina corrum- pitur prophana uocum nouitate et praetextu scientiae quae speciem quidem sapientiae habet reuera autem se ductio et impostura est. Intellexit autem Paulus quae cunque ab hominibus aliena a uera religione inueniun- tur et speciem quidem habent religionis et pietatis sed reuera ueritati repugnant. Cuiusmodi hodie mul- ta inter Pontifitios uoluntantur et reuolutantur tan- taque necessitate urgentur ut nisi quis nouis istis formis utatur protinus habeatur pro haeretico. Verum quis nam sit istarum quaestionum et argutationum fru- ctus ex eo patet quod sequitur, Quam nonnulli pro- fitentes circa fidem aberrauerunt, id quod et I. cap. inculcauit. Proinde cum nullum sit doctrinae genus an tiquius purius simplitius et melius quam quod autho re Christo tradiderunt apostoli, relictis omnibus aliis soli huic adhaereamus. Ad finem epistolae propria ma- nu fidei notam adiecit, Gratia sit tecum.  \pend\pstart  \pend
\endnumbering
\end{pages}
\end{document}
        